\documentclass{beamer}
\usetheme{Madrid}
\usecolortheme{default}
\usepackage[utf8]{inputenc}
\usepackage{amsmath}
\usepackage{amsfonts}
\usepackage{amssymb}
\usepackage{tikz}

\title{CRISC Hard Questions from ISACA QAE Manual}
\subtitle{Expanded Bank --- Answers on the Next Page}
\author{Compiled from ISACA CRISC Review Questions, Answers \& Explanations Manual, 6th Edition}
\date{\today}

\begin{document}
	
	\begin{frame}
		\titlepage
	\end{frame}
	
	\begin{frame}{Outline}
		\tableofcontents
	\end{frame}
	
	% ============================================================
	% DOMAIN 1 --- GOVERNANCE (26%)
	% ============================================================
	\section{Domain 1: Governance}
	
	\begin{frame}{D1-Q1: Defining Risk Management Strategies}
		\begin{block}{Question}
			Which of the following is MOST important to determine when defining risk management strategies?
		\end{block}
		\begin{enumerate}[(A)]
			\item Risk assessment criteria
			\item IT architecture complexity
			\item Enterprise disaster recovery plan
			\item Business objectives and operations
		\end{enumerate}
	\end{frame}
	\begin{frame}{D1-Q1: Answer}
		\begin{block}{Correct Answer: D --- Business objectives and operations}
			The risk practitioner must analyze the enterprise's objectives and risk tolerance and define a risk management framework based on this analysis. Some enterprises accept known risk; others invest in mitigating controls.
		\end{block}
	\end{frame}
	
	\begin{frame}{D1-Q2: Risk Management Strategic Plan}
		\begin{block}{Question}
			Which of the following is the MOST important information to include in a risk management strategic plan?
		\end{block}
		\begin{enumerate}[(A)]
			\item Risk management staffing requirements
			\item Risk management mission statement
			\item Risk mitigation investment plans
			\item Current state and desired future state
		\end{enumerate}
	\end{frame}
	\begin{frame}{D1-Q2: Answer}
		\begin{block}{Correct Answer: D --- Current state and desired future state}
			It is most important to paint a vision for the future and draw a road map from the starting point, which requires that the current state and desired future state be fully understood.
		\end{block}
	\end{frame}
	
	\begin{frame}{D1-Q3: Business Unit Risk Roles}
		\begin{block}{Question}
			Which BEST describes the risk-related roles of an organizational business unit (BU)? The BU management team:
		\end{block}
		\begin{enumerate}[(A)]
			\item Owns the mitigation plan while board members identify and assess risk
			\item Owns the risk and is responsible for identifying, assessing, mitigating, and reporting
			\item Carries out risk responsibilities, but ultimate accountability belongs to the board
			\item Is ultimately accountable, and the board owns the risk
		\end{enumerate}
	\end{frame}
	\begin{frame}{D1-Q3: Answer}
		\begin{block}{Correct Answer: B}
			The BU is responsible for owning the risk and its resulting actions. Risk owners identify, measure, monitor, control, and report on risk to executive management as established by the corporate risk framework.
		\end{block}
	\end{frame}
	
	\begin{frame}{D1-Q4: Outsourced IT in Foreign Country}
		\begin{block}{Question}
			An enterprise has outsourced the majority of its IT department to a third party whose servers are in a foreign country. Which is the MOST critical security consideration?
		\end{block}
		\begin{enumerate}[(A)]
			\item Security breach notification may be delayed due to the time difference
			\item Additional network intrusion detection sensors are required
			\item The enterprise could be unable to monitor compliance with internal security and privacy guidelines
			\item Laws and regulations of the country of origin may not be enforceable in the foreign country
		\end{enumerate}
	\end{frame}
	\begin{frame}{D1-Q4: Answer}
		\begin{block}{Correct Answer: D}
			Laws and regulations of the country of origin may not be enforceable in the foreign country. Conversely, the laws of the foreign vendor may affect the enterprise.
		\end{block}
	\end{frame}
	
	\begin{frame}{D1-Q5: Improving Risk Management Process}
		\begin{block}{Question}
			Which of the following is MOST beneficial to improvement of an enterprise's risk management process?
		\end{block}
		\begin{enumerate}[(A)]
			\item Key risk indicators
			\item External benchmarking
			\item Latest risk assessment
			\item A maturity model
		\end{enumerate}
	\end{frame}
	\begin{frame}{D1-Q5: Answer}
		\begin{block}{Correct Answer: D --- A maturity model}
			A maturity model helps identify the status quo and the path to the desired state and is most helpful when an enterprise desires to improve a business process such as risk management.
		\end{block}
	\end{frame}
	
	\begin{frame}{D1-Q6: Independent Risk Review}
		\begin{block}{Question}
			Which is the PRIMARY reason for subjecting the risk management process to review by independent risk auditors and assessors?
		\end{block}
		\begin{enumerate}[(A)]
			\item To ensure that the risk results are consistent
			\item To ensure that the risk factors and risk profile are well-defined
			\item To correct any mistakes in risk assessment
			\item To validate the control weaknesses for management reporting
		\end{enumerate}
	\end{frame}
	\begin{frame}{D1-Q6: Answer}
		\begin{block}{Correct Answer: B}
			Risk profile and risk factors are defined during the risk assessment process; an independent review helps ensure that the underlying process is effective and helps identify areas for future improvement.
		\end{block}
	\end{frame}
	
	\begin{frame}{D1-Q7: Evaluating Legal/Regulatory Impact}
		\begin{block}{Question}
			What is the MOST effective method to evaluate the potential impact of legal, regulatory, and contractual requirements on business objectives?
		\end{block}
		\begin{enumerate}[(A)]
			\item A compliance-oriented gap analysis
			\item Interviews with business process stakeholders
			\item A mapping of compliance requirements to policies and procedures
			\item A compliance-oriented business impact analysis
		\end{enumerate}
	\end{frame}
	\begin{frame}{D1-Q7: Answer}
		\begin{block}{Correct Answer: D}
			A compliance-oriented BIA identifies compliance requirements to which the enterprise is subject and assesses their effect on business objectives and activities.
		\end{block}
	\end{frame}
	
	\begin{frame}{D1-Q8: Maturity Assessment RFP}
		\begin{block}{Question}
			An enterprise is hiring a consultant to determine the maturity level of the risk management program. The MOST important element of the RFP is the:
		\end{block}
		\begin{enumerate}[(A)]
			\item Sample deliverable
			\item Past experience of the engagement team
			\item Methodology used in the assessment
			\item References from other enterprises
		\end{enumerate}
	\end{frame}
	\begin{frame}{D1-Q8: Answer}
		\begin{block}{Correct Answer: C --- Methodology}
			Methodology illustrates the consultant's process and offers a basis to align expectations with execution of the assessment.
		\end{block}
	\end{frame}
	
	\begin{frame}{D1-Q9: New Law Affecting HR System}
		\begin{block}{Question}
			Shortly after the annual policy review, a risk practitioner becomes aware that a new law may affect security requirements for the HR system. The risk practitioner should:
		\end{block}
		\begin{enumerate}[(A)]
			\item Analyze what systems and technology-related processes may be impacted
			\item Ensure necessary adjustments are implemented during the next review cycle
			\item Initiate an ad hoc revision of the corporate policy
			\item Notify the system custodian to implement changes
		\end{enumerate}
	\end{frame}
	\begin{frame}{D1-Q9: Answer}
		\begin{block}{Correct Answer: A}
			Assessing what systems and technology-related processes may be impacted is the best course of action. The analysis must also determine whether existing controls already address the new requirements.
		\end{block}
	\end{frame}
	
	\begin{frame}{D1-Q10: Outsourced Data Processing Violation}
		\begin{block}{Question}
			An enterprise outsourced personnel data processing to a supplier, and a regulatory violation occurs during processing. Who will be held legally responsible?
		\end{block}
		\begin{enumerate}[(A)]
			\item The supplier, because it has operational responsibility
			\item The enterprise, because it owns the data
			\item The enterprise and the supplier
			\item The supplier, because it did not comply with the contract
		\end{enumerate}
	\end{frame}
	\begin{frame}{D1-Q10: Answer}
		\begin{block}{Correct Answer: B}
			The enterprise retains responsibility for management of and adherence to policies, procedures, and regulatory requirements. The regulatory authority will generally hold the enterprise responsible.
		\end{block}
	\end{frame}
	
	\begin{frame}{D1-Q11: Evaluating Internal Controls Effectiveness}
		\begin{block}{Question}
			Which role is responsible for evaluating the effectiveness of existing internal information security controls within an enterprise?
		\end{block}
		\begin{enumerate}[(A)]
			\item Data owner
			\item Senior management
			\item End users
			\item System auditor
		\end{enumerate}
	\end{frame}
	\begin{frame}{D1-Q11: Answer}
		\begin{block}{Correct Answer: D --- System auditor}
			The system auditor provides continuous feedback to senior management about the effectiveness of internal controls.
		\end{block}
	\end{frame}
	
	\begin{frame}{D1-Q12: Overall Effectiveness of Risk Program}
		\begin{block}{Question}
			Which of the following BEST ensures the overall effectiveness of a risk management program?
		\end{block}
		\begin{enumerate}[(A)]
			\item Obtaining feedback from all end users
			\item Assigning a dedicated risk manager to run the program
			\item Applying quantitative risk methodologies
			\item Gaining the participation of relevant stakeholders
		\end{enumerate}
	\end{frame}
	\begin{frame}{D1-Q12: Answer}
		\begin{block}{Correct Answer: D}
			Without participation from enterprise stakeholders and their active role in supervision and risk monitoring, the effectiveness of a risk management program cannot be guaranteed.
		\end{block}
	\end{frame}
	
	\begin{frame}{D1-Q13: Achieving Risk-Based Objectives}
		\begin{block}{Question}
			Which approach BEST helps an enterprise achieve risk-based organizational objectives?
		\end{block}
		\begin{enumerate}[(A)]
			\item Ensure that asset owners perform annual risk assessments
			\item Review and update the risk register regularly
			\item Assign a steering committee to the risk management process
			\item Embed risk management activities into business processes
		\end{enumerate}
	\end{frame}
	\begin{frame}{D1-Q13: Answer}
		\begin{block}{Correct Answer: D}
			The primary objective of embedding risk management activities into business processes is to achieve risk-based organizational objectives in the most effective manner possible.
		\end{block}
	\end{frame}
	
	\begin{frame}{D1-Q14: Regulatory Reliance on Reviews}
		\begin{block}{Question}
			When assessing the capability of the risk management process, a regulatory body would place the GREATEST reliance on:
		\end{block}
		\begin{enumerate}[(A)]
			\item A peer review
			\item An internal review
			\item An external review
			\item A process capability review
		\end{enumerate}
	\end{frame}
	\begin{frame}{D1-Q14: Answer}
		\begin{block}{Correct Answer: C --- External review}
			Regulatory entities generally use assessments performed by an objective and independent third party.
		\end{block}
	\end{frame}
	
	\begin{frame}{D1-Q15: Evaluating Current Risk Profile}
		\begin{block}{Question}
			A company has been improving its security and compliance program since the last review one year ago. What should it do to evaluate its current risk profile?
		\end{block}
		\begin{enumerate}[(A)]
			\item Review previous findings and ensure all issues are resolved
			\item Conduct follow-up audits in previously deficient areas
			\item Monitor KRIs and use results to develop targeted assessments
			\item Perform a new enterprise risk assessment using an independent expert
		\end{enumerate}
	\end{frame}
	\begin{frame}{D1-Q15: Answer}
		\begin{block}{Correct Answer: D}
			The best way to ensure the security posture is still compliant is to conduct another risk assessment. A lot can change in a year, and an independent expert provides more objective results.
		\end{block}
	\end{frame}
	
	\begin{frame}{D1-Q16: Inadequate Data/System Ownership Policy}
		\begin{block}{Question}
			Which is the GREATEST risk of a policy that defines data and system ownership inadequately?
		\end{block}
		\begin{enumerate}[(A)]
			\item Audit recommendations may not be implemented
			\item Users may have unauthorized access to create, modify, or delete data
			\item User management coordination does not exist
			\item Specific user accountability cannot be established
		\end{enumerate}
	\end{frame}
	\begin{frame}{D1-Q16: Answer}
		\begin{block}{Correct Answer: B}
			Without a policy defining who grants access to specific data or systems, the risk that employees receive access without justified business purpose increases.
		\end{block}
	\end{frame}
	
	\begin{frame}{D1-Q17: Management's Role in Risk Strategy}
		\begin{block}{Question}
			Which BEST describes the role of management in implementing a risk management strategy?
		\end{block}
		\begin{enumerate}[(A)]
			\item Ensure that planning, budgeting, and performance of security components are appropriate
			\item Assess and incorporate the results of risk management into decision-making
			\item Identify, evaluate, and minimize risk to IT systems supporting the mission
			\item Understand the risk management process to develop training materials
		\end{enumerate}
	\end{frame}
	\begin{frame}{D1-Q17: Answer}
		\begin{block}{Correct Answer: B}
			Assessing and incorporating the results of risk management into the decision-making process best describes the role of senior management.
		\end{block}
	\end{frame}
	
	\begin{frame}{D1-Q18: Primary Focus of IT Business Risk}
		\begin{block}{Question}
			The PRIMARY focus of managing IT-related business risk is to protect:
		\end{block}
		\begin{enumerate}[(A)]
			\item Information
			\item Hardware
			\item Applications
			\item Databases
		\end{enumerate}
	\end{frame}
	\begin{frame}{D1-Q18: Answer}
		\begin{block}{Correct Answer: A --- Information}
			The primary objective for any enterprise is to protect mission-critical information based on a risk assessment.
		\end{block}
	\end{frame}
	
	\begin{frame}{D1-Q19: Best Risk Management Perspective}
		\begin{block}{Question}
			Which can provide the BEST perspective of risk management to an enterprise's employees and stockholders?
		\end{block}
		\begin{enumerate}[(A)]
			\item An interdisciplinary team within the enterprise
			\item A third-party risk assessment service provider
			\item The enterprise's IT department
			\item The enterprise's internal compliance department
		\end{enumerate}
	\end{frame}
	\begin{frame}{D1-Q19: Answer}
		\begin{block}{Correct Answer: A}
			An interdisciplinary team ensures all areas are adequately considered in risk assessment and provides an enterprise-wide perspective.
		\end{block}
	\end{frame}
	
	\begin{frame}{D1-Q20: Global Policy for Multiple Regions}
		\begin{block}{Question}
			Which approach to corporate policy BEST supports an enterprise's expansion to other regions where different local laws apply?
		\end{block}
		\begin{enumerate}[(A)]
			\item A global policy without provisions that might be disputed locally
			\item A global policy amended to comply with local laws
			\item A global policy that complies with laws at enterprise headquarters
			\item Local policies to accommodate laws within each region
		\end{enumerate}
	\end{frame}
	\begin{frame}{D1-Q20: Answer}
		\begin{block}{Correct Answer: B}
			A global policy including local amendments ensures alignment with local laws and regulations.
		\end{block}
	\end{frame}
	
	\begin{frame}{D1-Q21: Data Classification Responsibility}
		\begin{block}{Question}
			Who is MOST likely responsible for data classification?
		\end{block}
		\begin{enumerate}[(A)]
			\item Data user
			\item Data owner
			\item Data custodian
			\item System administrator
		\end{enumerate}
	\end{frame}
	\begin{frame}{D1-Q21: Answer}
		\begin{block}{Correct Answer: B --- Data owner}
			The data owner classifies data according to the enterprise's data classification scheme.
		\end{block}
	\end{frame}
	
	\begin{frame}{D1-Q22: Security Governance Model Impact}
		\begin{block}{Question}
			Which factor will have the GREATEST impact on the type of information security governance model an enterprise adopts?
		\end{block}
		\begin{enumerate}[(A)]
			\item The number of employees
			\item The enterprise's budget
			\item The organizational structure
			\item The type of technology used
		\end{enumerate}
	\end{frame}
	\begin{frame}{D1-Q22: Answer}
		\begin{block}{Correct Answer: C --- Organizational structure}
			Information security governance models depend significantly on the overall organizational structure.
		\end{block}
	\end{frame}
	
	\begin{frame}{D1-Q23: Information Security Procedures}
		\begin{block}{Question}
			Information security procedures should:
		\end{block}
		\begin{enumerate}[(A)]
			\item Be updated frequently as new software is released
			\item Underline the importance of security governance
			\item Define the allowable limits of behavior
			\item Describe security baselines for each platform
		\end{enumerate}
	\end{frame}
	\begin{frame}{D1-Q23: Answer}
		\begin{block}{Correct Answer: A}
			Security procedures must change frequently to keep up with software changes. Because a procedure is a how-to document, it must be kept current.
		\end{block}
	\end{frame}
	
	\begin{frame}{D1-Q24: Greatest Benefit of Risk-Aware Culture}
		\begin{block}{Question}
			Which is the GREATEST benefit of a risk-aware culture?
		\end{block}
		\begin{enumerate}[(A)]
			\item Issues are escalated when suspicious activity is noticed
			\item Controls are double-checked to anticipate any issues
			\item Individuals communicate with peers for knowledge sharing
			\item Employees are self-motivated to learn about costs and benefits
		\end{enumerate}
	\end{frame}
	\begin{frame}{D1-Q24: Answer}
		\begin{block}{Correct Answer: A}
			Management benefits most from an escalation process because risk and incidents are reported in a timely manner.
		\end{block}
	\end{frame}
	
	\begin{frame}{D1-Q25: Main Objective of IT Risk Management}
		\begin{block}{Question}
			The MAIN objective of IT risk management is to:
		\end{block}
		\begin{enumerate}[(A)]
			\item Prevent loss of IT assets
			\item Provide timely management reports
			\item Ensure regulatory compliance
			\item Enable risk-aware business decisions
		\end{enumerate}
	\end{frame}
	\begin{frame}{D1-Q25: Answer}
		\begin{block}{Correct Answer: D}
			IT risk management should be conducted as part of enterprise-wide risk management, whose ultimate objective is to support risk-aware business decisions.
		\end{block}
	\end{frame}
	
	\begin{frame}{D1-Q26: Multiple Governmental Jurisdictions}
		\begin{block}{Question}
			Which is the BEST approach for a global enterprise subject to regulation by multiple governmental jurisdictions with differing requirements?
		\end{block}
		\begin{enumerate}[(A)]
			\item Bring all locations into conformity with the aggregate requirements of all jurisdictions
			\item Bring all locations into conformity with generally accepted industry good practices
			\item Establish a baseline standard incorporating the requirements all jurisdictions have in common
			\item Establish baseline standards for all locations and add supplemental standards as required
		\end{enumerate}
	\end{frame}
	\begin{frame}{D1-Q26: Answer}
		\begin{block}{Correct Answer: D}
			It is more efficient to establish a baseline standard and then develop additional standards for locations that must meet specific requirements.
		\end{block}
	\end{frame}
	
	\begin{frame}{D1-Q27: Final Sign-off on Risk Plan}
		\begin{block}{Question}
			Who MUST give final sign-off on the IT risk management plan?
		\end{block}
		\begin{enumerate}[(A)]
			\item IT auditors performing the risk assessment
			\item Business process owners
			\item Senior managers
			\item IT security administrators
		\end{enumerate}
	\end{frame}
	\begin{frame}{D1-Q27: Answer}
		\begin{block}{Correct Answer: C --- Senior managers}
			Senior managers understand performance metrics, approved the policies and standards that govern the enterprise, and have final responsibility for risk.
		\end{block}
	\end{frame}
	
	\begin{frame}{D1-Q28: Security Boundary Before Assessment}
		\begin{block}{Question}
			Which is the PRIMARY reason a risk practitioner determines the security boundary prior to conducting a risk assessment?
		\end{block}
		\begin{enumerate}[(A)]
			\item To decide which laws and regulations apply
			\item To identify the scope of the risk assessment
			\item To identify the business owners of the system
			\item To decide whether quantitative or qualitative analysis is appropriate
		\end{enumerate}
	\end{frame}
	\begin{frame}{D1-Q28: Answer}
		\begin{block}{Correct Answer: B --- Identify the scope}
			Identifying the security boundary establishes the fundamental scope of inquiry, including which systems and components are subject to assessment.
		\end{block}
	\end{frame}
	
	\begin{frame}{D1-Q29: Who Manages/Executes Risk Program}
		\begin{block}{Question}
			Which group would be MOST effective in managing and executing an enterprise's risk program?
		\end{block}
		\begin{enumerate}[(A)]
			\item Mid-level managers
			\item Senior managers
			\item Frontline employees
			\item The incident response team
		\end{enumerate}
	\end{frame}
	\begin{frame}{D1-Q29: Answer}
		\begin{block}{Correct Answer: A --- Mid-level managers}
			Mid-level managers are most centrally located within the organizational hierarchy, combining sufficient breadth of influence with proximity to day-to-day operations.
		\end{block}
	\end{frame}
	
	\begin{frame}{D1-Q30: IT Steering Committee Representation}
		\begin{block}{Question}
			The IT steering committee of a one-year-old start-up makes all IT decisions, including technology budget. It will be BEST represented by:
		\end{block}
		\begin{enumerate}[(A)]
			\item Members of the executive board
			\item High-level members of the IT department
			\item IT experts from outside the enterprise
			\item Key members from each department
		\end{enumerate}
	\end{frame}
	\begin{frame}{D1-Q30: Answer}
		\begin{block}{Correct Answer: D}
			The IT steering committee should be comprised of individuals from each department to ensure the entire enterprise is represented.
		\end{block}
	\end{frame}
	
	\begin{frame}{D1-Q31: First Line of Defense Purpose}
		\begin{block}{Question}
			Which is one of the MAIN purposes of the first line of defense in the three lines of defense model?
		\end{block}
		\begin{enumerate}[(A)]
			\item Ensure that financial controls are in place
			\item Ensure control deficiencies are addressed
			\item Ensure risk management practices are effective
			\item Ensure compliance with rules and regulations
		\end{enumerate}
	\end{frame}
	\begin{frame}{D1-Q31: Answer}
		\begin{block}{Correct Answer: B}
			In the first line, operational managers own the risk and address it by mitigating control deficiencies.
		\end{block}
	\end{frame}
	
	\begin{frame}{D1-Q32: Selecting Risk Response Technique}
		\begin{block}{Question}
			The PRIMARY consideration when selecting a risk response technique is:
		\end{block}
		\begin{enumerate}[(A)]
			\item Coverage of all identified risk
			\item Availability of resources
			\item Enterprise goals and objectives
			\item Standards and industry good practices
		\end{enumerate}
	\end{frame}
	\begin{frame}{D1-Q32: Answer}
		\begin{block}{Correct Answer: C}
			The risk response will be based primarily on the goals and objectives of the enterprise.
		\end{block}
	\end{frame}
	
	\begin{frame}{D1-Q33: First Policy for Start-up}
		\begin{block}{Question}
			The board of a one-year-old start-up asked the CIO to create all IT policies and procedures. Which should the CIO create FIRST?
		\end{block}
		\begin{enumerate}[(A)]
			\item The strategic IT plan
			\item The data classification scheme
			\item The information architecture document
			\item The technology infrastructure plan
		\end{enumerate}
	\end{frame}
	\begin{frame}{D1-Q33: Answer}
		\begin{block}{Correct Answer: A --- Strategic IT plan}
			The strategic IT plan is the first policy to create when developing an enterprise's governance model.
		\end{block}
	\end{frame}
	
	\begin{frame}{D1-Q34: Most Concerning to Risk Practitioner}
		\begin{block}{Question}
			Which should be of MOST concern to a risk practitioner?
		\end{block}
		\begin{enumerate}[(A)]
			\item Failure to notify the public of an intrusion
			\item Failure to notify the police of an attempted intrusion
			\item Failure to internally report a successful attack
			\item Failure to examine access rights periodically
		\end{enumerate}
	\end{frame}
	\begin{frame}{D1-Q34: Answer}
		\begin{block}{Correct Answer: C}
			Failure to report a successful intrusion is a serious concern and could in some instances be interpreted as abetting.
		\end{block}
	\end{frame}
	
	\begin{frame}{D1-Q35: Insight into Risk Capabilities}
		\begin{block}{Question}
			Which review will provide the MOST insight into an enterprise's risk management capabilities?
		\end{block}
		\begin{enumerate}[(A)]
			\item A capability maturity model review
			\item A capability comparison with industry standards or regulations
			\item A self-assessment of capabilities
			\item An internal audit review of capabilities
		\end{enumerate}
	\end{frame}
	\begin{frame}{D1-Q35: Answer}
		\begin{block}{Correct Answer: A}
			Capability maturity modeling allows an enterprise to understand its level of maturity in its risk capabilities.
		\end{block}
	\end{frame}
	
	\begin{frame}{D1-Q36: Leveraging Internal Audit in ERM}
		\begin{block}{Question}
			As part of an ERM program, a risk practitioner BEST leverages the work performed by internal audit by having it:
		\end{block}
		\begin{enumerate}[(A)]
			\item Design, implement, and maintain the ERM process
			\item Manage and assess the overall risk awareness
			\item Evaluate ongoing changes to enterprise risk factors
			\item Assist in monitoring, evaluating, examining, and reporting on controls
		\end{enumerate}
	\end{frame}
	\begin{frame}{D1-Q36: Answer}
		\begin{block}{Correct Answer: D}
			The internal audit function is responsible for assisting management and the board in monitoring, evaluating, examining, and reporting on internal controls.
		\end{block}
	\end{frame}
	
	\begin{frame}{D1-Q37: Interconnections with External Systems}
		\begin{block}{Question}
			Which is the BIGGEST concern for a CISO regarding interconnections with systems outside the enterprise?
		\end{block}
		\begin{enumerate}[(A)]
			\item Requirements to comply with each other's contractual security obligations
			\item Uncertainty that the other system will be available as needed
			\item The ability to perform risk assessments on the other system
			\item Ensuring communications are encrypted through a VPN tunnel
		\end{enumerate}
	\end{frame}
	\begin{frame}{D1-Q37: Answer}
		\begin{block}{Correct Answer: A}
			Ensuring both systems comply with mutual contractual security obligations should be the primary concern.
		\end{block}
	\end{frame}
	
	\begin{frame}{D1-Q38: IT Organizational Structure}
		\begin{block}{Question}
			The CIO creates all IT policies and procedures, managed and approved by the IT steering committee, which makes all IT decisions including the technology budget. Which type of IT organizational structure is this?
		\end{block}
		\begin{enumerate}[(A)]
			\item Project-based
			\item Centralized
			\item Decentralized
			\item Divisional
		\end{enumerate}
	\end{frame}
	\begin{frame}{D1-Q38: Answer}
		\begin{block}{Correct Answer: B --- Centralized}
			Within a centralized IT organizational structure, one group makes all decisions for the entire enterprise.
		\end{block}
	\end{frame}
	
	\begin{frame}{D1-Q39: Communicating Continuous Monitoring}
		\begin{block}{Question}
			The aggregated results of continuous monitoring activities are BEST communicated to:
		\end{block}
		\begin{enumerate}[(A)]
			\item The risk owner
			\item Technical staff
			\item The audit department
			\item The information security manager
		\end{enumerate}
	\end{frame}
	\begin{frame}{D1-Q39: Answer}
		\begin{block}{Correct Answer: A --- The risk owner}
			The risk owner is accountable for the fact that appropriate risk responses are executed in alignment with the enterprise's risk appetite.
		\end{block}
	\end{frame}
	
	\begin{frame}{D1-Q40: IT Risk Awareness Program}
		\begin{block}{Question}
			Which is a PRIMARY consideration when developing an IT risk awareness program?
		\end{block}
		\begin{enumerate}[(A)]
			\item Why technology risk is owned by IT
			\item How technology risk can affect each attendee's area of business
			\item How business process owners can transfer technology risk
			\item Why technology risk is more difficult to manage than other risk
		\end{enumerate}
	\end{frame}
	\begin{frame}{D1-Q40: Answer}
		\begin{block}{Correct Answer: B}
			Stakeholders must understand how IT-related risk affects the overall business.
		\end{block}
	\end{frame}
	
	\begin{frame}{D1-Q41: IT Delivering Against Requirements}
		\begin{block}{Question}
			Management wants to ensure IT is successful in delivering against business requirements. Which BEST supports that effort?
		\end{block}
		\begin{enumerate}[(A)]
			\item An internal control system or framework
			\item A cost-benefit analysis
			\item A return on investment analysis
			\item A benchmark process
		\end{enumerate}
	\end{frame}
	\begin{frame}{D1-Q41: Answer}
		\begin{block}{Correct Answer: A}
			For IT to be successful in delivering against business requirements, management should develop an internal control system that supports its business requirements.
		\end{block}
	\end{frame}
	
	\begin{frame}{D1-Q42: Aligning Risk Appetite}
		\begin{block}{Question}
			It is MOST important that risk appetite be aligned with business objectives to ensure that:
		\end{block}
		\begin{enumerate}[(A)]
			\item Resources are directed toward areas of low risk tolerance
			\item Major risk is identified and eliminated
			\item IT and business goals are aligned
			\item The risk strategy is adequately communicated
		\end{enumerate}
	\end{frame}
	\begin{frame}{D1-Q42: Answer}
		\begin{block}{Correct Answer: A}
			Aligning risk appetite with business objectives allows an enterprise to evaluate and deploy valuable resources toward those objectives with low risk tolerance for loss.
		\end{block}
	\end{frame}
	
	\begin{frame}{D1-Q43: Out-of-Date Security Policy}
		\begin{block}{Question}
			Which causes the GREATEST concern to a risk practitioner reviewing a corporate information security policy that is out of date? The policy:
		\end{block}
		\begin{enumerate}[(A)]
			\item Was not reviewed within the last three years
			\item Is missing newer technologies/platforms
			\item Was not updated to account for new locations
			\item Does not enforce control monitoring
		\end{enumerate}
	\end{frame}
	\begin{frame}{D1-Q43: Answer}
		\begin{block}{Correct Answer: A}
			Not reviewing the policy for three years and updating it as necessary does not follow good practices and is the greatest concern.
		\end{block}
	\end{frame}
	
	\begin{frame}{D1-Q44: Accountable for IT Business Risk}
		\begin{block}{Question}
			Who is accountable for business risk related to IT?
		\end{block}
		\begin{enumerate}[(A)]
			\item Chief information officer
			\item Chief financial officer
			\item Users of IT services
			\item Chief architect
		\end{enumerate}
	\end{frame}
	\begin{frame}{D1-Q44: Answer}
		\begin{block}{Correct Answer: C --- Users of IT services}
			The enterprise (i.e., the users of IT services) ultimately owns business-related risk, including the risk related to the use of IT.
		\end{block}
	\end{frame}
	
	\begin{frame}{D1-Q45: Expanding to Nearby Domestic Locations}
		\begin{block}{Question}
			An enterprise is expanding into new nearby domestic locations. Which is MOST important for a risk practitioner to report on?
		\end{block}
		\begin{enumerate}[(A)]
			\item Competitor analysis
			\item Legal and regulatory requirements
			\item Political issues
			\item Potential of natural disasters
		\end{enumerate}
	\end{frame}
	\begin{frame}{D1-Q45: Answer}
		\begin{block}{Correct Answer: B}
			Legal and regulatory requirements are most likely to change when moving to a nearby location because each municipality may enforce significantly different regulations.
		\end{block}
	\end{frame}
	
	\begin{frame}{D1-Q46: True Statement About IT Risk}
		\begin{block}{Question}
			Which of the following is true about IT risk?
		\end{block}
		\begin{enumerate}[(A)]
			\item IT risk cannot be assessed and measured quantitatively
			\item IT risk should be calculated separately from business risk
			\item IT risk management is the responsibility of the IT department
			\item IT risk exists regardless of whether it is detected or recognized
		\end{enumerate}
	\end{frame}
	\begin{frame}{D1-Q46: Answer}
		\begin{block}{Correct Answer: D}
			The enterprise must identify, acknowledge, and respond to risk; ignorance of risk is not acceptable.
		\end{block}
	\end{frame}
	
	\begin{frame}{D1-Q47: Selecting Risk Management Methodology}
		\begin{block}{Question}
			Which is MOST important when selecting an appropriate risk management methodology?
		\end{block}
		\begin{enumerate}[(A)]
			\item Risk culture
			\item Countermeasure analysis
			\item Cost-benefit analysis
			\item Risk transfer strategy
		\end{enumerate}
	\end{frame}
	\begin{frame}{D1-Q47: Answer}
		\begin{block}{Correct Answer: A --- Risk culture}
			Without understanding risk culture---how and why an enterprise makes decisions regarding risk---one cannot select a risk management methodology.
		\end{block}
	\end{frame}
	
	\begin{frame}{D1-Q48: Compliance with Risk Appetite}
		\begin{block}{Question}
			Which BEST determines compliance with the risk appetite of an enterprise?
		\end{block}
		\begin{enumerate}[(A)]
			\item Balance between preventive and detective controls
			\item Inherent risk and acceptable risk level
			\item Residual risk and acceptable risk level
			\item Balance between countermeasures to threats and preventive controls
		\end{enumerate}
	\end{frame}
	\begin{frame}{D1-Q48: Answer}
		\begin{block}{Correct Answer: C}
			Considering residual risk in terms of acceptable risk yields risk that is appropriately balanced after the application of controls.
		\end{block}
	\end{frame}
	
	\begin{frame}{D1-Q49: Monitoring Control Effectiveness vs External Requirements}
		\begin{block}{Question}
			A key objective when monitoring information systems control effectiveness against the enterprise's external requirements is to:
		\end{block}
		\begin{enumerate}[(A)]
			\item Design the applicable information security controls for external audits
			\item Create the enterprise's information security policy provisions for third parties
			\item Ensure that the enterprise's legal obligations have been satisfied
			\item Identify those legal obligations that apply to the enterprise's security practices
		\end{enumerate}
	\end{frame}
	\begin{frame}{D1-Q49: Answer}
		\begin{block}{Correct Answer: C}
			Legal obligations are one of the principal external requirements that necessitate compliance monitoring.
		\end{block}
	\end{frame}
	
	\begin{frame}{D1-Q50: Overall Risk Status}
		\begin{block}{Question}
			Risk monitoring provides timely information on the actual status of the enterprise with regard to risk. Which choice provides an overall risk status of the enterprise?
		\end{block}
		\begin{enumerate}[(A)]
			\item Risk management
			\item Risk analysis
			\item Risk appetite
			\item Risk profile
		\end{enumerate}
	\end{frame}
	\begin{frame}{D1-Q50: Answer}
		\begin{block}{Correct Answer: D --- Risk profile}
			The risk profile provides the current overall portfolio of identified risk to which the enterprise is exposed.
		\end{block}
	\end{frame}
	
	\begin{frame}{D1-Q51: Firewall Deviation from Security Policy}
		\begin{block}{Question}
			Which activity should a risk professional perform to determine whether firewall deployments are deviating from the enterprise's information security policy?
		\end{block}
		\begin{enumerate}[(A)]
			\item Review the firewall parameter settings
			\item Review the firewall intrusion prevention system logs
			\item Review the firewall hardening procedures
			\item Analyze the firewall log file for recent attacks
		\end{enumerate}
	\end{frame}
	\begin{frame}{D1-Q51: Answer}
		\begin{block}{Correct Answer: A}
			Firewall parameter settings tie into configurations linked to the governing security policy.
		\end{block}
	\end{frame}
	
	\begin{frame}{D1-Q52: Continuous Risk Management Benefit}
		\begin{block}{Question}
			Which is the MOST significant benefit of implementing a continuous risk management process in an enterprise?
		\end{block}
		\begin{enumerate}[(A)]
			\item It complies with external regulators' requirements
			\item It captures changes to the enterprise's risk profile
			\item It ensures stakeholders are linked to the identified risk
			\item It ensures timely action is taken to mitigate risk
		\end{enumerate}
	\end{frame}
	\begin{frame}{D1-Q52: Answer}
		\begin{block}{Correct Answer: B}
			A continuous risk management process ensures changes to the enterprise's risk profile are captured.
		\end{block}
	\end{frame}
	
	\begin{frame}{D1-Q53: High Risk vs Moderate Appetite}
		\begin{block}{Question}
			Senior management defined the enterprise risk appetite as moderate. A business-critical application poses a high risk. What is the NEXT action?
		\end{block}
		\begin{enumerate}[(A)]
			\item Remove the high-risk application and replace it with another system
			\item Recommend that management increase its acceptable risk level
			\item Assess the impact of planned controls to the application
			\item Restrict access to the application only to trusted users
		\end{enumerate}
	\end{frame}
	\begin{frame}{D1-Q53: Answer}
		\begin{block}{Correct Answer: C}
			The risk practitioner should determine whether planned implementation of new controls may lower the risk from high to moderate or low before taking further action.
		\end{block}
	\end{frame}
	
	\begin{frame}{D1-Q54: Developing the Risk Profile}
		\begin{block}{Question}
			Which BEST assists in the development of the risk profile?
		\end{block}
		\begin{enumerate}[(A)]
			\item The presence of preventive and detective controls
			\item Inherent risk and detection risk
			\item Cost-benefit analysis of controls
			\item Likelihood and impact of risk
		\end{enumerate}
	\end{frame}
	\begin{frame}{D1-Q54: Answer}
		\begin{block}{Correct Answer: D}
			Likelihood and impact of risk help in the development of the risk profile.
		\end{block}
	\end{frame}
	
	\begin{frame}{D1-Q55: Employee Handbook in Multiple Regions}
		\begin{block}{Question}
			An enterprise expanded into Europe, Asia, and Latin America and has a single-version, multiple-language handbook last updated three years ago. Which is of MOST concern?
		\end{block}
		\begin{enumerate}[(A)]
			\item The handbook may not have been correctly translated
			\item Newer policies may not be included
			\item Expired policies may be included
			\item The handbook may violate local laws and regulations
		\end{enumerate}
	\end{frame}
	\begin{frame}{D1-Q55: Answer}
		\begin{block}{Correct Answer: D}
			Because customs and laws affect an enterprise's ability to operate in a location and vary by state and country, it is critical for the employee handbook to acknowledge and account for regional and national differences.
		\end{block}
	\end{frame}
	
	\begin{frame}{D1-Q56: Risk Owner Purpose}
		\begin{block}{Question}
			The primary reason for having a risk owner is:
		\end{block}
		\begin{enumerate}[(A)]
			\item Leverage their resources
			\item Ensure accountability
			\item Provide flexibility
			\item Seek assistance
		\end{enumerate}
	\end{frame}
	\begin{frame}{D1-Q56: Answer}
		\begin{block}{Correct Answer: B --- Ensure accountability}
			The risk owner drives accountability in risk remediation.
		\end{block}
	\end{frame}
	
	\begin{frame}{D1-Q57: Risk Utilities}
		\begin{block}{Question}
			Which BEST describes the utility of a risk?
		\end{block}
		\begin{enumerate}[(A)]
			\item The finance incentive behind the risk
			\item The potential opportunity of the risk
			\item The mechanics of how a risk works
			\item The usefulness of the risk to individuals or groups
		\end{enumerate}
	\end{frame}
	\begin{frame}{D1-Q57: Answer}
		\begin{block}{Correct Answer: D}
			Risk utility describes the usefulness of a particular risk to an individual or group.
		\end{block}
	\end{frame}
	
	\begin{frame}{D1-Q58: Scalability of Monitoring Tools}
		\begin{block}{Question}
			Which monitoring tool aspect ensures ability to keep up with enterprise growth?
		\end{block}
		\begin{enumerate}[(A)]
			\item Scalability
			\item Customizability
			\item Sustainability
			\item Impact on performance
		\end{enumerate}
	\end{frame}
	\begin{frame}{D1-Q58: Answer}
		\begin{block}{Correct Answer: A}
			Scalability ensures monitoring tools can meet anticipated growth in process, complexity, or transaction volumes.
		\end{block}
	\end{frame}
	
	\begin{frame}{D1-Q59: Moderate Risk Definition}
		\begin{block}{Question}
			A risk that is a noticeable failure threatening the success of certain goals exists at which level?
		\end{block}
		\begin{enumerate}[(A)]
			\item Moderate risk
			\item High risk
			\item Extremely high risk
			\item Low risk
		\end{enumerate}
	\end{frame}
	\begin{frame}{D1-Q59: Answer}
		\begin{block}{Correct Answer: A}
			Moderate risks are noticeable failures threatening the success of certain goals.
		\end{block}
	\end{frame}
	
	\begin{frame}{D1-Q60: Grouping Risks by Common Causes}
		\begin{block}{Question}
			What is the primary advantage to group risks by common causes during qualitative risk analysis?
		\end{block}
		\begin{enumerate}[(A)]
			\item Realize areas most laden with risks
			\item Assist in developing effective risk responses
			\item Save time by collecting related resources
			\item Create risk categories unique to each project
		\end{enumerate}
	\end{frame}
	\begin{frame}{D1-Q60: Answer}
		\begin{block}{Correct Answer: B}
			By grouping risks by categories, the project team can develop effective risk responses. Related risks often have common causal factors.
		\end{block}
	\end{frame}
	
	\begin{frame}{D1-Q61: Risk Communication Definition}
		\begin{block}{Question}
			"Process of exchanging information and views about risks among stakeholders" describes which process?
		\end{block}
		\begin{enumerate}[(A)]
			\item Risk governance
			\item Risk identification
			\item Risk response planning
			\item Risk communication
		\end{enumerate}
	\end{frame}
	\begin{frame}{D1-Q61: Answer}
		\begin{block}{Correct Answer: D}
			Risk communication is exactly defined as exchanging information and views about risks among stakeholders.
		\end{block}
	\end{frame}
	
	\begin{frame}{D1-Q62: Risk Register Output of Identify Risks}
		\begin{block}{Question}
			Which is a key output of the Identify Risks process?
		\end{block}
		\begin{enumerate}[(A)]
			\item Risk Register
			\item Risk Management Plan
			\item Risk Breakdown Structure
			\item Risk Categories
		\end{enumerate}
	\end{frame}
	\begin{frame}{D1-Q62: Answer}
		\begin{block}{Correct Answer: A}
			The primary output from Identify Risks is the initial entries into the risk register.
		\end{block}
	\end{frame}
	
	\begin{frame}{D1-Q63: Scenario Actors}
		\begin{block}{Question}
			Which risk scenario component has potential to generate internal or external threat?
		\end{block}
		\begin{enumerate}[(A)]
			\item Timing dimension
			\item Events
			\item Assets
			\item Actors
		\end{enumerate}
	\end{frame}
	\begin{frame}{D1-Q63: Answer}
		\begin{block}{Correct Answer: D}
			Actors are those components that have the potential to generate the threat (internal or external, human or non-human).
		\end{block}
	\end{frame}
	
	\begin{frame}{D1-Q64: Before Relying on Controls}
		\begin{block}{Question}
			What should you do before relying on any control?
		\end{block}
		\begin{enumerate}[(A)]
			\item Review performance data
			\item Discover risk exposure
			\item Conduct pilot testing
			\item Articulate risk
		\end{enumerate}
	\end{frame}
	\begin{frame}{D1-Q64: Answer}
		\begin{block}{Correct Answer: A and C}
			Pilot testing and reviewing performance data to verify operation against design are done before relying on control.
		\end{block}
	\end{frame}
	
	\begin{frame}{D1-Q65: Maturity Level 1}
		\begin{block}{Question}
			Which statement is NOT true for risk management capability maturity level 1?
		\end{block}
		\begin{enumerate}[(A)]
			\item Risk is viewed as technical issue
			\item Decisions involving risk lack credible information
			\item Risk appetite applied only during episodic assessments
			\item Risk management skills exist on ad hoc basis
		\end{enumerate}
	\end{frame}
	\begin{frame}{D1-Q65: Answer}
		\begin{block}{Correct Answer: B}
			Decisions lacking credible information is a characteristic of Level 0, not Level 1.
		\end{block}
	\end{frame}
	
	\begin{frame}{D1-Q66: Most Important Stakeholder for Risk Review}
		\begin{block}{Question}
			Which stakeholder is MOST important for reviewing risk response options?
		\end{block}
		\begin{enumerate}[(A)]
			\item Information security managers
			\item Internal auditors
			\item Incident response team members
			\item Business managers
		\end{enumerate}
	\end{frame}
	\begin{frame}{D1-Q66: Answer}
		\begin{block}{Correct Answer: D}
			Business managers are accountable for managing associated risk and determine actions based on information provided by others.
		\end{block}
	\end{frame}
	
	\begin{frame}{D1-Q67: Fault Tree Analysis Purpose}
		\begin{block}{Question}
			Which technique provides systematic description of unwanted occurrences in a system?
		\end{block}
		\begin{enumerate}[(A)]
			\item Sensitivity analysis
			\item Scenario analysis
			\item Fault tree analysis
			\item Cause and effect analysis
		\end{enumerate}
	\end{frame}
	\begin{frame}{D1-Q67: Answer}
		\begin{block}{Correct Answer: C}
			Fault Tree Analysis (FTA) provides systematic description of possible occurrences that can result in undesirable outcomes.
		\end{block}
	\end{frame}
	
	\begin{frame}{D1-Q68: Risk Treatment Definition}
		\begin{block}{Question}
			What is the process for selecting and implementing measures to impact risk called?
		\end{block}
		\begin{enumerate}[(A)]
			\item Risk Treatment
			\item Control
			\item Risk Assessment
			\item Risk Management
		\end{enumerate}
	\end{frame}
	\begin{frame}{D1-Q68: Answer}
		\begin{block}{Correct Answer: A}
			Risk treatment is the process for selecting and implementing measures for impacting risk in the environment.
		\end{block}
	\end{frame}
	
	\begin{frame}{D1-Q69: SOX Section 302}
		\begin{block}{Question}
			Which SOX section specifies "Periodic financial reports must be certified by CEO and CFO"?
		\end{block}
		\begin{enumerate}[(A)]
			\item Section 302
			\item Section 404
			\item Section 203
			\item Section 409
		\end{enumerate}
	\end{frame}
	\begin{frame}{D1-Q69: Answer}
		\begin{block}{Correct Answer: A --- Section 302}
			Section 302 requires corporate responsibility for financial reports to be certified by CEO, CFO, or designated representative.
		\end{block}
	\end{frame}
	
	\begin{frame}{D1-Q70: Control Assessment Primary Need}
		\begin{block}{Question}
			What is the PRIMARY need for effectively assessing controls?
		\end{block}
		\begin{enumerate}[(A)]
			\item Control's alignment with operating environment
			\item Control's design effectiveness
			\item Control's objective achievement
			\item Control's operating effectiveness
		\end{enumerate}
	\end{frame}
	\begin{frame}{D1-Q70: Answer}
		\begin{block}{Correct Answer: C}
			Controls can be effectively assessed only by determining how accurately the control objective is achieved within the operating environment.
		\end{block}
	\end{frame}
	
	\begin{frame}{D1-Q71: Fast Tracking Effect}
		\begin{block}{Question}
			When you fast track the project, what is likely to increase?
		\end{block}
		\begin{enumerate}[(A)]
			\item Human resource needs
			\item Quality control concerns
			\item Costs
			\item Risks
		\end{enumerate}
	\end{frame}
	\begin{frame}{D1-Q71: Answer}
		\begin{block}{Correct Answer: D --- Risks}
			Fast tracking allows phases to overlap and generally increases risks within the project.
		\end{block}
	\end{frame}
	
	\begin{frame}{D1-Q72: Risk Mitigation Action}
		\begin{block}{Question}
			David takes actions to ensure risk event doesn't happen, costing \$10,000. What risk response is this?
		\end{block}
		\begin{enumerate}[(A)]
			\item Avoidance
			\item Mitigation
			\item Acceptance
			\item Transfer
		\end{enumerate}
	\end{frame}
	\begin{frame}{D1-Q72: Answer}
		\begin{block}{Correct Answer: B}
			Taking operational controls to reduce likelihood and impact of risk is risk mitigation.
		\end{block}
	\end{frame}
	
	\begin{frame}{D1-Q73: IS Control Objective}
		\begin{block}{Question}
			What is the MOST important objective of information system control?
		\end{block}
		\begin{enumerate}[(A)]
			\item Business objectives achieved, undesired risk detected/corrected
			\item Ensuring effective and efficient operations
			\item Developing business continuity plans
			\item Safeguarding assets
		\end{enumerate}
	\end{frame}
	\begin{frame}{D1-Q73: Answer}
		\begin{block}{Correct Answer: A}
			The basic purpose of IS control is to ensure business objectives are achieved and undesired risk events are detected and corrected.
		\end{block}
	\end{frame}
	
	\begin{frame}{D1-Q74: Business Continuity Strategy}
		\begin{block}{Question}
			Which is prepared by business and serves as starting point for IT Service Continuity Strategy?
		\end{block}
		\begin{enumerate}[(A)]
			\item Business Continuity Strategy
			\item Index of Disaster-Relevant Information
			\item Disaster Invocation Guideline
			\item Testing Schedule
		\end{enumerate}
	\end{frame}
	\begin{frame}{D1-Q74: Answer}
		\begin{block}{Correct Answer: A}
			Business Continuity Strategy is outline of approach to ensure continuity of Vital Business Functions, prepared by business.
		\end{block}
	\end{frame}
	
	\begin{frame}{D1-Q75: Maturity Level 5}
		\begin{block}{Question}
			"Real-time monitoring of risk events and control exceptions exists, as does automation of policy management" describes which level?
		\end{block}
		\begin{enumerate}[(A)]
			\item Level 3
			\item Level 0
			\item Level 5
			\item Level 2
		\end{enumerate}
	\end{frame}
	\begin{frame}{D1-Q75: Answer}
		\begin{block}{Correct Answer: C}
			Risk management capability maturity level 5 has real-time monitoring and automation of policy management.
		\end{block}
	\end{frame}
	
	\begin{frame}{D1-Q76: Cost Performance Index}
		\begin{block}{Question}
			Which is true for Cost Performance Index (CPI)?
		\end{block}
		\begin{enumerate}[(A)]
			\item If CPI greater than 1, it indicates better than expected performance
			\item CPI = EV × AC
			\item CPI is used to measure schedule performance
			\item If CPI = 1, it indicates poor performance
		\end{enumerate}
	\end{frame}
	\begin{frame}{D1-Q76: Answer}
		\begin{block}{Correct Answer: A}
			CPI = EV/AC. If CPI greater than 1, indicates better than expected performance. CPI = 1 indicates right on target.
		\end{block}
	\end{frame}
	
	\begin{frame}{D1-Q77: NOT Indirect Information}
		\begin{block}{Question}
			Which does NOT provide indirect information?
		\end{block}
		\begin{enumerate}[(A)]
			\item Information about propriety of cutoff
			\item Reports of orders rejected for credit limitations
			\item Reports of unusual deviations
			\item No significant differences between perpetual and actual levels
		\end{enumerate}
	\end{frame}
	\begin{frame}{D1-Q77: Answer}
		\begin{block}{Correct Answer: A}
			Information about propriety of cutoff is direct information, not indirect.
		\end{block}
	\end{frame}
	
	\begin{frame}{D1-Q78: Salience Model}
		\begin{block}{Question}
			Which activity best describes a salience model?
		\end{block}
		\begin{enumerate}[(A)]
			\item Power, urgency, and legitimacy
			\item Power and interest
			\item Influence and impact
			\item Power and influence
		\end{enumerate}
	\end{frame}
	\begin{frame}{D1-Q78: Answer}
		\begin{block}{Correct Answer: A}
			Salience model describes stakeholders based on power, urgency, and legitimacy.
		\end{block}
	\end{frame}
	
	\begin{frame}{D1-Q79: First Step in Risk Assessment}
		\begin{block}{Question}
			What is the first step in the risk assessment process?
		\end{block}
		\begin{enumerate}[(A)]
			\item Identification of assets
			\item Identification of threats
			\item Identification of threat sources
			\item Identification of vulnerabilities
		\end{enumerate}
	\end{frame}
	\begin{frame}{D1-Q79: Answer}
		\begin{block}{Correct Answer: A}
			Asset identification is the most crucial and first step in the risk assessment process.
		\end{block}
	\end{frame}
	
	\begin{frame}{D1-Q80: Risk Threshold Matrix}
		\begin{block}{Question}
			Which matrix is used to specify risk thresholds?
		\end{block}
		\begin{enumerate}[(A)]
			\item Risk indicator matrix
			\item Impact matrix
			\item Risk scenario matrix
			\item Probability matrix
		\end{enumerate}
	\end{frame}
	\begin{frame}{D1-Q80: Answer}
		\begin{block}{Correct Answer: A}
			Risk indicators are metrics used to indicate risk thresholds - when risk level is approaching high or unacceptable level.
		\end{block}
	\end{frame}
	
	\begin{frame}{D1-Q81: Risk Appetite Factors}
		\begin{block}{Question}
			What are the two MAJOR factors for deciding risk appetite? (Choose two)
		\end{block}
		\begin{enumerate}[(A)]
			\item Amount of loss enterprise wants to accept
			\item Alignment with risk-culture
			\item Risk-aware decisions
			\item Capacity to absorb loss
		\end{enumerate}
	\end{frame}
	\begin{frame}{D1-Q81: Answer}
		\begin{block}{Correct Answer: A and D}
			Two major factors: (1) capacity to absorb loss, (2) culture towards risk taking - amount of loss enterprise wants to accept.
		\end{block}
	\end{frame}
	
	\begin{frame}{D1-Q82: Risk Monitoring Process}
		\begin{block}{Question}
			Which process is responsible for identifying new risks, changing risks, and outdated risks?
		\end{block}
		\begin{enumerate}[(A)]
			\item Risk planning
			\item Risk monitoring and controlling
			\item Risk identification
			\item Risk analysis
		\end{enumerate}
	\end{frame}
	\begin{frame}{D1-Q82: Answer}
		\begin{block}{Correct Answer: B}
			Risk monitoring and controlling is responsible for identifying new risks, determining status of changed risks, and outdated risks.
		\end{block}
	\end{frame}
	
	\begin{frame}{D1-Q83: SLE Calculation}
		\begin{block}{Question}
			Asset valued at \$125,000 subjected to exposure factor of 25\%. What is Single Loss Expectancy?
		\end{block}
		\begin{enumerate}[(A)]
			\item \$125,025
			\item \$31,250
			\item \$5,000
			\item \$3,125,000
		\end{enumerate}
	\end{frame}
	\begin{frame}{D1-Q83: Answer}
		\begin{block}{Correct Answer: B}
			SLE = Asset Value × Exposure Factor = \$125,000 × 0.25 = \$31,250
		\end{block}
	\end{frame}
	
	\begin{frame}{D1-Q84: Access Control Principles}
		\begin{block}{Question}
			Which are principles of access controls? (Choose three)
		\end{block}
		\begin{enumerate}[(A)]
			\item Confidentiality
			\item Availability
			\item Reliability
			\item Integrity
		\end{enumerate}
	\end{frame}
	\begin{frame}{D1-Q84: Answer}
		\begin{block}{Correct Answer: A, B, D}
			Access controls focus on Confidentiality, Integrity, and Availability (CIA triad).
		\end{block}
	\end{frame}
	
	\begin{frame}{D1-Q85: Maintaining KRIs Reason}
		\begin{block}{Question}
			What is the MOST important reason to maintain Key Risk Indicators?
		\end{block}
		\begin{enumerate}[(A)]
			\item Risk reports need to be timely
			\item Complex metrics require fine-tuning
			\item Threats and vulnerabilities change over time
			\item They help to avoid risk
		\end{enumerate}
	\end{frame}
	\begin{frame}{D1-Q85: Answer}
		\begin{block}{Correct Answer: C}
			Since internal and external environments constantly change, KRIs need to effectively capture these changes.
		\end{block}
	\end{frame}
	
	\begin{frame}{D1-Q86: NOT Technical Control}
		\begin{block}{Question}
			Which control does NOT come under technical class of control?
		\end{block}
		\begin{enumerate}[(A)]
			\item Program management control
			\item System and Communications Protection control
			\item Identification and Authentication control
			\item Access Control
		\end{enumerate}
	\end{frame}
	\begin{frame}{D1-Q86: Answer}
		\begin{block}{Correct Answer: A}
			Program Management control comes under management class of controls, not technical.
		\end{block}
	\end{frame}
	
	\begin{frame}{D1-Q87: Brainstorming Method}
		\begin{block}{Question}
			Mary uses a facilitator to help generate ideas about project risks. What method is she using?
		\end{block}
		\begin{enumerate}[(A)]
			\item Delphi Techniques
			\item Expert judgment
			\item Brainstorming
			\item Checklist analysis
		\end{enumerate}
	\end{frame}
	\begin{frame}{D1-Q87: Answer}
		\begin{block}{Correct Answer: C}
			Brainstorming attempts to create a comprehensive list of risks and is often led by a moderator or facilitator.
		\end{block}
	\end{frame}
	
	\begin{frame}{D1-Q88: Administrative Control}
		\begin{block}{Question}
			Which is an administrative control?
		\end{block}
		\begin{enumerate}[(A)]
			\item Water detection
			\item Reasonableness check
			\item Data loss prevention program
			\item Session timeout
		\end{enumerate}
	\end{frame}
	\begin{frame}{D1-Q88: Answer}
		\begin{block}{Correct Answer: C}
			Data Loss Prevention program is an administrative control (program consists of projects, policies).
		\end{block}
	\end{frame}
	
	\begin{frame}{D1-Q89: Plan Creation}
		\begin{block}{Question}
			What document defines how risks will be identified, quantified, and contingency plans implemented?
		\end{block}
		\begin{enumerate}[(A)]
			\item Project plan
			\item Resource management plan
			\item Project management plan
			\item Risk management plan
		\end{enumerate}
	\end{frame}
	\begin{frame}{D1-Q89: Answer}
		\begin{block}{Correct Answer: D}
			Risk management plan defines how risks will be identified, analyzed, monitored, controlled, and responded to.
		\end{block}
	\end{frame}
	
	\begin{frame}{D1-Q90: Where Risks Documented}
		\begin{block}{Question}
			Where are all risks and risk responses documented as project progresses?
		\end{block}
		\begin{enumerate}[(A)]
			\item Risk management plan
			\item Project management plan
			\item Risk response plan
			\item Risk register
		\end{enumerate}
	\end{frame}
	\begin{frame}{D1-Q90: Answer}
		\begin{block}{Correct Answer: D}
			All risks, their responses, and other characteristics are documented in the risk register.
		\end{block}
	\end{frame}
	
	\begin{frame}{D1-Q91: Hiring Third Party}
		\begin{block}{Question}
			Hiring a company to deal with all hardware work is which type of risk response?
		\end{block}
		\begin{enumerate}[(A)]
			\item Transference
			\item Mitigation
			\item Avoidance
			\item Exploit
		\end{enumerate}
	\end{frame}
	\begin{frame}{D1-Q91: Answer}
		\begin{block}{Correct Answer: A}
			Hiring a third party to own risk is risk transference.
		\end{block}
	\end{frame}
	
	\begin{frame}{D1-Q92: Activity Duration Estimates}
		\begin{block}{Question}
			Which input of identify risks process helps identify risks associated with time allowances?
		\end{block}
		\begin{enumerate}[(A)]
			\item Activity duration estimates
			\item Activity cost estimates
			\item Risk management plan
			\item Schedule management plan
		\end{enumerate}
	\end{frame}
	\begin{frame}{D1-Q92: Answer}
		\begin{block}{Correct Answer: A}
			Activity duration estimates review helps identify risks associated with time allowances.
		\end{block}
	\end{frame}
	
	\begin{frame}{D1-Q93: Loss of Integrity Events}
		\begin{block}{Question}
			Which events refer to loss of integrity? (Choose three)
		\end{block}
		\begin{enumerate}[(A)]
			\item Someone sees company's secret formula
			\item Someone makes unauthorized changes to a Web site
			\item An e-mail message is modified in transit
			\item A virus infects a file
		\end{enumerate}
	\end{frame}
	\begin{frame}{D1-Q93: Answer}
		\begin{block}{Correct Answer: B, C, D}
			Loss of integrity includes: message modified in transit, virus infects file, unauthorized Web site changes.
		\end{block}
	\end{frame}
	
	\begin{frame}{D1-Q94: First Step IS Control Design}
		\begin{block}{Question}
			What should be PRIMARILY considered while designing information systems controls?
		\end{block}
		\begin{enumerate}[(A)]
			\item The IT strategic plan
			\item The existing IT environment
			\item The organizational strategic plan
			\item The present IT budget
		\end{enumerate}
	\end{frame}
	\begin{frame}{D1-Q94: Answer}
		\begin{block}{Correct Answer: C}
			Review of enterprise's strategic plan is the first step in designing effective IS controls.
		\end{block}
	\end{frame}
	
	\begin{frame}{D1-Q95: Communication Inhibitor}
		\begin{block}{Question}
			Which is the MOST effective inhibitor of relevant and efficient communication?
		\end{block}
		\begin{enumerate}[(A)]
			\item False sense of confidence at top
			\item Perception of covering up known risk
			\item Existence of a blame culture
			\item Misalignment between risk appetite and policies
		\end{enumerate}
	\end{frame}
	\begin{frame}{D1-Q95: Answer}
		\begin{block}{Correct Answer: C}
			Blame culture is the most effective inhibitor. Business units point fingers at IT when projects fail.
		\end{block}
	\end{frame}
	
	\begin{frame}{D1-Q96: Low Impact Risks}
		\begin{block}{Question}
			What should you do with small risks that won't affect the project much?
		\end{block}
		\begin{enumerate}[(A)]
			\item These risks can be dismissed
			\item These risks can be accepted
			\item Added to low priority risk watch list
			\item All risks must have valid documented risk response
		\end{enumerate}
	\end{frame}
	\begin{frame}{D1-Q96: Answer}
		\begin{block}{Correct Answer: C}
			Low-impact, low-probability risks can be added to the low priority risk watch list.
		\end{block}
	\end{frame}
	
	\begin{frame}{D1-Q97: IDS Control Type}
		\begin{block}{Question}
			What type of control is an intrusion detection system (IDS)?
		\end{block}
		\begin{enumerate}[(A)]
			\item Detective
			\item Corrective
			\item Preventative
			\item Recovery
		\end{enumerate}
	\end{frame}
	\begin{frame}{D1-Q97: Answer}
		\begin{block}{Correct Answer: A}
			IDS detects and gives warning when policy violation occurs; it is a detective control.
		\end{block}
	\end{frame}
	
	\begin{frame}{D1-Q98: Audit and Accountability}
		\begin{block}{Question}
			What are functions of audit and accountability control? (Choose three)
		\end{block}
		\begin{enumerate}[(A)]
			\item Provides details on how to protect audit logs
			\item Implement effective access control
			\item Implement an effective audit program
			\item Provides details on how to determine what to audit
		\end{enumerate}
	\end{frame}
	\begin{frame}{D1-Q98: Answer}
		\begin{block}{Correct Answer: A, C, D}
			Audit and accountability helps: implement audit program, determine what to audit, protect audit logs, non-repudiation.
		\end{block}
	\end{frame}
	
	\begin{frame}{D1-Q99: Risk Response Trigger}
		\begin{block}{Question}
			Which acts as a trigger for risk response process?
		\end{block}
		\begin{enumerate}[(A)]
			\item Risk level increases above risk appetite
			\item Risk level increases above risk tolerance
			\item Risk level equates risk appetite
			\item Risk level equates risk tolerance
		\end{enumerate}
	\end{frame}
	\begin{frame}{D1-Q99: Answer}
		\begin{block}{Correct Answer: B}
			Risk response process is triggered when a risk exceeds the enterprise's risk tolerance level.
		\end{block}
	\end{frame}
	
	\begin{frame}{D1-Q100: Exposure Factor Value}
		\begin{block}{Question}
			What is the value of exposure factor if the asset is lost completely?
		\end{block}
		\begin{enumerate}[(A)]
			\item 1
			\item Infinity
			\item 10
			\item 0
		\end{enumerate}
	\end{frame}
	\begin{frame}{D1-Q100: Answer}
		\begin{block}{Correct Answer: A}
			Exposure Factor = 1.0 when asset is completely lost (100\% loss).
		\end{block}
	\end{frame}
	
	% ============================================================
	% DOMAIN 2 --- IT RISK ASSESSMENT (20%)
	% ============================================================
	\section{Domain 2: IT Risk Assessment}
	
	\begin{frame}{D2-Q101: Annual Loss Expectancy Formula}
		\begin{block}{Question}
			An asset's annual loss expectancy is calculated as the:
		\end{block}
		\begin{enumerate}[(A)]
			\item Exposure factor (EF) multiplied by the annualized rate of occurrence (ARO)
			\item Single loss expectancy (SLE) multiplied by the EF
			\item SLE multiplied by the ARO
			\item Asset value multiplied by the SLE
		\end{enumerate}
	\end{frame}
	\begin{frame}{D2-Q101: Answer}
		\begin{block}{Correct Answer: C --- ALE = SLE × ARO}
			ALE is calculated by multiplying the single loss expectancy (SLE) by the annualized rate of occurrence (ARO).
		\end{block}
	\end{frame}
	
	\begin{frame}{D2-Q102: Effective Information Risk Program}
		\begin{block}{Question}
			Which is the BEST indicator of an effective information risk management program?
		\end{block}
		\begin{enumerate}[(A)]
			\item The security policy is made widely available
			\item Risk is considered before all decisions
			\item Security procedures are updated annually
			\item Risk assessments occur on an annual basis
		\end{enumerate}
	\end{frame}
	\begin{frame}{D2-Q102: Answer}
		\begin{block}{Correct Answer: B}
			Defining information risk in advance of business decisions best ensures that risk tolerance remains at approved levels.
		\end{block}
	\end{frame}
	
	\begin{frame}{D2-Q103: Risk Management Strategic Plans}
		\begin{block}{Question}
			Risk management strategic plans are MOST effective when developed for:
		\end{block}
		\begin{enumerate}[(A)]
			\item The enterprise as a whole
			\item Each individual system based on technology used
			\item Every location based on geographic threats
			\item End-to-end business processes
		\end{enumerate}
	\end{frame}
	\begin{frame}{D2-Q103: Answer}
		\begin{block}{Correct Answer: A}
			Risk management strategic plans are most effective when they are created and followed by the entire enterprise.
		\end{block}
	\end{frame}
	
	\begin{frame}{D2-Q104: Considering Risk Appetite}
		\begin{block}{Question}
			Which is MOST important when considering the risk appetite of an enterprise?
		\end{block}
		\begin{enumerate}[(A)]
			\item The capacity of the enterprise to absorb loss
			\item The definition of responsibilities for risk management
			\item The line of business and the typical risk of the industry
			\item The culture and predisposition toward risk taking
		\end{enumerate}
	\end{frame}
	\begin{frame}{D2-Q104: Answer}
		\begin{block}{Correct Answer: D}
			When considering risk appetite, two major factors are relevant: the management culture and the predisposition toward risk taking.
		\end{block}
	\end{frame}
	
	\begin{frame}{D2-Q105: Primary Purpose of ERM Framework}
		\begin{block}{Question}
			The PRIMARY purpose of adopting an enterprise-wide risk management framework is to:
		\end{block}
		\begin{enumerate}[(A)]
			\item Allow flexibility to adjust the risk response strategy
			\item Centralize responsibility for the maintenance of the risk response program
			\item Enable a consistent approach to risk response throughout the enterprise
			\item Avoid higher costs for risk reduction and audit strategies
		\end{enumerate}
	\end{frame}
	\begin{frame}{D2-Q105: Answer}
		\begin{block}{Correct Answer: C}
			Enabling consistent risk response is a key objective of the risk management framework.
		\end{block}
	\end{frame}
	
	\begin{frame}{D2-Q106: 15\% Project Deviation}
		\begin{block}{Question}
			Projects frequently exceed timelines or budget by nearly 10 percent; management advises a deviation of 15 percent is acceptable. This is an example of:
		\end{block}
		\begin{enumerate}[(A)]
			\item Risk avoidance
			\item Risk tolerance
			\item Risk acceptance
			\item Risk mitigation
		\end{enumerate}
	\end{frame}
	\begin{frame}{D2-Q106: Answer}
		\begin{block}{Correct Answer: B --- Risk tolerance}
			Risk tolerance is the permissible deviation from declared risk appetite levels.
		\end{block}
	\end{frame}
	
	\begin{frame}{D2-Q107: Computing Annual Loss Exposure}
		\begin{block}{Question}
			Which is MOST useful when computing annual loss exposure?
		\end{block}
		\begin{enumerate}[(A)]
			\item The cost of existing controls
			\item The number of vulnerabilities
			\item The net present value of the asset
			\item The business value of the asset
		\end{enumerate}
	\end{frame}
	\begin{frame}{D2-Q107: Answer}
		\begin{block}{Correct Answer: D}
			Annual loss exposure is a function of the value of the information asset and the impact if a given potential risk should materialize.
		\end{block}
	\end{frame}
	
	\begin{frame}{D2-Q108: GRC Tool Major Risk}
		\begin{block}{Question}
			Which is a MAJOR risk associated with the use of governance, risk, and compliance (GRC) tools?
		\end{block}
		\begin{enumerate}[(A)]
			\item Misinterpretation of the dashboard's output
			\item Incomplete or inaccurate coverage in the enterprise for GRC
			\item Obsolescence of content
			\item Complex integration of diverse requirements
		\end{enumerate}
	\end{frame}
	\begin{frame}{D2-Q108: Answer}
		\begin{block}{Correct Answer: C --- Obsolescence of content}
			A GRC application must be updated regularly with current regulations and policies.
		\end{block}
	\end{frame}
	
	\begin{frame}{D2-Q109: Need to Review Risk Practices}
		\begin{block}{Question}
			Which signifies the need to review an enterprise's risk practices?
		\end{block}
		\begin{enumerate}[(A)]
			\item Business owners regularly challenge risk assessment findings
			\item Manufacturing assigns its own internal risk management roles
			\item The finance department finds exceptions during its yearly risk review
			\item Sales department risk procedures were last reviewed 11 months ago
		\end{enumerate}
	\end{frame}
	\begin{frame}{D2-Q109: Answer}
		\begin{block}{Correct Answer: A}
			An enterprise's risk management practices must be clearly understood and supported by business stakeholders.
		\end{block}
	\end{frame}
	
	\begin{frame}{D2-Q110: Correct Info Not Received in Time}
		\begin{block}{Question}
			The correct information was not received by the necessary recipients in a suitable time to allow proper action. This can be categorized as:
		\end{block}
		\begin{enumerate}[(A)]
			\item Integrity risk
			\item Availability risk
			\item Access risk
			\item Relevance risk
		\end{enumerate}
	\end{frame}
	\begin{frame}{D2-Q110: Answer}
		\begin{block}{Correct Answer: D --- Relevance risk}
			Relevance risk is a composite form of business risk, requiring both integrity and availability to be addressed.
		\end{block}
	\end{frame}
	
	\begin{frame}{D2-Q111: Justifying Security Program}
		\begin{block}{Question}
			Which risk assessment output is MOST suitable to help justify an enterprise information security program?
		\end{block}
		\begin{enumerate}[(A)]
			\item An inventory of risk that may impact the enterprise
			\item Documented threats to the enterprise
			\item Evaluation of the consequences
			\item A list of appropriate controls for addressing risk
		\end{enumerate}
	\end{frame}
	\begin{frame}{D2-Q111: Answer}
		\begin{block}{Correct Answer: D}
			A list of information security controls corresponding to risk scenarios is one of the primary deliverables of the risk assessment exercise.
		\end{block}
	\end{frame}
	
	\begin{frame}{D2-Q112: Ultimate Accountability for Risk}
		\begin{block}{Question}
			Accountability for risk ultimately belongs to the:
		\end{block}
		\begin{enumerate}[(A)]
			\item Chief risk officer
			\item Compliance officer
			\item Chief financial officer
			\item Board of directors
		\end{enumerate}
	\end{frame}
	\begin{frame}{D2-Q112: Answer}
		\begin{block}{Correct Answer: D --- Board of directors}
			The board of directors has ultimate accountability to shareholders, customers, employees, and the general public.
		\end{block}
	\end{frame}
	
	\begin{frame}{D2-Q113: Outdated IT Standards}
		\begin{block}{Question}
			A risk assessment reveals that many corporate IT standards have not been updated. The BEST course of action is to:
		\end{block}
		\begin{enumerate}[(A)]
			\item Review the standards against current requirements and determine their adequacy
			\item Determine that the standards should be updated annually
			\item Report that IT standards are adequate and do not need updating
			\item Review the IT policy document to see how frequently standards should be updated
		\end{enumerate}
	\end{frame}
	\begin{frame}{D2-Q113: Answer}
		\begin{block}{Correct Answer: A}
			The risk practitioner should verify that the standards are still adequate; standards that are lacking should be updated.
		\end{block}
	\end{frame}
	
	\begin{frame}{D2-Q114: Accountable for Third-Party App Risk}
		\begin{block}{Question}
			Marketing procures a third-party application found to pose data privacy regulatory risk in some regions. Who will be accountable for the risk if implemented globally?
		\end{block}
		\begin{enumerate}[(A)]
			\item The IT department
			\item The data privacy officer
			\item The chief risk officer
			\item The marketing department
		\end{enumerate}
	\end{frame}
	\begin{frame}{D2-Q114: Answer}
		\begin{block}{Correct Answer: D --- Marketing department}
			Marketing is the business owner of the application and therefore must be accountable.
		\end{block}
	\end{frame}
	
	\begin{frame}{D2-Q115: Owns Risk Treatment Decisions}
		\begin{block}{Question}
			A risk practitioner is conducting a risk assessment of the local office's network infrastructure. Who owns the risk treatment decisions?
		\end{block}
		\begin{enumerate}[(A)]
			\item Control owner
			\item Chief information officer
			\item IT security manager
			\item Senior management
		\end{enumerate}
	\end{frame}
	\begin{frame}{D2-Q115: Answer}
		\begin{block}{Correct Answer: D --- Senior management}
			Senior management ultimately owns the risk and would make risk treatment decisions.
		\end{block}
	\end{frame}
	
	\begin{frame}{D2-Q116: Responsible for Third-Party App Risk}
		\begin{block}{Question}
			Marketing procures a third-party application that poses data privacy regulatory risk. Which role will be responsible for the risk it poses to the business if implemented globally?
		\end{block}
		\begin{enumerate}[(A)]
			\item The marketing department
			\item The IT department
			\item The data privacy officer
			\item The chief risk officer
		\end{enumerate}
	\end{frame}
	\begin{frame}{D2-Q116: Answer}
		\begin{block}{Correct Answer: B --- IT department}
			The IT department is responsible for the risk posed by this application.
		\end{block}
	\end{frame}
	
	\begin{frame}{D2-Q117: Supporting Encryption Recommendation}
		\begin{block}{Question}
			Which provides the GREATEST support to a risk practitioner recommending encryption of corporate laptops and removable media as a risk mitigation measure?
		\end{block}
		\begin{enumerate}[(A)]
			\item Benchmarking with peers
			\item Evaluating public reports on encryption algorithms
			\item Developing a business case
			\item Scanning unencrypted systems for vulnerabilities
		\end{enumerate}
	\end{frame}
	\begin{frame}{D2-Q117: Answer}
		\begin{block}{Correct Answer: C --- Developing a business case}
			A business case documents how encryption addresses risk and illustrates that risk losses can be reduced by recommended controls.
		\end{block}
	\end{frame}
	
	\begin{frame}{D2-Q118: Additional Stakeholder for Privacy Risk}
		\begin{block}{Question}
			Marketing procures a third-party application that poses data privacy regulatory risk in certain regions. Which additional stakeholder should be included in reporting the risk?
		\end{block}
		\begin{enumerate}[(A)]
			\item Chief privacy officer
			\item Third-party vendor
			\item Business continuity manager
			\item Audit manager
		\end{enumerate}
	\end{frame}
	\begin{frame}{D2-Q118: Answer}
		\begin{block}{Correct Answer: A --- Chief privacy officer}
			The chief privacy officer supports defining controls that will be implemented to address and mitigate data transfer and data privacy requirements.
		\end{block}
	\end{frame}
	
	\begin{frame}{D2-Q119: Cloud Email Retention Contract}
		\begin{block}{Question}
			A company learned its cloud email provider was never contracted to back up messages, despite a three-year retention policy. Which would have BEST safeguarded the company?
		\end{block}
		\begin{enumerate}[(A)]
			\item Providing the contractor with the record retention policy up front
			\item Validating the company policies to the provider's contract
			\item Providing the contractor with the email retention policy up front
			\item Backing up the data on the company's internal network nightly
		\end{enumerate}
	\end{frame}
	\begin{frame}{D2-Q119: Answer}
		\begin{block}{Correct Answer: B}
			The initial review of third-party services should confirm that vendors are contractually required to enforce all internal policies.
		\end{block}
	\end{frame}
	
	\begin{frame}{D2-Q120: No Policy for Personal Devices}
		\begin{block}{Question}
			If an enterprise does not have a formal policy regarding personal devices in the workplace, which should the risk practitioner recommend?
		\end{block}
		\begin{enumerate}[(A)]
			\item Introduce remote wipe functionality
			\item Implement an exception process based on appropriate approvals
			\item Update incident response procedures
			\item Create and maintain an inventory of personal devices
		\end{enumerate}
	\end{frame}
	\begin{frame}{D2-Q120: Answer}
		\begin{block}{Correct Answer: B}
			If employees can attach personal devices without a formal policy, employees should be required to gain approval from management.
		\end{block}
	\end{frame}
	
	\begin{frame}{D2-Q121: Applying Risk Management}
		\begin{block}{Question}
			To be effective, risk management should be applied to:
		\end{block}
		\begin{enumerate}[(A)]
			\item The elements identified by a risk assessment
			\item Any area that exceeds acceptable risk levels
			\item All enterprise activities
			\item Only those areas that have potential impact
		\end{enumerate}
	\end{frame}
	\begin{frame}{D2-Q121: Answer}
		\begin{block}{Correct Answer: C}
			Although not all enterprise activities entail unacceptable risk, the practice of risk management is ideally applied to all enterprise activities.
		\end{block}
	\end{frame}
	
	\begin{frame}{D2-Q122: Basis for Corporate Security Policy}
		\begin{block}{Question}
			Development of corporate information security policy should PRIMARILY be based on:
		\end{block}
		\begin{enumerate}[(A)]
			\item Vulnerabilities
			\item Threats
			\item Assets
			\item Impacts
		\end{enumerate}
	\end{frame}
	\begin{frame}{D2-Q122: Answer}
		\begin{block}{Correct Answer: C --- Assets}
			The corporate information security policy is based on management's commitment to protect the assets of the enterprise.
		\end{block}
	\end{frame}
	
	\begin{frame}{D2-Q123: Risk-Aware Business Decisions}
		\begin{block}{Question}
			Which MOST enables risk-aware business decisions?
		\end{block}
		\begin{enumerate}[(A)]
			\item Robust information security policies
			\item Exchange of accurate and timely information
			\item Skilled risk management personnel
			\item Skilled risk management personnel
		\end{enumerate}
	\end{frame}
	\begin{frame}{D2-Q123: Answer}
		\begin{block}{Correct Answer: B}
			Robust exchange of information enables management to optimize risk-related decisions.
		\end{block}
	\end{frame}
	
	\begin{frame}{D2-Q124: Initial Stages of Risk Program}
		\begin{block}{Question}
			Which requirement MUST be met during the initial stages of developing a risk management program?
		\end{block}
		\begin{enumerate}[(A)]
			\item Management establishes ownership of identified risk
			\item Information security policies and standards are established
			\item A management committee exists to provide program oversight
			\item The context and purpose of the program are defined
		\end{enumerate}
	\end{frame}
	\begin{frame}{D2-Q124: Answer}
		\begin{block}{Correct Answer: D}
			Initial requirements to determine the enterprise's purpose for creating an information security risk management program include determining the desired outcomes and defining objectives.
		\end{block}
	\end{frame}
	
	\begin{frame}{D2-Q125: Effective BYOD Policy}
		\begin{block}{Question}
			An enterprise implemented a corporate BYOD policy to reduce the risk of data leakage. Which approach MOST enables the policy to be effective?
		\end{block}
		\begin{enumerate}[(A)]
			\item Obtaining signed acceptance from users
			\item Educating users on acceptable and unacceptable practices
			\item Requiring users to read the policy and future updates
			\item Clearly stating disciplinary action for noncompliance
		\end{enumerate}
	\end{frame}
	\begin{frame}{D2-Q125: Answer}
		\begin{block}{Correct Answer: B}
			A BYOD policy that is approved and communicated will not be effective without proper training of all users.
		\end{block}
	\end{frame}
	
	\begin{frame}{D2-Q126: New Regulation First Action}
		\begin{block}{Question}
			A new regulation for safeguarding information processed by a specific transaction type comes to an IT manager's attention. The manager should FIRST:
		\end{block}
		\begin{enumerate}[(A)]
			\item Meet with stakeholders to decide how to comply
			\item Analyze the key risk in the compliance process
			\item Update the existing security/privacy policy
			\item Assess whether existing controls meet the regulation
		\end{enumerate}
	\end{frame}
	\begin{frame}{D2-Q126: Answer}
		\begin{block}{Correct Answer: D}
			The first step is to understand the impact and requirements of the new regulation, which includes assessing how the enterprise will comply and to what extent existing controls support compliance.
		\end{block}
	\end{frame}
	
	\begin{frame}{D2-Q127: Document Identifying Compliance Risk}
		\begin{block}{Question}
			Which document BEST identifies an enterprise's compliance risk and corrective actions in progress to meet regulatory requirements?
		\end{block}
		\begin{enumerate}[(A)]
			\item An internal audit report
			\item A risk register
			\item An external audit report
			\item A risk assessment report
		\end{enumerate}
	\end{frame}
	\begin{frame}{D2-Q127: Answer}
		\begin{block}{Correct Answer: B --- Risk register}
			A risk register provides a report of all current identified risk within an enterprise (including compliance risk) and tracks the status of corrective actions.
		\end{block}
	\end{frame}
	
	\begin{frame}{D2-Q128: Impact on Security Governance Models}
		\begin{block}{Question}
			Which will have the MOST significant impact on standard information security governance models?
		\end{block}
		\begin{enumerate}[(A)]
			\item Number of employees and consulting staff
			\item Cultural differences between physical locations
			\item Complexity of the organizational structure
			\item Uncertainty in legislative requirements
		\end{enumerate}
	\end{frame}
	\begin{frame}{D2-Q128: Answer}
		\begin{block}{Correct Answer: C}
			Information security governance models are highly dependent on the complexity of the organizational structure.
		\end{block}
	\end{frame}
	
	\begin{frame}{D2-Q129: Highest Tolerance for Public Cloud}
		\begin{block}{Question}
			Senior management will MOST likely have the highest tolerance for moving which to a public cloud?
		\end{block}
		\begin{enumerate}[(A)]
			\item Credit-card processing
			\item Research and development
			\item The legacy financial system
			\item The corporate email system
		\end{enumerate}
	\end{frame}
	\begin{frame}{D2-Q129: Answer}
		\begin{block}{Correct Answer: D}
			Of the options, the corporate email system has the least competitive distinction, complexity, and sensitive/highly classified information.
		\end{block}
	\end{frame}
	
	\begin{frame}{D2-Q130: Risk Profile Consideration}
		\begin{block}{Question}
			Which item is MOST important to consider in relation to a risk profile?
		\end{block}
		\begin{enumerate}[(A)]
			\item A summary of regional loss events
			\item Aggregated risk to the enterprise
			\item A description of critical risk
			\item An analysis of historical loss events
		\end{enumerate}
	\end{frame}
	\begin{frame}{D2-Q130: Answer}
		\begin{block}{Correct Answer: B}
			The risk profile is based on the aggregated risk to the enterprise, including historical risk, critical risk, and emerging risk.
		\end{block}
	\end{frame}
	
	\begin{frame}{D2-Q131: Collaborative Research Cloud}
		\begin{block}{Question}
			Which cloud deployment model is MOST appropriate for a collaborative research program between universities?
		\end{block}
		\begin{enumerate}[(A)]
			\item Private cloud
			\item Community cloud
			\item Public cloud
			\item Hybrid cloud
		\end{enumerate}
	\end{frame}
	\begin{frame}{D2-Q131: Answer}
		\begin{block}{Correct Answer: B --- Community cloud}
			A community cloud is shared between several enterprises and supports a specific community with a common mission or interest.
		\end{block}
	\end{frame}
	
	\begin{frame}{D2-Q132: Control Environment Review Line}
		\begin{block}{Question}
			Which line of defense is MOST likely involved when conducting a review of the adequacy of the control environment?
		\end{block}
		\begin{enumerate}[(A)]
			\item The first line of defense
			\item The second line of defense
			\item The third line of defense
			\item Third-party line of defense
		\end{enumerate}
	\end{frame}
	\begin{frame}{D2-Q132: Answer}
		\begin{block}{Correct Answer: A --- First line}
			A control environment review is a first line of defense activity.
		\end{block}
	\end{frame}
	
	\begin{frame}{D2-Q133: Board Concern on Risk Capability}
		\begin{block}{Question}
			Which issue would be of MOST concern to the board when assessing risk management capability?
		\end{block}
		\begin{enumerate}[(A)]
			\item The third line acts independently of the other lines
			\item The second line is checking and challenging the first line
			\item Internal audit is preparing the risk management strategy
			\item Each line of defense is conducting its own planning independently
		\end{enumerate}
	\end{frame}
	\begin{frame}{D2-Q133: Answer}
		\begin{block}{Correct Answer: D}
			Joint planning across the three lines of defense is key to achieving an effective risk management capability.
		\end{block}
	\end{frame}
	
	\begin{frame}{D2-Q134: Supporting Code of Ethics}
		\begin{block}{Question}
			Which MOST effectively supports adherence to an enterprise's code of ethics?
		\end{block}
		\begin{enumerate}[(A)]
			\item Ensuring periodic training, evaluation, and attestation of employees
			\item Performing background checks at hire
			\item Providing clear enterprise policies and standards
			\item Providing continuous awareness of ethical requirements
		\end{enumerate}
	\end{frame}
	\begin{frame}{D2-Q134: Answer}
		\begin{block}{Correct Answer: A}
			Attestation to comply with the code of conduct during ethics training ensures employees have a clear idea of what is expected.
		\end{block}
	\end{frame}
	
	\begin{frame}{D2-Q135: Explaining Zero-Day Exploit}
		\begin{block}{Question}
			Who is responsible for explaining the ramifications of a new zero-day exploit to senior management?
		\end{block}
		\begin{enumerate}[(A)]
			\item Chief operating officer
			\item Chief risk officer
			\item Chief information security officer
			\item Chief information officer
		\end{enumerate}
	\end{frame}
	\begin{frame}{D2-Q135: Answer}
		\begin{block}{Correct Answer: B --- Chief risk officer}
			The chief risk officer is the most senior official accountable for all aspects of risk management across the enterprise.
		\end{block}
	\end{frame}
	
	\begin{frame}{D2-Q136: Information Classification Responsibility}
		\begin{block}{Question}
			The application of information classification is the responsibility of the:
		\end{block}
		\begin{enumerate}[(A)]
			\item Information security officer
			\item Information owner
			\item Information systems auditor
			\item Information custodian
		\end{enumerate}
	\end{frame}
	\begin{frame}{D2-Q136: Answer}
		\begin{block}{Correct Answer: B --- Information owner}
			The information owner determines classification based on the criticality and sensitivity of information.
		\end{block}
	\end{frame}
	
	\begin{frame}{D2-Q137: Three Lines of Defense Benefit}
		\begin{block}{Question}
			The MOST significant benefit of using the three lines of defense model in a risk management framework is that it:
		\end{block}
		\begin{enumerate}[(A)]
			\item Ensures ongoing success of the risk management initiative
			\item Enhances communication between various stakeholders
			\item Clarifies essential roles of the key stakeholders
			\item Helps risk owners in selecting appropriate control owners
		\end{enumerate}
	\end{frame}
	\begin{frame}{D2-Q137: Answer}
		\begin{block}{Correct Answer: C}
			The three lines of defense model is designed to assign roles related to risk, controls, and security while better defining the relationships between them.
		\end{block}
	\end{frame}
	
	\begin{frame}{D2-Q138: First Line of Defense Role}
		\begin{block}{Question}
			Which risk management role is part of the first line of defense?
		\end{block}
		\begin{enumerate}[(A)]
			\item Chief risk officer
			\item Risk steering committee
			\item Risk owner
			\item Board of directors
		\end{enumerate}
	\end{frame}
	\begin{frame}{D2-Q138: Answer}
		\begin{block}{Correct Answer: C --- Risk owner}
			The first line of defense is operational management, and risk owners are part of operational management.
		\end{block}
	\end{frame}
	
	\begin{frame}{D2-Q139: Ultimate Objective of Risk Governance}
		\begin{block}{Question}
			What is the ULTIMATE objective of risk governance?
		\end{block}
		\begin{enumerate}[(A)]
			\item Benefits realization
			\item Risk optimization
			\item Resource optimization
			\item Value creation
		\end{enumerate}
	\end{frame}
	\begin{frame}{D2-Q139: Answer}
		\begin{block}{Correct Answer: D --- Value creation}
			Value creation is the main objective of risk governance and is achieved when the three underlying objectives are balanced.
		\end{block}
	\end{frame}
	
	\begin{frame}{D2-Q140: Second Line of Defense Activity}
		\begin{block}{Question}
			Which activity would be a direct responsibility of the second line of defense?
		\end{block}
		\begin{enumerate}[(A)]
			\item Establish a risk management framework
			\item Identify, assess, and respond to risk
			\item Implement information security controls
			\item Test and report control effectiveness
		\end{enumerate}
	\end{frame}
	\begin{frame}{D2-Q140: Answer}
		\begin{block}{Correct Answer: A}
			Establishing a risk management framework, providing awareness training, and supervising overall risk management are responsibilities of the second line of defense.
		\end{block}
	\end{frame}
	
	\begin{frame}{D2-Q141: Ethical Behavior Contribution}
		\begin{block}{Question}
			Developing and practicing ethical behavior within an enterprise contributes the most to:
		\end{block}
		\begin{enumerate}[(A)]
			\item Determining the risk appetite
			\item Building the risk culture
			\item Enhancing risk management
			\item Increasing risk tolerance
		\end{enumerate}
	\end{frame}
	\begin{frame}{D2-Q141: Answer}
		\begin{block}{Correct Answer: B --- Building the risk culture}
			Practicing professional ethics within an enterprise helps stakeholders review risk associated with unethical behavior.
		\end{block}
	\end{frame}
	
	\begin{frame}{D2-Q142: Benefits of Professional Ethics}
		\begin{block}{Question}
			Which is the MOST important benefit when stakeholders promote and follow professional ethics?
		\end{block}
		\begin{enumerate}[(A)]
			\item It increases reliability in risk management processes
			\item It increases risk management effectiveness within an enterprise
			\item It improves compliance with regulatory requirements
			\item It improves the enterprise's reputation and business growth
		\end{enumerate}
	\end{frame}
	\begin{frame}{D2-Q142: Answer}
		\begin{block}{Correct Answer: B}
			When stakeholders follow ethical practices, actions and decisions are based on reviewing the results and impact of those actions.
		\end{block}
	\end{frame}
	
	\begin{frame}{D2-Q143: Best Proactive Ethics Approach}
		\begin{block}{Question}
			Which is the BEST proactive approach for practicing professional ethics within an enterprise?
		\end{block}
		\begin{enumerate}[(A)]
			\item Develop an ethics policy
			\item Implement processes to report unethical behavior
			\item Define procedures for disciplinary action
			\item Provide ethics awareness training
		\end{enumerate}
	\end{frame}
	\begin{frame}{D2-Q143: Answer}
		\begin{block}{Correct Answer: D}
			Awareness training is the most proactive approach to improve professional ethics within an enterprise.
		\end{block}
	\end{frame}
	
	\begin{frame}{D2-Q144: Ensuring 3LoD Effectiveness}
		\begin{block}{Question}
			Which is the MOST critical to ensure the three lines of defense work together effectively?
		\end{block}
		\begin{enumerate}[(A)]
			\item A documented disaster recovery plan
			\item Automated operational processes
			\item Independent external audits
			\item Clear, detailed roles and responsibilities
		\end{enumerate}
	\end{frame}
	\begin{frame}{D2-Q144: Answer}
		\begin{block}{Correct Answer: D}
			Without clear, detailed job descriptions for employees working in each line of defense, there can be confusion on what is expected at each level.
		\end{block}
	\end{frame}
	
	\begin{frame}{D2-Q145: What Should Drive IT Plan}
		\begin{block}{Question}
			Which should drive the IT plan?
		\end{block}
		\begin{enumerate}[(A)]
			\item Strategic planning and business requirements
			\item Technology and operational procedures
			\item Compliance with laws and regulations
			\item Project plans and stakeholder requirements
		\end{enumerate}
	\end{frame}
	\begin{frame}{D2-Q145: Answer}
		\begin{block}{Correct Answer: A}
			IT exists to support business objectives. Management of enterprise IT should align the IT plan closely with the business.
		\end{block}
	\end{frame}
	
	\begin{frame}{D2-Q146: Quantitative Risk Analysis Drawback}
		\begin{block}{Question}
			MOST significant drawback of using quantitative risk analysis instead of qualitative risk analysis is:
		\end{block}
		\begin{enumerate}[(A)]
			\item Lower objectivity
			\item Greater reliance on expertise
			\item Less management buy-in
			\item Higher cost
		\end{enumerate}
	\end{frame}
	\begin{frame}{D2-Q146: Answer}
		\begin{block}{Correct Answer: D}
			Quantitative risk analysis is generally more complex and, therefore, more costly than qualitative risk analysis.
		\end{block}
	\end{frame}
	
	\begin{frame}{D2-Q147: Risk Scenarios Analysis Purpose}
		\begin{block}{Question}
			Risk scenarios are analyzed to determine:
		\end{block}
		\begin{enumerate}[(A)]
			\item Strength of controls
			\item Likelihood and impact
			\item Current risk profile
			\item Scenario root cause
		\end{enumerate}
	\end{frame}
	\begin{frame}{D2-Q147: Answer}
		\begin{block}{Correct Answer: B}
			Risk scenarios are descriptions of events that can lead to a business impact and are evaluated to determine the likelihood and impact should the event occur.
		\end{block}
	\end{frame}
	
	\begin{frame}{D2-Q148: Risk Ownership of Critical System}
		\begin{block}{Question}
			The risk to an information system that supports a critical business process is owned by:
		\end{block}
		\begin{enumerate}[(A)]
			\item IT director
			\item Senior management
			\item Risk management department
			\item System users
		\end{enumerate}
	\end{frame}
	\begin{frame}{D2-Q148: Answer}
		\begin{block}{Correct Answer: B}
			Senior management is responsible for the acceptance and mitigation of all risk.
		\end{block}
	\end{frame}
	
	\begin{frame}{D2-Q149: Periodic Risk Assessment Reason}
		\begin{block}{Question}
			PRIMARY reason risk assessments should be repeated at regular intervals is:
		\end{block}
		\begin{enumerate}[(A)]
			\item Omissions in earlier assessments can be addressed
			\item Periodic assessments allow various methodologies
			\item Business threats are constantly changing
			\item They help raise risk awareness among staff
		\end{enumerate}
	\end{frame}
	\begin{frame}{D2-Q149: Answer}
		\begin{block}{Correct Answer: C}
			As business objectives and methods change, the nature and relevance of threats also change.
		\end{block}
	\end{frame}
	
	\begin{frame}{D2-Q150: Measuring Risk Management Development}
		\begin{block}{Question}
			Which BEST assists in measuring existing level of development of risk management processes against desired state?
		\end{block}
		\begin{enumerate}[(A)]
			\item Capability maturity model (CMM)
			\item Risk management audit reports
			\item Balanced scorecard
			\item Enterprise security architecture
		\end{enumerate}
	\end{frame}
	\begin{frame}{D2-Q150: Answer}
		\begin{block}{Correct Answer: A}
			The capability maturity model (CMM) grades processes on a scale of 0 to 5, based on their maturity.
		\end{block}
	\end{frame}
	
	\begin{frame}{D2-Q151: Identifying Control Deficiencies}
		\begin{block}{Question}
			Which BEST helps identify information systems control deficiencies?
		\end{block}
		\begin{enumerate}[(A)]
			\item Gap analysis
			\item Current IT risk profile
			\item IT controls framework
			\item Countermeasure analysis
		\end{enumerate}
	\end{frame}
	\begin{frame}{D2-Q151: Answer}
		\begin{block}{Correct Answer: A}
			Controls are deployed to achieve desired control objectives based on risk assessments and business requirements. The gap between desired and actual identifies control deficiencies.
		\end{block}
	\end{frame}
	
	\begin{frame}{D2-Q152: Statistical Methods Risk Analysis}
		\begin{block}{Question}
			Deriving likelihood and impact of risk scenarios through statistical methods is MOST LIKELY associated with which type of risk analysis?
		\end{block}
		\begin{enumerate}[(A)]
			\item Risk scenario
			\item Qualitative
			\item Quantitative
			\item Semiquantitative
		\end{enumerate}
	\end{frame}
	\begin{frame}{D2-Q152: Answer}
		\begin{block}{Correct Answer: C}
			The essence of quantitative risk assessment is to derive the likelihood and impact of risk scenarios based on statistical methods and data.
		\end{block}
	\end{frame}
	
	\begin{frame}{D2-Q153: Reviewing Risk Analysis Results}
		\begin{block}{Question}
			Which review is BEST suited for review of IT risk analysis results before sending to management?
		\end{block}
		\begin{enumerate}[(A)]
			\item Internal audit review
			\item Peer review
			\item Compliance review
			\item Risk policy review
		\end{enumerate}
	\end{frame}
	\begin{frame}{D2-Q153: Answer}
		\begin{block}{Correct Answer: B}
			It is effective, efficient and good practice to perform a peer review of IT risk analysis results before sending them to management.
		\end{block}
	\end{frame}
	
	\begin{frame}{D2-Q154: Threat and Vulnerability Assessment}
		\begin{block}{Question}
			Which uses risk scenarios when estimating the likelihood and impact of significant risk to the enterprise?
		\end{block}
		\begin{enumerate}[(A)]
			\item An IT audit
			\item A security gap analysis
			\item A threat and vulnerability assessment
			\item An IT security assessment
		\end{enumerate}
	\end{frame}
	\begin{frame}{D2-Q154: Answer}
		\begin{block}{Correct Answer: C}
			A threat and vulnerability assessment typically evaluates all elements of a business process for threats and vulnerabilities and identifies likelihood and business impact.
		\end{block}
	\end{frame}
	
	\begin{frame}{D2-Q155: Rapid Implementation Deviating from Policy}
		\begin{block}{Question}
			Because of its importance to the business, an enterprise wants to quickly implement a technical solution that deviates from company policies. The risk practitioner should:
		\end{block}
		\begin{enumerate}[(A)]
			\item Recommend against implementation because it violates policies
			\item Recommend revision of the current policy
			\item Conduct a risk assessment and allow or disallow based on the outcome
			\item Recommend a risk assessment and subsequent implementation only if residual risk is accepted
		\end{enumerate}
	\end{frame}
	\begin{frame}{D2-Q155: Answer}
		\begin{block}{Correct Answer: D}
			A risk assessment should be conducted to clarify the risk whenever the enterprise's policies cannot be followed. The solution should be implemented only if the related risk is formally accepted.
		\end{block}
	\end{frame}
	
	\begin{frame}{D2-Q156: Comprehensive Qualitative Risk Analysis}
		\begin{block}{Question}
			Which will produce comprehensive results when performing a qualitative risk analysis?
		\end{block}
		\begin{enumerate}[(A)]
			\item A vulnerability assessment
			\item Scenarios with threats and impacts
			\item The value of information assets
			\item Estimated productivity losses
		\end{enumerate}
	\end{frame}
	\begin{frame}{D2-Q156: Answer}
		\begin{block}{Correct Answer: B}
			Using a list of possible scenarios with threats and impacts will better frame the range of risk and facilitate a more informed discussion and decision.
		\end{block}
	\end{frame}
	
	\begin{frame}{D2-Q157: Removing Risk from Register}
		\begin{block}{Question}
			Risk may be removed from the risk register when the risk:
		\end{block}
		\begin{enumerate}[(A)]
			\item Has been eliminated
			\item Is transferred to the vendor
			\item Threshold is exceeded
			\item Is no longer relevant
		\end{enumerate}
	\end{frame}
	\begin{frame}{D2-Q157: Answer}
		\begin{block}{Correct Answer: D}
			If a risk is relevant, it will be listed in the risk register. Once the risk is no longer relevant, it may be removed.
		\end{block}
	\end{frame}
	
	\begin{frame}{D2-Q158: Maintaining Accurate Risk Register}
		\begin{block}{Question}
			Which is the BEST way to ensure that an accurate risk register is maintained over time?
		\end{block}
		\begin{enumerate}[(A)]
			\item Monitor key risk indicators and record the findings
			\item Publish the risk register centrally with workflow features that periodically poll risk assessors
			\item Distribute the risk register to business process owners for review and updating
			\item Use audit personnel to perform regular audits and to maintain the risk register
		\end{enumerate}
	\end{frame}
	\begin{frame}{D2-Q158: Answer}
		\begin{block}{Correct Answer: B}
			Centrally publishing the risk register and enabling periodic polling of risk assessors through workflow features will ensure accuracy of content.
		\end{block}
	\end{frame}
	
	\begin{frame}{D2-Q159: Adding Emerging Risk}
		\begin{block}{Question}
			An emerging risk should be added to the risk register by the risk practitioner when:
		\end{block}
		\begin{enumerate}[(A)]
			\item The impact of the risk can be quantified
			\item The probability of occurrence is high
			\item A competitor has added the risk to its register
			\item The activity that triggers the risk initiates
		\end{enumerate}
	\end{frame}
	\begin{frame}{D2-Q159: Answer}
		\begin{block}{Correct Answer: D}
			Risk identification usually starts when planning an activity, and risk identified at planning needs to be added to the risk register.
		\end{block}
	\end{frame}
	
	\begin{frame}{D2-Q160: Primary Objective of Risk Program}
		\begin{block}{Question}
			Which is the PRIMARY objective of a risk management program?
		\end{block}
		\begin{enumerate}[(A)]
			\item Maintain residual risk at an acceptable level
			\item Implement preventive controls for every threat
			\item Remove all identified risk
			\item Reduce inherent risk to zero
		\end{enumerate}
	\end{frame}
	\begin{frame}{D2-Q160: Answer}
		\begin{block}{Correct Answer: A}
			Ensuring that all residual risk is maintained at a level acceptable to the business is the objective of a risk management program.
		\end{block}
	\end{frame}
	
	\begin{frame}{D2-Q161: Risk Assessment Techniques Purpose}
		\begin{block}{Question}
			Risk assessment techniques should be used by a risk practitioner to:
		\end{block}
		\begin{enumerate}[(A)]
			\item Maximize the return on investment
			\item Provide documentation for auditors and regulators
			\item Justify the selection of risk mitigation strategies
			\item Quantify the risk that would otherwise be subjective
		\end{enumerate}
	\end{frame}
	\begin{frame}{D2-Q161: Answer}
		\begin{block}{Correct Answer: C}
			A risk practitioner should use risk assessment techniques to justify and implement a risk mitigation strategy as efficiently as possible.
		\end{block}
	\end{frame}
	
	\begin{frame}{D2-Q162: Printer Hard Disk Risk}
		\begin{block}{Question}
			A procurement employee notices new printer models keep a copy of all printed documents on an internal hard disk. Considering the risk of unintentionally disclosing confidential data, the employee should:
		\end{block}
		\begin{enumerate}[(A)]
			\item Proceed with the order and configure printers to wipe disks after each job
			\item Notify the security manager to conduct a risk assessment for the new equipment
			\item Seek another vendor that offers printers without hard disk drives
			\item Procure printers with hard disks and notify staff to wipe disks when decommissioning
		\end{enumerate}
	\end{frame}
	\begin{frame}{D2-Q162: Answer}
		\begin{block}{Correct Answer: B}
			Risk assessment is most appropriate because it yields risk mitigation techniques that are appropriate for enterprise risk context and appetite.
		\end{block}
	\end{frame}
	
	\begin{frame}{D2-Q163: Repeated Risk Assessments}
		\begin{block}{Question}
			Risk assessments should be repeated at regular intervals because:
		\end{block}
		\begin{enumerate}[(A)]
			\item Omissions in earlier assessments can be addressed
			\item Periodic assessments allow various methodologies
			\item Business threats are constantly changing
			\item They help raise risk awareness among staff
		\end{enumerate}
	\end{frame}
	\begin{frame}{D2-Q163: Answer}
		\begin{block}{Correct Answer: C}
			As business objectives and methods change, the nature and relevance of threats also change.
		\end{block}
	\end{frame}
	
	\begin{frame}{D2-Q164: Assessing Information Systems Risk}
		\begin{block}{Question}
			Assessing information systems risk is BEST achieved by:
		\end{block}
		\begin{enumerate}[(A)]
			\item Using the enterprise's past actual loss experience
			\item Reviewing published loss statistics from comparable enterprises
			\item Evaluating threats associated with existing information systems assets
			\item Reviewing information systems control weaknesses in audit reports
		\end{enumerate}
	\end{frame}
	\begin{frame}{D2-Q164: Answer}
		\begin{block}{Correct Answer: C}
			To assess IT risk, threats and vulnerabilities need to be evaluated using qualitative or quantitative risk assessment approaches.
		\end{block}
	\end{frame}
	
	\begin{frame}{D2-Q165: Lack of Adequate Controls}
		\begin{block}{Question}
			A lack of adequate controls represents:
		\end{block}
		\begin{enumerate}[(A)]
			\item A vulnerability
			\item An impact
			\item An asset
			\item A threat
		\end{enumerate}
	\end{frame}
	\begin{frame}{D2-Q165: Answer}
		\begin{block}{Correct Answer: A}
			Lack of adequate controls represents a vulnerability, exposing sensitive information and data to the risk of malicious damage, attack, or unauthorized access.
		\end{block}
	\end{frame}
	
	\begin{frame}{D2-Q166: Learning of Security Breach Elsewhere}
		\begin{block}{Question}
			An enterprise learns of a security breach at another entity using similar network technology. The MOST important action for a risk practitioner is to:
		\end{block}
		\begin{enumerate}[(A)]
			\item Assess the likelihood of the incident occurring at the risk practitioner's enterprise
			\item Discontinue the use of the vulnerable technology
			\item Report to senior management that the enterprise is not affected
			\item Remind staff that no similar security breaches have taken place
		\end{enumerate}
	\end{frame}
	\begin{frame}{D2-Q166: Answer}
		\begin{block}{Correct Answer: A}
			The likelihood of a similar incident occurring at the risk practitioner's enterprise should be assessed first, based on available information.
		\end{block}
	\end{frame}
	
	\begin{frame}{D2-Q167: Trigger for Comprehensive Risk Assessment}
		\begin{block}{Question}
			The MOST likely trigger for conducting a comprehensive risk assessment is detection of changes in:
		\end{block}
		\begin{enumerate}[(A)]
			\item The asset inventory
			\item Asset classification levels
			\item The business environment
			\item Information security policies
		\end{enumerate}
	\end{frame}
	\begin{frame}{D2-Q167: Answer}
		\begin{block}{Correct Answer: C}
			Changes in the business environment, including new threats, vulnerabilities, or changes to information asset deployment, will trigger comprehensive periodic risk assessment.
		\end{block}
	\end{frame}
	
	\begin{frame}{D2-Q168: Anonymous Risk Identification}
		\begin{block}{Question}
			Which is the BEST risk identification technique to support an enterprise that allows employees to identify risk anonymously?
		\end{block}
		\begin{enumerate}[(A)]
			\item The Delphi technique
			\item Isolated pilot groups
			\item A SWOT analysis
			\item A root cause analysis
		\end{enumerate}
	\end{frame}
	\begin{frame}{D2-Q168: Answer}
		\begin{block}{Correct Answer: A}
			With the Delphi technique, polling or information gathering is done either anonymously or privately between the interviewer and interviewee.
		\end{block}
	\end{frame}
	
	\begin{frame}{D2-Q169: Risk Register Content}
		\begin{block}{Question}
			Which BEST describes the information needed for each risk in a risk register?
		\end{block}
		\begin{enumerate}[(A)]
			\item Risk scenario including date, description, impact, probability, risk score, mitigation action and owner
			\item Risk scenario including date, description, risk score, cost to remediate, communication plan and owner
			\item Risk scenario including date, description, impact, cost to remediate and owner
			\item Various activities leading to risk management planning
		\end{enumerate}
	\end{frame}
	\begin{frame}{D2-Q169: Answer}
		\begin{block}{Correct Answer: A}
			Information required for each risk in a risk register includes date, description, impact, probability, risk score, mitigation action and owner.
		\end{block}
	\end{frame}
	
	\begin{frame}{D2-Q170: BIA Greatest Advantage}
		\begin{block}{Question}
			The GREATEST advantage in performing a business impact analysis is that it:
		\end{block}
		\begin{enumerate}[(A)]
			\item Does not have to be updated
			\item Promotes continuity awareness in the enterprise
			\item Requires only qualitative estimates
			\item Eliminates the need for risk analysis
		\end{enumerate}
	\end{frame}
	\begin{frame}{D2-Q170: Answer}
		\begin{block}{Correct Answer: B}
			A BIA raises enterprise-wide awareness of risk to business recovery and continuity.
		\end{block}
	\end{frame}
	
	\begin{frame}{D2-Q171: Primary Advantage of Risk Register}
		\begin{block}{Question}
			The PRIMARY advantage of creating and updating a risk register is to:
		\end{block}
		\begin{enumerate}[(A)]
			\item Ensure that an inventory of identified risk is maintained
			\item Record all risk scenarios considered during risk identification
			\item Collect similar data on all risk identified within the enterprise
			\item Run reports based on various risk scenarios
		\end{enumerate}
	\end{frame}
	\begin{frame}{D2-Q171: Answer}
		\begin{block}{Correct Answer: A}
			Once assets and risk are identified, the risk register is used as an inventory of that risk. The risk register can accelerate risk decision-making and establish accountability for specific risk.
		\end{block}
	\end{frame}
	
	\begin{frame}{D2-Q172: Measuring Risk Management Development}
		\begin{block}{Question}
			Which BEST assists a risk practitioner in measuring the existing level of development of risk management processes against the desired state?
		\end{block}
		\begin{enumerate}[(A)]
			\item A capability maturity model
			\item Risk management audit reports
			\item A balanced scorecard
			\item Enterprise security architecture
		\end{enumerate}
	\end{frame}
	\begin{frame}{D2-Q172: Answer}
		\begin{block}{Correct Answer: A}
			The capability maturity model grades processes on a scale of 0 to 5, based on their maturity.
		\end{block}
	\end{frame}
	
	\begin{frame}{D2-Q173: Most Effective Risk Assessment}
		\begin{block}{Question}
			Which is MOST effective in assessing business risk?
		\end{block}
		\begin{enumerate}[(A)]
			\item A use case analysis
			\item A business case analysis
			\item Risk scenarios
			\item A risk plan
		\end{enumerate}
	\end{frame}
	\begin{frame}{D2-Q173: Answer}
		\begin{block}{Correct Answer: C}
			Risk scenarios are the most effective technique in assessing business risk. Scenarios help determine the likelihood and impact of an identified risk.
		\end{block}
	\end{frame}
	
	\begin{frame}{D2-Q174: Preparation of Risk Register}
		\begin{block}{Question}
			The preparation of a risk register begins in which risk management process?
		\end{block}
		\begin{enumerate}[(A)]
			\item Risk response planning
			\item Risk monitoring and control
			\item Risk management planning
			\item Risk identification
		\end{enumerate}
	\end{frame}
	\begin{frame}{D2-Q174: Answer}
		\begin{block}{Correct Answer: D}
			The risk register details all identified risk, including description, category, cause, probability of occurring, impact on objectives, proposed responses, owners and current status. The primary outputs of risk identification are the initial entries into the risk register.
		\end{block}
	\end{frame}
	
	\begin{frame}{D2-Q175: BIA Primary Use}
		\begin{block}{Question}
			A business impact analysis is PRIMARILY used to:
		\end{block}
		\begin{enumerate}[(A)]
			\item Estimate the resources required to resume normal operations
			\item Evaluate the impact of disruption on ability to operate over time
			\item Calculate the likelihood and impact of known threats
			\item Evaluate high-level business requirements
		\end{enumerate}
	\end{frame}
	\begin{frame}{D2-Q175: Answer}
		\begin{block}{Correct Answer: B}
			A BIA is primarily intended to evaluate the impact of disruption over time to an enterprise's ability to operate. It determines the urgency of each business activity.
		\end{block}
	\end{frame}
	
	\begin{frame}{D2-Q176: Risk Scenarios Enable Assessment}
		\begin{block}{Question}
			Risk scenarios enable the risk assessment process because they:
		\end{block}
		\begin{enumerate}[(A)]
			\item Cover a wide range of potential risk
			\item Minimize the need for quantitative risk analysis techniques
			\item Segregate IT risk from business risk
			\item Help estimate the frequency and impact of risk
		\end{enumerate}
	\end{frame}
	\begin{frame}{D2-Q176: Answer}
		\begin{block}{Correct Answer: D}
			Risk scenarios help to estimate the frequency and impact of risk—two key elements of the risk assessment process.
		\end{block}
	\end{frame}
	
	\begin{frame}{D2-Q177: Risk Register for Scenarios}
		\begin{block}{Question}
			Which information in the risk register BEST helps in developing proper risk scenarios? A list of:
		\end{block}
		\begin{enumerate}[(A)]
			\item Potential threats to assets
			\item Residual risk on individual assets
			\item Accepted risk
			\item Security incidents
		\end{enumerate}
	\end{frame}
	\begin{frame}{D2-Q177: Answer}
		\begin{block}{Correct Answer: A}
			Identifying potential threats to business assets will help isolate vulnerabilities and associated risk, all of which contribute to developing proper risk scenarios.
		\end{block}
	\end{frame}
	
	\begin{frame}{D2-Q178: Identifying IS Control Deficiencies}
		\begin{block}{Question}
			Which BEST helps identify information systems control deficiencies?
		\end{block}
		\begin{enumerate}[(A)]
			\item Gap analysis
			\item The current IT risk profile
			\item The IT controls framework
			\item Countermeasure analysis
		\end{enumerate}
	\end{frame}
	\begin{frame}{D2-Q178: Answer}
		\begin{block}{Correct Answer: A}
			The gap between desired control objectives and actual control design and operational effectiveness identifies control deficiencies in information systems.
		\end{block}
	\end{frame}
	
	\begin{frame}{D2-Q179: Reliable Assessment of Server Performance}
		\begin{block}{Question}
			When assessing the performance of a critical application server, the MOST reliable assessment results may be obtained from:
		\end{block}
		\begin{enumerate}[(A)]
			\item Activation of native database auditing
			\item Documentation of performance objectives
			\item Continuous monitoring
			\item Documentation of security modules
		\end{enumerate}
	\end{frame}
	\begin{frame}{D2-Q179: Answer}
		\begin{block}{Correct Answer: C}
			With continuous monitoring it is possible to track key performance metrics and also possible to address critical issues with minimum impact.
		\end{block}
	\end{frame}
	
	\begin{frame}{D2-Q180: Updating Risk Register Greatest Risk}
		\begin{block}{Question}
			Which presents the GREATEST risk when updating the risk register? Updates are:
		\end{block}
		\begin{enumerate}[(A)]
			\item Carried out jointly with other functions
			\item Carried out following incidents
			\item Carried out annually
			\item Subject to approval by the chief information security officer
		\end{enumerate}
	\end{frame}
	\begin{frame}{D2-Q180: Answer}
		\begin{block}{Correct Answer: C}
			Updating the risk register only annually means the risk register does not reflect the true status of IT risk in the enterprise.
		\end{block}
	\end{frame}
	
	\begin{frame}{D2-Q181: IT Risk Measurement}
		\begin{block}{Question}
			IT risk is measured by its:
		\end{block}
		\begin{enumerate}[(A)]
			\item Level of damage to IT systems
			\item Impact on business operations
			\item Cost of countermeasures
			\item Annual loss expectancy
		\end{enumerate}
	\end{frame}
	\begin{frame}{D2-Q181: Answer}
		\begin{block}{Correct Answer: B}
			IT risk includes information and communication technology risk but is primarily measured by its impact on the business.
		\end{block}
	\end{frame}
	
	\begin{frame}{D2-Q182: Statistical Methods Analysis}
		\begin{block}{Question}
			Deriving the likelihood and impact of risk scenarios through statistical methods is BEST described as:
		\end{block}
		\begin{enumerate}[(A)]
			\item Quantitative risk analysis
			\item Risk scenario analysis
			\item Qualitative risk analysis
			\item Probabilistic risk assessment
		\end{enumerate}
	\end{frame}
	\begin{frame}{D2-Q182: Answer}
		\begin{block}{Correct Answer: A}
			Quantitative risk analysis derives the probability and impact of risk scenarios from statistical methods and data.
		\end{block}
	\end{frame}
	
	\begin{frame}{D2-Q183: Local Management Mitigation}
		\begin{block}{Question}
			During an internal risk assessment in a global enterprise, a risk manager notes that local management has proactively mitigated some high-level risk related to the global purchasing process. This means that:
		\end{block}
		\begin{enumerate}[(A)]
			\item Local management is now responsible for the risk
			\item The risk owner is the corporate chief risk officer
			\item The risk owner is the local purchasing manager
			\item Corporate management remains responsible for the risk
		\end{enumerate}
	\end{frame}
	\begin{frame}{D2-Q183: Answer}
		\begin{block}{Correct Answer: D}
			Corporate management remains responsible for the risk, even when the risk response is executed at a lower organizational level.
		\end{block}
	\end{frame}
	
	\begin{frame}{D2-Q184: Estimating Likelihood of Significant Events}
		\begin{block}{Question}
			Which BEST estimates the likelihood of significant events affecting an enterprise?
		\end{block}
		\begin{enumerate}[(A)]
			\item Threat analysis
			\item Cost-benefit analysis
			\item Scenario analysis
			\item Countermeasure analysis
		\end{enumerate}
	\end{frame}
	\begin{frame}{D2-Q184: Answer}
		\begin{block}{Correct Answer: C}
			Scenario analysis, along with vulnerability analysis, best determines whether a particular risk is relevant to the enterprise and helps estimate the likelihood that significant events will affect the enterprise.
		\end{block}
	\end{frame}
	
	\begin{frame}{D2-Q185: Improving Decision-Making}
		\begin{block}{Question}
			Which BEST improves decision-making related to risk?
		\end{block}
		\begin{enumerate}[(A)]
			\item Maintaining a documented risk register of all possible risk
			\item Risk awareness training in line with the risk culture
			\item Maintaining updated security policies and procedures
			\item Allocating accountability of risk to the department as a whole
		\end{enumerate}
	\end{frame}
	\begin{frame}{D2-Q185: Answer}
		\begin{block}{Correct Answer: A}
			Maintaining a documented risk register improves decision-making related to risk response because a risk register captures the population of relevant risk scenarios and provides a basis for prioritization.
		\end{block}
	\end{frame}
	
	\begin{frame}{D2-Q186: Independent Risk Management Review}
		\begin{block}{Question}
			The PRIMARY reason to have the risk management process reviewed by an independent risk management professional is to:
		\end{block}
		\begin{enumerate}[(A)]
			\item Validate cost-effective solutions for mitigating risk
			\item Validate control weaknesses detected by the internal team
			\item Assess the validity of the end-to-end process
			\item Assess whether the risk profile and risk factors are properly defined
		\end{enumerate}
	\end{frame}
	\begin{frame}{D2-Q186: Answer}
		\begin{block}{Correct Answer: C}
			The independent risk professional will be unbiased to review the risk management process end to end.
		\end{block}
	\end{frame}
	
	\begin{frame}{D2-Q187: Review of Risk Analysis Results}
		\begin{block}{Question}
			Which is BEST suited for the review of IT risk analysis results before the results are sent to management?
		\end{block}
		\begin{enumerate}[(A)]
			\item An internal audit review
			\item A peer review
			\item A compliance review
			\item A risk policy review
		\end{enumerate}
	\end{frame}
	\begin{frame}{D2-Q187: Answer}
		\begin{block}{Correct Answer: B}
			It is effective, efficient and good practice to perform a peer review of IT risk analysis results before sending them to management.
		\end{block}
	\end{frame}
	
	\begin{frame}{D2-Q188: First Step Identifying IT Risk}
		\begin{block}{Question}
			The FIRST step in identifying and assessing IT risk is to:
		\end{block}
		\begin{enumerate}[(A)]
			\item Confirm the risk tolerance level of the enterprise
			\item Identify threats and vulnerabilities
			\item Gather information on the current and future environment
			\item Review past incident reports and response activity
		\end{enumerate}
	\end{frame}
	\begin{frame}{D2-Q188: Answer}
		\begin{block}{Correct Answer: C}
			The first step in any risk assessment is to gather information about the current state and pending internal and external changes to the enterprise's environment.
		\end{block}
	\end{frame}
	
	\begin{frame}{D2-Q189: Basis for Risk Scenarios}
		\begin{block}{Question}
			Risk scenarios should be created PRIMARILY based on which of the following?
		\end{block}
		\begin{enumerate}[(A)]
			\item Input from senior management
			\item Previous security incidents
			\item Threats that the enterprise faces
			\item Results of the risk analysis
		\end{enumerate}
	\end{frame}
	\begin{frame}{D2-Q189: Answer}
		\begin{block}{Correct Answer: C}
			When creating risk scenarios, the most important factor to consider is the likelihood of threats or threat actions occurring due to the risk.
		\end{block}
	\end{frame}
	
	\begin{frame}{D2-Q190: Trigger for Ad Hoc Risk Assessment}
		\begin{block}{Question}
			Which of the following triggers performance of an internal ad hoc risk assessment before the annual occurrence?
		\end{block}
		\begin{enumerate}[(A)]
			\item A new chief information officer is hired
			\item Senior management adjusts risk appetite
			\item Risk changes on a frequent basis
			\item A new system is introduced into the environment
		\end{enumerate}
	\end{frame}
	\begin{frame}{D2-Q190: Answer}
		\begin{block}{Correct Answer: D}
			Introduction of new systems adds to overall risk of business objectives. The level of new or added risk should be determined via an ad hoc risk assessment.
		\end{block}
	\end{frame}
	
	\begin{frame}{D2-Q191: Emerging Threat}
		\begin{block}{Question}
			When a start-up company becomes popular, it suddenly is the target of hackers. This is considered:
		\end{block}
		\begin{enumerate}[(A)]
			\item An emerging vulnerability
			\item A vulnerability event
			\item An emerging threat
			\item An environmental risk factor
		\end{enumerate}
	\end{frame}
	\begin{frame}{D2-Q191: Answer}
		\begin{block}{Correct Answer: C}
			A threat is any event in which an asset is exposed to a threat condition or action that has the potential to directly result in harm.
		\end{block}
	\end{frame}
	
	\begin{frame}{D2-Q192: External Team Reviews Documentation}
		\begin{block}{Question}
			The PRIMARY reason an external risk assessment team reviews documentation as the first step in the risk assessment is to gain a thorough understanding of:
		\end{block}
		\begin{enumerate}[(A)]
			\item The technologies used
			\item Gaps in the documentation
			\item The enterprise's business processes
			\item The risk assessment plan
		\end{enumerate}
	\end{frame}
	\begin{frame}{D2-Q192: Answer}
		\begin{block}{Correct Answer: C}
			In order to evaluate risk, the external assessment team should thoroughly understand the enterprise's business processes before the assessment.
		\end{block}
	\end{frame}
	
	\begin{frame}{D2-Q193: Maturity Model Dimension}
		\begin{block}{Question}
			Which capability dimension is MOST important when using a maturity model for assessing the risk management process?
		\end{block}
		\begin{enumerate}[(A)]
			\item Effectiveness
			\item Efficiency
			\item Profitability
			\item Performance
		\end{enumerate}
	\end{frame}
	\begin{frame}{D2-Q193: Answer}
		\begin{block}{Correct Answer: D}
			Performance is achieved when the implemented process fulfills its purpose; thus, performance is the most important capability dimension.
		\end{block}
	\end{frame}
	
	\begin{frame}{D2-Q194: Lapsed Insurance Premium}
		\begin{block}{Question}
			A small start-up software development company has been flooded and insurance does not pay because the premium has lapsed. In relation to risk management, the lapsed premium is considered a:
		\end{block}
		\begin{enumerate}[(A)]
			\item Risk
			\item Vulnerability
			\item Threat
			\item Negligence
		\end{enumerate}
	\end{frame}
	\begin{frame}{D2-Q194: Answer}
		\begin{block}{Correct Answer: B}
			A lapsed insurance premium describes a vulnerability. A vulnerability is a weakness in the design, implementation, operation or internal control of a process.
		\end{block}
	\end{frame}
	
	\begin{frame}{D2-Q195: Value of Risk Register}
		\begin{block}{Question}
			Which statement BEST describes the value of a risk register?
		\end{block}
		\begin{enumerate}[(A)]
			\item It captures the risk inventory
			\item It drives the risk response plan
			\item It is a risk reporting tool
			\item It lists internal risk and external risk
		\end{enumerate}
	\end{frame}
	\begin{frame}{D2-Q195: Answer}
		\begin{block}{Correct Answer: B}
			Risk registers serve as the main reference for all risk-related information, supporting risk-related decisions such as risk response activities and their prioritization.
		\end{block}
	\end{frame}
	
	\begin{frame}{D2-Q196: Main Objective of Risk Identification}
		\begin{block}{Question}
			What is the MAIN objective of risk identification?
		\end{block}
		\begin{enumerate}[(A)]
			\item To detect possible threats that may affect the business
			\item To ensure that risk factors and root causes are managed
			\item To enable the review of the key performance indicators
			\item To provide qualitative impact values to stakeholders
		\end{enumerate}
	\end{frame}
	\begin{frame}{D2-Q196: Answer}
		\begin{block}{Correct Answer: A}
			Risk identification is the process of determining and documenting the risk that an enterprise faces. The identification of risk is based on the recognition of threats, vulnerabilities, assets and controls in the enterprise's operational environment.
		\end{block}
	\end{frame}
	
	\begin{frame}{D2-Q197: Maturity Model Benefit for Data Management}
		\begin{block}{Question}
			The PRIMARY benefit of using a maturity model to assess the enterprise's data management process is that it:
		\end{block}
		\begin{enumerate}[(A)]
			\item Can be used for benchmarking
			\item Helps identify gaps
			\item Provides goals and objectives
			\item Enforces continuous improvement
		\end{enumerate}
	\end{frame}
	\begin{frame}{D2-Q197: Answer}
		\begin{block}{Correct Answer: B}
			Maturity models can be used to help identify gaps between the current and the desired state to help enterprises determine necessary remediation efforts.
		\end{block}
	\end{frame}
	
	\begin{frame}{D2-Q198: Goal of IT Risk Analysis}
		\begin{block}{Question}
			The goal of IT risk analysis is to:
		\end{block}
		\begin{enumerate}[(A)]
			\item Enable the alignment of IT risk management with enterprise risk management
			\item Enable the prioritization of risk responses
			\item Satisfy legal and regulatory compliance requirements
			\item Identify known threats and vulnerabilities to information assets
		\end{enumerate}
	\end{frame}
	\begin{frame}{D2-Q198: Answer}
		\begin{block}{Correct Answer: B}
			Risk analysis estimates the likelihood and magnitude of IT risk scenarios. Risk analysis helps ensure that areas with the greatest risk likelihood and impact are prioritized above those with lower likelihood and impact.
		\end{block}
	\end{frame}
	
	\begin{frame}{D2-Q199: Factor Assessed After Likelihood}
		\begin{block}{Question}
			Which of the following factors should be assessed after the likelihood of a loss event has been determined?
		\end{block}
		\begin{enumerate}[(A)]
			\item Risk tolerance
			\item Magnitude of impact
			\item Residual risk
			\item Compensating controls
		\end{enumerate}
	\end{frame}
	\begin{frame}{D2-Q199: Answer}
		\begin{block}{Correct Answer: B}
			Once likelihood has been determined, the next step is to determine magnitude of impact.
		\end{block}
	\end{frame}
	
	\begin{frame}{D2-Q200: IT Risk Register Greatest Benefit}
		\begin{block}{Question}
			The GREATEST benefit of using an IT risk register is that it is:
		\end{block}
		\begin{enumerate}[(A)]
			\item A list of potential events that have been identified to understand their impact
			\item A document used to track risk that has been identified, analyzed and prioritized
			\item A list of risk that can be shared with all stakeholders in an easy-to-read format
			\item The basis for choosing a commercial, off-the-shelf risk management tool
		\end{enumerate}
	\end{frame}
	\begin{frame}{D2-Q200: Answer}
		\begin{block}{Correct Answer: B}
			The greatest benefit of a risk register is that it provides information about the likelihood, impact and prioritization of all identified IT risk.
		\end{block}
	\end{frame}
	
	% ============================================================
	% DOMAIN 3 --- RISK RESPONSE AND REPORTING (32%)
	% ============================================================
	\section{Domain 3: Risk Response and Reporting}
	
	\begin{frame}{D3-Q201: Outsourced Noncore Process Concern}
		\begin{block}{Question}
			Which outcome of outsourcing noncore processes is of GREATEST concern to management?
		\end{block}
		\begin{enumerate}[(A)]
			\item Total cost of ownership exceeds projections
			\item Internal information systems experience is lost
			\item Employees of the vendor are disloyal to the client
			\item Processing of sensitive data is subcontracted by the vendor
		\end{enumerate}
	\end{frame}
	\begin{frame}{D3-Q201: Answer}
		\begin{block}{Correct Answer: D}
			The greatest risk in third-party relationships is ceding direct control of IS processes. Subcontracting increases this risk.
		\end{block}
	\end{frame}
	
	\begin{frame}{D3-Q202: Proposing Risk Treatment}
		\begin{block}{Question}
			When proposing the implementation of a specific risk treatment, a risk practitioner PRIMARILY uses a:
		\end{block}
		\begin{enumerate}[(A)]
			\item Technical evaluation report
			\item Business case
			\item Vulnerability assessment report
			\item Budgetary requirements
		\end{enumerate}
	\end{frame}
	\begin{frame}{D3-Q202: Answer}
		\begin{block}{Correct Answer: B --- Business case}
			A manager should base any proposed risk response on a risk assessment in the context of business objectives. The manager must document costs of controls and compare them against benefits.
		\end{block}
	\end{frame}
	
	\begin{frame}{D3-Q203: Most Important for Effective Risk Management}
		\begin{block}{Question}
			Which is MOST important for effective risk management?
		\end{block}
		\begin{enumerate}[(A)]
			\item Assignment of risk owners to identified risk
			\item Ensuring compliance with regulatory requirements
			\item Integration of risk management into operational processes
			\item Implementation of a risk avoidance strategy
		\end{enumerate}
	\end{frame}
	\begin{frame}{D3-Q203: Answer}
		\begin{block}{Correct Answer: A}
			It is of utmost importance to assign risk to individual owners and therein maximize accountability.
		\end{block}
	\end{frame}
	
	\begin{frame}{D3-Q204: Most Important Outsourcing Question}
		\begin{block}{Question}
			An enterprise outsourced several business functions including IT development, data hosting, and support. What is the MOST important question regarding the outsourcing arrangements?
		\end{block}
		\begin{enumerate}[(A)]
			\item Are policies and procedures in place to handle security exceptions?
			\item Is the supplier meeting the terms of the SLAs?
			\item Is the provider's security program based on an international standard?
			\item Are specific security controls mandated in the outsourcing contract?
		\end{enumerate}
	\end{frame}
	\begin{frame}{D3-Q204: Answer}
		\begin{block}{Correct Answer: D}
			Without enumerating security requirements directly in the outsourcing contract, the outsourcing company has no assurance that the provider will comply with specific security requirements.
		\end{block}
	\end{frame}
	
	\begin{frame}{D3-Q205: Basis for Establishing Risk Ownership}
		\begin{block}{Question}
			Which activity provides the BEST basis for establishing risk ownership?
		\end{block}
		\begin{enumerate}[(A)]
			\item Documenting interdependencies between departments
			\item Mapping identified risk to a specific business process
			\item Referring to available RACI charts for business processes
			\item Distributing risk among identified business process owners
		\end{enumerate}
	\end{frame}
	\begin{frame}{D3-Q205: Answer}
		\begin{block}{Correct Answer: B}
			Mapping identified risk to a specific business process helps identify the process owner. Aggregation of related business processes results in identification of the prospective risk owner.
		\end{block}
	\end{frame}
	
	\begin{frame}{D3-Q206: Risk Response Report Recommendations}
		\begin{block}{Question}
			A risk response report includes recommendations for:
		\end{block}
		\begin{enumerate}[(A)]
			\item Acceptance
			\item Assessment
			\item Evaluation
			\item Quantification
		\end{enumerate}
	\end{frame}
	\begin{frame}{D3-Q206: Answer}
		\begin{block}{Correct Answer: A --- Acceptance}
			Acceptance of a risk is an alternative to be considered in the risk response process.
		\end{block}
	\end{frame}
	
	\begin{frame}{D3-Q207: Reduced to Achieve Acceptable Risk}
		\begin{block}{Question}
			Which is MOST likely to be reduced to achieve acceptable risk?
		\end{block}
		\begin{enumerate}[(A)]
			\item Risk appetite
			\item Control risk
			\item Residual risk
			\item Inherent risk
		\end{enumerate}
	\end{frame}
	\begin{frame}{D3-Q207: Answer}
		\begin{block}{Correct Answer: C --- Residual risk}
			Residual risk is the remaining risk after management has implemented a risk response. Acceptable risk is achieved when residual risk is reduced to levels within the enterprise's risk appetite.
		\end{block}
	\end{frame}
	
	\begin{frame}{D3-Q208: Denial-of-Service Vulnerability Decision}
		\begin{block}{Question}
			A global financial institution decided not to take further action on a denial-of-service vulnerability. The MOST likely reason is:
		\end{block}
		\begin{enumerate}[(A)]
			\item The needed countermeasure is too complicated to deploy
			\item There are sufficient safeguards to prevent this risk
			\item The likelihood of the risk occurring is unknown
			\item The cost of countermeasure outweighs the value of the asset and potential loss
		\end{enumerate}
	\end{frame}
	\begin{frame}{D3-Q208: Answer}
		\begin{block}{Correct Answer: D}
			An enterprise may decide to accept a specific risk because the protection would cost more than the potential loss.
		\end{block}
	\end{frame}
	
	\begin{frame}{D3-Q209: Two-Factor Authentication Cost-Benefit}
		\begin{block}{Question}
			Which is MOST relevant to include in a cost-benefit analysis of a two-factor authentication system?
		\end{block}
		\begin{enumerate}[(A)]
			\item The approved budget of the project
			\item The frequency of incidents
			\item The annual loss expectancy of security incidents
			\item The total cost of ownership
		\end{enumerate}
	\end{frame}
	\begin{frame}{D3-Q209: Answer}
		\begin{block}{Correct Answer: D}
			Total cost of ownership is the most relevant piece of information because it establishes a cost baseline that must be considered for the full life cycle of the control.
		\end{block}
	\end{frame}
	
	\begin{frame}{D3-Q210: Most Important Outsourcing Contract Part}
		\begin{block}{Question}
			Which is the MOST important part of any outsourcing contract?
		\end{block}
		\begin{enumerate}[(A)]
			\item The right to audit the outsourcing provider
			\item Provisions to assess the compliance of the provider
			\item Procedures for dealing with incident notification
			\item Requirements to encrypt hosted data
		\end{enumerate}
	\end{frame}
	\begin{frame}{D3-Q210: Answer}
		\begin{block}{Correct Answer: B}
			If a contract contains no provision to monitor and hold a supplier accountable for security, then the outsourcing enterprise cannot ensure compliance.
		\end{block}
	\end{frame}
	
	\begin{frame}{D3-Q211: Greatest Risk with Cloud Provider}
		\begin{block}{Question}
			Which poses the GREATEST risk to an enterprise that recently engaged a cloud provider?
		\end{block}
		\begin{enumerate}[(A)]
			\item The cloud provider's primary facility is in the same vicinity
			\item The service level agreement is ambiguous
			\item References from other customers were not obtained
			\item Auditing the vendor requires dependence on a third-party audit firm
		\end{enumerate}
	\end{frame}
	\begin{frame}{D3-Q211: Answer}
		\begin{block}{Correct Answer: B}
			If the service level agreement is ambiguous, it will be difficult to determine whether the provider complies.
		\end{block}
	\end{frame}
	
	\begin{frame}{D3-Q212: Insuring Against Unmitigable Risk}
		\begin{block}{Question}
			When a risk cannot be sufficiently mitigated through manual or automatic controls, which option will BEST protect the enterprise from the potential financial impact?
		\end{block}
		\begin{enumerate}[(A)]
			\item Insuring against the risk
			\item Updating the IT risk registry
			\item Improving staff training in the risk area
			\item Outsourcing the related process to a third party
		\end{enumerate}
	\end{frame}
	\begin{frame}{D3-Q212: Answer}
		\begin{block}{Correct Answer: A}
			An insurance policy can compensate the enterprise monetarily for the impact of the risk by transferring the risk to the insurance company.
		\end{block}
	\end{frame}
	
	\begin{frame}{D3-Q213: Cost of Mitigation Exceeds Benefit}
		\begin{block}{Question}
			After a risk assessment, it is determined that the cost to mitigate the risk is much greater than the benefit to be derived. The BEST risk response is that the risk be:
		\end{block}
		\begin{enumerate}[(A)]
			\item Treated
			\item Terminated
			\item Accepted
			\item Transferred
		\end{enumerate}
	\end{frame}
	\begin{frame}{D3-Q213: Answer}
		\begin{block}{Correct Answer: C --- Accepted}
			When the cost of control is more than the cost of the potential impact, the risk should be accepted.
		\end{block}
	\end{frame}
	
	\begin{frame}{D3-Q214: Improving Continuity of Operations}
		\begin{block}{Question}
			A data center provider's clients complain of frequent service disruptions. Management sets up a project team to improve continuity. Which risk management activity will MOST benefit this initiative?
		\end{block}
		\begin{enumerate}[(A)]
			\item Develop key risk indicators
			\item Determine risk and control ownership
			\item Update the risk register
			\item Develop IT risk scenarios
		\end{enumerate}
	\end{frame}
	\begin{frame}{D3-Q214: Answer}
		\begin{block}{Correct Answer: D}
			The relationship between IT risk scenarios and business impact must be established to understand the effects of possible adverse events on enterprise objectives.
		\end{block}
	\end{frame}
	
	\begin{frame}{D3-Q215: Reciprocal Agreement for Equipment}
		\begin{block}{Question}
			Critical IT equipment will be delivered late due to flooding, but a reciprocal agreement exists with another company for a replacement until it arrives. This is an example of risk:
		\end{block}
		\begin{enumerate}[(A)]
			\item Transfer
			\item Avoidance
			\item Acceptance
			\item Mitigation
		\end{enumerate}
	\end{frame}
	\begin{frame}{D3-Q215: Answer}
		\begin{block}{Correct Answer: D --- Mitigation}
			Risk mitigation attempts to reduce the impact when a risk event occurs. Making plans such as a reciprocal arrangement reduces the consequence of the risk event.
		\end{block}
	\end{frame}
	
	\begin{frame}{D3-Q216: Selecting an Appropriate Risk Response}
		\begin{block}{Question}
			Which would BEST help an enterprise select an appropriate risk response?
		\end{block}
		\begin{enumerate}[(A)]
			\item An analysis of change in the risk environment
			\item An analysis of risk that can be transferred
			\item An analysis of likelihood and impact of risk scenarios
			\item An analysis of control costs and benefits
		\end{enumerate}
	\end{frame}
	\begin{frame}{D3-Q216: Answer}
		\begin{block}{Correct Answer: D}
			An analysis of costs and benefits for controls helps an enterprise understand if it can mitigate the risk to an acceptable level.
		\end{block}
	\end{frame}
	
	\begin{frame}{D3-Q217: Optimal Return on Security Investment}
		\begin{block}{Question}
			Which leads to the BEST optimal return on security investment?
		\end{block}
		\begin{enumerate}[(A)]
			\item Deploying maximum protection across all information assets
			\item Focusing on the most important information assets and then determining their protection
			\item Deploying minimum protection across all information assets
			\item Investing only after a major security incident is reported
		\end{enumerate}
	\end{frame}
	\begin{frame}{D3-Q217: Answer}
		\begin{block}{Correct Answer: B}
			To optimize return on security investment, focus on identifying important information assets and protect them appropriately.
		\end{block}
	\end{frame}
	
	\begin{frame}{D3-Q218: Compromised Administrator Account Impact}
		\begin{block}{Question}
			Which control can reduce the potential impact of a malicious hacker who gains access to an administrator account?
		\end{block}
		\begin{enumerate}[(A)]
			\item Multifactor authentication
			\item Audit logging
			\item Least privilege
			\item Password policy
		\end{enumerate}
	\end{frame}
	\begin{frame}{D3-Q218: Answer}
		\begin{block}{Correct Answer: C --- Least privilege}
			Least privilege can reduce the impact of a compromised account within the scope of its intended use and limit impact to the enterprise as a whole.
		\end{block}
	\end{frame}
	
	\begin{frame}{D3-Q219: Critical Element of Effective Risk Program}
		\begin{block}{Question}
			Which is critical to the risk practitioner for an effective risk management program?
		\end{block}
		\begin{enumerate}[(A)]
			\item Risk response strategy
			\item Risk register content
			\item Risk profile content
			\item Risk owners and accountability
		\end{enumerate}
	\end{frame}
	\begin{frame}{D3-Q219: Answer}
		\begin{block}{Correct Answer: D}
			The identification of risk owners is critical because risk owners must make informed and cost-effective decisions regarding appropriate controls to mitigate their owned risk.
		\end{block}
	\end{frame}
	
	\begin{frame}{D3-Q220: Anti-Malware Cost Exceeds Loss}
		\begin{block}{Question}
			Where the cost of anti-malware exceeds the loss expectancy of malware threats, what is the MOST viable risk response?
		\end{block}
		\begin{enumerate}[(A)]
			\item Risk elimination
			\item Risk acceptance
			\item Risk transfer
			\item Risk mitigation
		\end{enumerate}
	\end{frame}
	\begin{frame}{D3-Q220: Answer}
		\begin{block}{Correct Answer: B --- Risk acceptance}
			When the cost of a risk response exceeds the loss expectancy, the most viable risk response is risk acceptance.
		\end{block}
	\end{frame}
	
	\begin{frame}{D3-Q221: Best Example of Risk Avoidance}
		\begin{block}{Question}
			Which is the BEST example of risk avoidance behavior?
		\end{block}
		\begin{enumerate}[(A)]
			\item Taking no action against the risk
			\item Outsourcing the related process
			\item Insuring against a specific event
			\item Exiting the process that gives rise to risk
		\end{enumerate}
	\end{frame}
	\begin{frame}{D3-Q221: Answer}
		\begin{block}{Correct Answer: D}
			Avoidance means exiting the activities or conditions that give rise to risk.
		\end{block}
	\end{frame}
	
	\begin{frame}{D3-Q222: Addressing Critical Dependencies}
		\begin{block}{Question}
			Which would ensure that critical dependencies are addressed in the risk treatment plan?
		\end{block}
		\begin{enumerate}[(A)]
			\item Implement the risk treatment strategy for all possible risk
			\item Verify the accomplishment of business objectives through a top-down process review
			\item Treat each risk independently
			\item Verify the accomplishment of business objectives through a bottom-up process review
		\end{enumerate}
	\end{frame}
	\begin{frame}{D3-Q222: Answer}
		\begin{block}{Correct Answer: B}
			Development of an overall risk treatment strategy should be a top-down process, driven jointly by the need to achieve business objectives and to apply economically feasible constraints.
		\end{block}
	\end{frame}
	
	\begin{frame}{D3-Q223: CISO Recommending Controls}
		\begin{block}{Question}
			A CISO recommends several controls, such as anti-malware, to protect the enterprise's information systems. Which approach to handling risk is being recommended?
		\end{block}
		\begin{enumerate}[(A)]
			\item Risk transference
			\item Risk mitigation
			\item Risk acceptance
			\item Risk avoidance
		\end{enumerate}
	\end{frame}
	\begin{frame}{D3-Q223: Answer}
		\begin{block}{Correct Answer: B --- Risk mitigation}
			By implementing controls, the company is trying to reduce the risk to an acceptable level, thereby mitigating risk.
		\end{block}
	\end{frame}
	
	\begin{frame}{D3-Q224: Best Practice for Outsourcing Risk}
		\begin{block}{Question}
			Which practice BEST mitigates the risk associated with outsourcing a business function?
		\end{block}
		\begin{enumerate}[(A)]
			\item Performing audits to verify compliance with contract requirements
			\item Requiring all vendor staff to complete annual awareness training
			\item Retaining copies of all sensitive data on internal systems
			\item Reviewing the financial records of the vendor
		\end{enumerate}
	\end{frame}
	\begin{frame}{D3-Q224: Answer}
		\begin{block}{Correct Answer: A}
			When an outsourcing relationship is established, the risk of noncompliance with the agreement must be met through review, monitoring, and enforcement of contract terms.
		\end{block}
	\end{frame}
	
	\begin{frame}{D3-Q225: Best Addressed by Transferring Risk}
		\begin{block}{Question}
			Which is BEST addressed by transferring risk?
		\end{block}
		\begin{enumerate}[(A)]
			\item An antiquated fire suppression system in the computer room
			\item The threat of disgruntled employee sabotage
			\item The threat of disgruntled employee theft
			\item A building located in a 100-year flood plain
		\end{enumerate}
	\end{frame}
	\begin{frame}{D3-Q225: Answer}
		\begin{block}{Correct Answer: D}
			Purchasing an insurance policy transfers the risk of a flood. Risk transfer is the process of assigning risk to another entity.
		\end{block}
	\end{frame}
	
	\begin{frame}{D3-Q226: Finalizing Risk Treatment Plan}
		\begin{block}{Question}
			Which would BEST help finalize the risk treatment plan?
		\end{block}
		\begin{enumerate}[(A)]
			\item Vulnerability analysis
			\item Impact analysis
			\item Cost-benefit analysis
			\item SWOT analysis
		\end{enumerate}
	\end{frame}
	\begin{frame}{D3-Q226: Answer}
		\begin{block}{Correct Answer: C}
			A cost-benefit analysis helps determine if the benefit of a control outweighs the cost of implementing the control.
		\end{block}
	\end{frame}
	
	\begin{frame}{D3-Q227: Greatest Benefit of Risk Treatment Plan}
		\begin{block}{Question}
			The GREATEST benefit of implementing a risk treatment plan is:
		\end{block}
		\begin{enumerate}[(A)]
			\item To reduce the impact and likelihood of risk occurrence
			\item To identify unmitigated risk that can be transferred
			\item To exploit the risk to test organizational preparedness
			\item To enhance the overall risk appetite of the enterprise
		\end{enumerate}
	\end{frame}
	\begin{frame}{D3-Q227: Answer}
		\begin{block}{Correct Answer: A}
			Implementing the risk treatment plan reduces the negative impact and likelihood of a risk occurrence.
		\end{block}
	\end{frame}
	
	\begin{frame}{D3-Q228: Validating Patching Program}
		\begin{block}{Question}
			What is the BEST method to validate the effectiveness of an enterprise's patching program?
		\end{block}
		\begin{enumerate}[(A)]
			\item Conduct penetration testing
			\item Conduct a risk identification initiative
			\item Carry out vulnerability scans
			\item Review the requests for change
		\end{enumerate}
	\end{frame}
	\begin{frame}{D3-Q228: Answer}
		\begin{block}{Correct Answer: C}
			Performing vulnerability scans enables the IT risk practitioner to determine if patches are being installed on a timely basis.
		\end{block}
	\end{frame}
	
	\begin{frame}{D3-Q229: Designing IS Controls in Complex Environment}
		\begin{block}{Question}
			Which is the MOST important factor when designing information systems controls in a complex environment?
		\end{block}
		\begin{enumerate}[(A)]
			\item Development methodologies
			\item Scalability of the solution
			\item Technical platform interfaces
			\item Stakeholder requirements
		\end{enumerate}
	\end{frame}
	\begin{frame}{D3-Q229: Answer}
		\begin{block}{Correct Answer: D}
			The most important factor when designing IS controls is their ability to advance the interests of the business by addressing stakeholder requirements.
		\end{block}
	\end{frame}
	
	\begin{frame}{D3-Q230: EFT Testing Focus}
		\begin{block}{Question}
			A financial institution is testing its EFT system after major enhancements. The risk practitioner would MOST benefit from a test that:
		\end{block}
		\begin{enumerate}[(A)]
			\item Identifies the introduction of potential new gaps in security
			\item Verifies adequate system recovery in case of failure
			\item Ensures the system performs to expectation
			\item Ensures the system can support transaction volume
		\end{enumerate}
	\end{frame}
	\begin{frame}{D3-Q230: Answer}
		\begin{block}{Correct Answer: A}
			The security of an EFT system is of the most interest to a risk practitioner. Testing should identify the introduction of potential new gaps in information security.
		\end{block}
	\end{frame}
	
	\begin{frame}{D3-Q231: Creating KRIs for Quarterly Reporting}
		\begin{block}{Question}
			What is the BEST approach for creating KRIs for quarterly reporting to senior leadership?
		\end{block}
		\begin{enumerate}[(A)]
			\item Survey senior leaders about primary risk concerns
			\item Identify the most common network vulnerabilities
			\item Determine KRIs used in similar industry verticals
			\item Identify the enterprise risk appetite and metrics of current risk
		\end{enumerate}
	\end{frame}
	\begin{frame}{D3-Q231: Answer}
		\begin{block}{Correct Answer: D}
			KRIs are measurable and quantifiable indicators of risk. Creation of KRIs should be based on the enterprise's risk appetite and the measures available.
		\end{block}
	\end{frame}
	
	\begin{frame}{D3-Q232: Gauging Stakeholder Concern}
		\begin{block}{Question}
			An enterprise contracted an external supplier for critical components. Risk tolerance levels for the outsourced component are documented and approved. Which can serve to gauge risk that may trigger stakeholder concern?
		\end{block}
		\begin{enumerate}[(A)]
			\item Indicators with approved thresholds
			\item Approved status report of completed milestones
			\item List of controls to be implemented by the supplier
			\item Number of findings in external audit reports
		\end{enumerate}
	\end{frame}
	\begin{frame}{D3-Q232: Answer}
		\begin{block}{Correct Answer: A}
			Indicators with approved thresholds demonstrate the acceptable risk levels stakeholders are willing to tolerate.
		\end{block}
	\end{frame}
	
	\begin{frame}{D3-Q233: Revoking Temporary User Access}
		\begin{block}{Question}
			Which control within the user provision process BEST ensures revocation of system access for contractors and temporary users when no longer required?
		\end{block}
		\begin{enumerate}[(A)]
			\item Log account usage and send it to managers
			\item Establish predetermined, automatic expiration dates
			\item Ensure each user signs a security acknowledgment
			\item Require managers to email security when the user leaves
		\end{enumerate}
	\end{frame}
	\begin{frame}{D3-Q233: Answer}
		\begin{block}{Correct Answer: B}
			Predetermined expiration dates are the most effective means of removing systems access for temporary users.
		\end{block}
	\end{frame}
	
	\begin{frame}{D3-Q234: Presenting Risk Profile to Board}
		\begin{block}{Question}
			Which BEST helps while presenting the current risk profile to executive management and the board?
		\end{block}
		\begin{enumerate}[(A)]
			\item Risk response dashboard
			\item Emerging risk report
			\item Risk register dashboard
			\item Key risk indicators report
		\end{enumerate}
	\end{frame}
	\begin{frame}{D3-Q234: Answer}
		\begin{block}{Correct Answer: C}
			A risk register dashboard provides a comprehensive overview of the risk profile of the enterprise.
		\end{block}
	\end{frame}
	
	\begin{frame}{D3-Q235: Senior Management Monitoring KCIs}
		\begin{block}{Question}
			What is the MAIN reason senior management monitors and analyzes trends in key control indicators?
		\end{block}
		\begin{enumerate}[(A)]
			\item It provides feedback on the overall control environment
			\item It helps in identifying redundant controls
			\item It proactively identifies impacts to the risk profile
			\item It helps determine if additional controls are required
		\end{enumerate}
	\end{frame}
	\begin{frame}{D3-Q235: Answer}
		\begin{block}{Correct Answer: C}
			The primary objective of KCIs is to ensure that controls actually mitigate risk at an effective level.
		\end{block}
	\end{frame}
	
	\begin{frame}{D3-Q236: Risk Treatment Plans Purpose}
		\begin{block}{Question}
			Risk treatment plans are necessary to describe how the:
		\end{block}
		\begin{enumerate}[(A)]
			\item Identified risk is further analyzed
			\item Chosen treatment options will be implemented
			\item Accepted risk is treated
			\item Risk indicators will monitor the risk
		\end{enumerate}
	\end{frame}
	\begin{frame}{D3-Q236: Answer}
		\begin{block}{Correct Answer: B}
			A risk treatment plan includes the plan of action for the chosen treatment, how it will be implemented, who will implement it, key dates, and resource requirements.
		\end{block}
	\end{frame}
	
	\begin{frame}{D3-Q237: Cost-Benefit Analysis Stage}
		\begin{block}{Question}
			In the risk management process, a cost-benefit analysis is MAINLY performed:
		\end{block}
		\begin{enumerate}[(A)]
			\item As part of an initial risk assessment
			\item As part of risk-response planning
			\item During an information asset valuation
			\item When insurance is calculated for risk transfer
		\end{enumerate}
	\end{frame}
	\begin{frame}{D3-Q237: Answer}
		\begin{block}{Correct Answer: B}
			In risk response, a range of controls that can mitigate risk will be identified; a cost-benefit analysis helps identify the right controls to address the risk at acceptable levels within the budget.
		\end{block}
	\end{frame}
	
	\begin{frame}{D3-Q238: Insufficient Resources for Treatment}
		\begin{block}{Question}
			In the event that available resources for risk treatment are not sufficient, the risk treatment plan should:
		\end{block}
		\begin{enumerate}[(A)]
			\item Define the priorities across all treatments to assist in resource allocation
			\item Recommend postponing treatment until resources are available
			\item Recommend reassessing the risk treatment options
			\item Suggest increasing the priority of the treatment
		\end{enumerate}
	\end{frame}
	\begin{frame}{D3-Q238: Answer}
		\begin{block}{Correct Answer: A}
			When available resources for treating the risk are not sufficient, the risk treatment plan should include prioritization of the risk so resources can be properly allocated.
		\end{block}
	\end{frame}
	
	\begin{frame}{D3-Q239: System Backup and Restore Classification}
		\begin{block}{Question}
			System backup and restore procedures can BEST be classified as:
		\end{block}
		\begin{enumerate}[(A)]
			\item Technical controls
			\item Detective controls
			\item Corrective controls
			\item Deterrent controls
		\end{enumerate}
	\end{frame}
	\begin{frame}{D3-Q239: Answer}
		\begin{block}{Correct Answer: C --- Corrective controls}
			Corrective controls remediate the impact from negative events. If a system suffers harm so extensive that processing cannot continue, backup restore procedures enable that system to be recovered.
		\end{block}
	\end{frame}
	
	\begin{frame}{D3-Q240: Primary Reason for Policy Exception}
		\begin{block}{Question}
			A PRIMARY reason for initiating a policy-exception process is:
		\end{block}
		\begin{enumerate}[(A)]
			\item The risk is justified by the benefit
			\item Policy compliance is difficult to enforce
			\item Operations are too busy to comply
			\item Users may initially be inconvenienced
		\end{enumerate}
	\end{frame}
	\begin{frame}{D3-Q240: Answer}
		\begin{block}{Correct Answer: A}
			Exceptions to policies are warranted if the benefits outweigh the costs of policy compliance; however, the enterprise needs to assess both the tangible and intangible risk.
		\end{block}
	\end{frame}
	
	\begin{frame}{D3-Q241: Measuring Operational Controls}
		\begin{block}{Question}
			Which would BEST measure the effectiveness of operational controls?
		\end{block}
		\begin{enumerate}[(A)]
			\item Control matrix
			\item Key performance indicator
			\item Statement of applicability
			\item Key control indicator
		\end{enumerate}
	\end{frame}
	\begin{frame}{D3-Q241: Answer}
		\begin{block}{Correct Answer: D --- KCI}
			Key control indicators, also referred to as control effectiveness indicators, are metrics that provide information on the degree to which a control is working.
		\end{block}
	\end{frame}
	
	\begin{frame}{D3-Q242: Selecting Security Technologies}
		\begin{block}{Question}
			Security technologies should be selected PRIMARILY on the basis of their:
		\end{block}
		\begin{enumerate}[(A)]
			\item Evaluation in security publications
			\item Compliance with industry standards
			\item Ability to mitigate risk to organizational objectives
			\item Cost compared to the enterprise's IT budget
		\end{enumerate}
	\end{frame}
	\begin{frame}{D3-Q242: Answer}
		\begin{block}{Correct Answer: C}
			The most fundamental criterion for selecting security technology is the capacity to reduce risk to organizational objectives.
		\end{block}
	\end{frame}
	
	\begin{frame}{D3-Q243: Before Black Box Penetration Test}
		\begin{block}{Question}
			Which should be in place before a black box penetration test begins?
		\end{block}
		\begin{enumerate}[(A)]
			\item A clearly stated definition of scope
			\item Previous test results
			\item Proper communication and awareness training
			\item An incident response plan
		\end{enumerate}
	\end{frame}
	\begin{frame}{D3-Q243: Answer}
		\begin{block}{Correct Answer: A}
			A clearly stated definition of scope ensures a proper understanding of risk and success criteria.
		\end{block}
	\end{frame}
	
	\begin{frame}{D3-Q244: Securing Corporate Network Against External Attack}
		\begin{block}{Question}
			What is the BEST way to ensure that a corporate network is adequately secured against external attack?
		\end{block}
		\begin{enumerate}[(A)]
			\item Use an intrusion detection system
			\item Establish minimum security baselines
			\item Implement vendor recommended settings
			\item Perform periodic penetration testing
		\end{enumerate}
	\end{frame}
	\begin{frame}{D3-Q244: Answer}
		\begin{block}{Correct Answer: D}
			Penetration testing is the best way to ensure that perimeter security is adequate.
		\end{block}
	\end{frame}
	
	\begin{frame}{D3-Q245: Before Establishing KCIs}
		\begin{block}{Question}
			Which is PRIMARILY defined before establishing key control indicators?
		\end{block}
		\begin{enumerate}[(A)]
			\item Desired capacity
			\item Desired risk threshold
			\item Desired tolerances
			\item Desired risk appetite
		\end{enumerate}
	\end{frame}
	\begin{frame}{D3-Q245: Answer}
		\begin{block}{Correct Answer: C --- Desired tolerances}
			The goal of KCIs is to track performance of control actions relative to tolerances, providing insight into the ongoing adequacy of a control in keeping risk within acceptable levels.
		\end{block}
	\end{frame}
	
	\begin{frame}{D3-Q246: Substantive Test of Tape Inventory}
		\begin{block}{Question}
			A substantive test to verify that tape library inventory records are accurate involves:
		\end{block}
		\begin{enumerate}[(A)]
			\item Determining whether bar code readers are installed
			\item Conducting a physical count of the tape inventory
			\item Checking whether receipts and issues of tapes are accurately recorded
			\item Determining whether movement of tapes is authorized
		\end{enumerate}
	\end{frame}
	\begin{frame}{D3-Q246: Answer}
		\begin{block}{Correct Answer: B}
			A substantive test involves gathering evidence to evaluate the integrity of individual transactions, data, or other information. A physical count of the tape inventory is a substantive test.
		\end{block}
	\end{frame}
	
	\begin{frame}{D3-Q247: Ensuring Mitigation of Vulnerabilities}
		\begin{block}{Question}
			Which BEST ensures that appropriate mitigation occurs on identified information systems vulnerabilities?
		\end{block}
		\begin{enumerate}[(A)]
			\item Presenting root cause analysis to management
			\item Implementing software to input action points
			\item Incorporating findings into the annual report
			\item Assigning action plans with deadlines to responsible personnel
		\end{enumerate}
	\end{frame}
	\begin{frame}{D3-Q247: Answer}
		\begin{block}{Correct Answer: D}
			Assigning mitigation to personnel establishes responsibility for its completion within the deadline.
		\end{block}
	\end{frame}
	
	\begin{frame}{D3-Q248: Audit Report on Uncritical Functions}
		\begin{block}{Question}
			What is the BEST action if an audit report recommends supplemental controls for uncritical functions?
		\end{block}
		\begin{enumerate}[(A)]
			\item Appoint a project lead to start implementation
			\item Eliminate the risk by purchasing risk insurance
			\item Request that IT implement all supplemental controls
			\item Request a cost-benefit analysis to determine business need
		\end{enumerate}
	\end{frame}
	\begin{frame}{D3-Q248: Answer}
		\begin{block}{Correct Answer: D}
			A cost-benefit analysis may help determine if supplemental controls are needed, ensuring the cost of the control does not exceed the impact from the risk being realized.
		\end{block}
	\end{frame}
	
	\begin{frame}{D3-Q249: Peer Review of Risk Management}
		\begin{block}{Question}
			Which BEST enables a peer review of an enterprise's risk management process?
		\end{block}
		\begin{enumerate}[(A)]
			\item A balanced scorecard
			\item An industry survey
			\item A capability maturity model
			\item A framework
		\end{enumerate}
	\end{frame}
	\begin{frame}{D3-Q249: Answer}
		\begin{block}{Correct Answer: C}
			A capability maturity model describes essential elements and criteria for effective processes and outlines an evolutionary improvement path from ad hoc to disciplined, mature processes.
		\end{block}
	\end{frame}
	
	\begin{frame}{D3-Q250: Risk Treatment Plan With Resolution}
		\begin{block}{Question}
			Which is the MOST important information to include in a risk treatment plan that already has an appropriate resolution and a date for completion?
		\end{block}
		\begin{enumerate}[(A)]
			\item Responsible personnel
			\item Mitigating factors
			\item Likelihood of occurrence
			\item Cost of completion
		\end{enumerate}
	\end{frame}
	\begin{frame}{D3-Q250: Answer}
		\begin{block}{Correct Answer: A}
			Risk response activities must be assigned to a responsible person or group; without this assignment, it will be unclear who will implement the countermeasure.
		\end{block}
	\end{frame}
	
	\begin{frame}{D3-Q251: Ensuring Risk Remains Acceptable}
		\begin{block}{Question}
			Which BEST ensures that identified risk remains at an acceptable level?
		\end{block}
		\begin{enumerate}[(A)]
			\item Reviewing controls periodically according to the risk treatment plan
			\item Listing each risk as a separate entry in the risk register
			\item Creating a separate risk register for every department
			\item Maintaining a KRI for assets in the risk register
		\end{enumerate}
	\end{frame}
	\begin{frame}{D3-Q251: Answer}
		\begin{block}{Correct Answer: A}
			Controls deployed according to the risk treatment plan should provide the desired results, because the plan is based on management's acceptance of residual risk and approval of deployment steps.
		\end{block}
	\end{frame}
	
	\begin{frame}{D3-Q252: Risk Response Increasing Liability}
		\begin{block}{Question}
			Which risk response option is MOST likely to increase the liability of the enterprise?
		\end{block}
		\begin{enumerate}[(A)]
			\item Risk acceptance
			\item Risk reduction
			\item Risk transfer
			\item Risk avoidance
		\end{enumerate}
	\end{frame}
	\begin{frame}{D3-Q252: Answer}
		\begin{block}{Correct Answer: A --- Risk acceptance}
			An enterprise may choose to accept risk without knowing the correct level of risk being accepted; this may result in accusations of negligence.
		\end{block}
	\end{frame}
	
	\begin{frame}{D3-Q253: Best Reason Not to Reduce Risk}
		\begin{block}{Question}
			Which is the BEST reason for an enterprise to decide not to reduce an identified risk?
		\end{block}
		\begin{enumerate}[(A)]
			\item There is no regulatory requirement to reduce the risk
			\item There are mitigating controls in place
			\item The cost of mitigation exceeds the risk
			\item The budget for risk mitigation is limited
		\end{enumerate}
	\end{frame}
	\begin{frame}{D3-Q253: Answer}
		\begin{block}{Correct Answer: C}
			Enterprises will accept the risk when the cost of mitigation exceeds the risk.
		\end{block}
	\end{frame}
	
	\begin{frame}{D3-Q254: Outsourcing Part of IT Activities}
		\begin{block}{Question}
			An enterprise decides to address risk associated with an IT project by outsourcing part of the IT activities to a third party with a specialized skill set. In relation to the project, this is an example of:
		\end{block}
		\begin{enumerate}[(A)]
			\item Risk transfer
			\item Risk avoidance
			\item Risk acceptance
			\item Risk mitigation
		\end{enumerate}
	\end{frame}
	\begin{frame}{D3-Q254: Answer}
		\begin{block}{Correct Answer: D --- Risk mitigation}
			When specific activities are outsourced to an entity with a specialized skill set, the inherent risk of the activity is reduced. The outsourcing process works as a control in this case.
		\end{block}
	\end{frame}
	
	\begin{frame}{D3-Q255: Facilitating Cost-Effective Risk Response}
		\begin{block}{Question}
			Which BEST facilitates cost-effective risk response?
		\end{block}
		\begin{enumerate}[(A)]
			\item Prioritizing and addressing risk according to the risk management strategy
			\item Mitigating risk on the basis of likelihood and magnitude of impact
			\item Performing countermeasure analysis for each control deployed
			\item Selecting controls that are at zero or near-zero costs
		\end{enumerate}
	\end{frame}
	\begin{frame}{D3-Q255: Answer}
		\begin{block}{Correct Answer: A}
			Prioritizing and addressing risk in line with the risk treatment strategy balances the costs and benefits of managing the IT risk.
		\end{block}
	\end{frame}
	
	\begin{frame}{D3-Q256: Recommending DLP Device}
		\begin{block}{Question}
			Which can BEST be used as a basis for recommending a data loss prevention (DLP) device as a security control?
		\end{block}
		\begin{enumerate}[(A)]
			\item Benchmarking with peers on DLP deployment
			\item A business case for DLP to protect data
			\item Evaluation report of popular DLP solutions
			\item Data leakage scenario in risk register
		\end{enumerate}
	\end{frame}
	\begin{frame}{D3-Q256: Answer}
		\begin{block}{Correct Answer: B}
			A business case with costs versus benefits provides the business reasoning for why the DLP solution addresses the risk and explains how risk losses could be reduced.
		\end{block}
	\end{frame}
	
	\begin{frame}{D3-Q257: Cost Greater Than Opportunity}
		\begin{block}{Question}
			When the cost of risk related to a specific business process is greater than the potential opportunity, the BEST risk response is:
		\end{block}
		\begin{enumerate}[(A)]
			\item Transfer
			\item Acceptance
			\item Mitigation
			\item Avoidance
		\end{enumerate}
	\end{frame}
	\begin{frame}{D3-Q257: Answer}
		\begin{block}{Correct Answer: D --- Avoidance}
			Risk avoidance is the process for systematically avoiding risk, constituting one approach to managing risk.
		\end{block}
	\end{frame}
	
	\begin{frame}{D3-Q258: First Step When Outsourcing All IT}
		\begin{block}{Question}
			What is the FIRST step for a risk practitioner when an enterprise decides to outsource all IT services and support to a third party?
		\end{block}
		\begin{enumerate}[(A)]
			\item Validate that the internal systems of the provider are secure
			\item Enforce transborder data flow regulations
			\item Ensure security requirements are addressed in all contracts and agreements
			\item Build a business case to perform an on-site audit of the vendor
		\end{enumerate}
	\end{frame}
	\begin{frame}{D3-Q258: Answer}
		\begin{block}{Correct Answer: C}
			A contract only covers topics listed in the contract. If security is not explicitly included in the contract terms, the enterprise may not be properly protected.
		\end{block}
	\end{frame}
	
	\begin{frame}{D3-Q259: Justifying IT Control Cost}
		\begin{block}{Question}
			Which BEST helps ensure that the cost is justifiable when selecting an IT control?
		\end{block}
		\begin{enumerate}[(A)]
			\item The investment cost is within budget
			\item The risk likelihood and impact are reflected
			\item The net present value of the IT control cost is high
			\item Low-cost open source technology is used
		\end{enumerate}
	\end{frame}
	\begin{frame}{D3-Q259: Answer}
		\begin{block}{Correct Answer: B}
			The total cost of ownership of a control should not exceed the projected likelihood times the impact of the risk it is intended to mitigate.
		\end{block}
	\end{frame}
	
	\begin{frame}{D3-Q260: Testing Compliance of Response Plan}
		\begin{block}{Question}
			Testing compliance of a response and recovery plan should begin by conducting a:
		\end{block}
		\begin{enumerate}[(A)]
			\item Tabletop exercise
			\item Review of archived logs
			\item Penetration test
			\item Business impact analysis
		\end{enumerate}
	\end{frame}
	\begin{frame}{D3-Q260: Answer}
		\begin{block}{Correct Answer: A}
			Tabletop exercises simulate incidents to test the response capability of an enterprise.
		\end{block}
	\end{frame}
	
	\begin{frame}{D3-Q261: Capability Maturity Model Basis}
		\begin{block}{Question}
			The use of a capability maturity model is based on:
		\end{block}
		\begin{enumerate}[(A)]
			\item The training of staff
			\item The development of new controls
			\item The application of standard, repeatable processes that can be measured
			\item Users developing new innovative solutions
		\end{enumerate}
	\end{frame}
	\begin{frame}{D3-Q261: Answer}
		\begin{block}{Correct Answer: C}
			The use of maturity models requires development, implementation, and measuring of consistent procedures and activities.
		\end{block}
	\end{frame}
	
	\begin{frame}{D3-Q262: Internal Control Requirements Source}
		\begin{block}{Question}
			An enterprise is implementing controls to protect employee details from unauthorized exposure. The internal control requirements will come from:
		\end{block}
		\begin{enumerate}[(A)]
			\item The risk management team
			\item Internal audit
			\item IT management
			\item Process owners
		\end{enumerate}
	\end{frame}
	\begin{frame}{D3-Q262: Answer}
		\begin{block}{Correct Answer: D --- Process owners}
			Process owners provide the internal control requirements based on the business needs and objectives.
		\end{block}
	\end{frame}
	
	\begin{frame}{D3-Q263: Costs Balanced with Potential Loss}
		\begin{block}{Question}
			Which costs MUST be balanced with the potential for loss to ensure effective risk management?
		\end{block}
		\begin{enumerate}[(A)]
			\item Regulatory fines and penalties
			\item Total cost of ownership
			\item Breach management cost
			\item Business operation cost
		\end{enumerate}
	\end{frame}
	\begin{frame}{D3-Q263: Answer}
		\begin{block}{Correct Answer: B}
			Total cost of ownership considers implementation and operating costs for a control, which must be balanced with the potential for loss.
		\end{block}
	\end{frame}
	
	\begin{frame}{D3-Q264: Controls Design Reduction}
		\begin{block}{Question}
			Controls are most effective when they are designed to reduce:
		\end{block}
		\begin{enumerate}[(A)]
			\item Threats
			\item Likelihood
			\item Uncertainty
			\item Vulnerabilities
		\end{enumerate}
	\end{frame}
	\begin{frame}{D3-Q264: Answer}
		\begin{block}{Correct Answer: D --- Vulnerabilities}
			Controls are most effective when designed to reduce vulnerabilities affecting the enterprise.
		\end{block}
	\end{frame}
	
	\begin{frame}{D3-Q265: Mitigating Patch Interoperability Risk}
		\begin{block}{Question}
			System administrators may release critical patches into production without testing. Which would BEST mitigate the risk of interoperability issues?
		\end{block}
		\begin{enumerate}[(A)]
			\item Ensure a reliable system rollback plan is in place
			\item Test the patch on the least critical systems first
			\item Only allow updates to occur after hours
			\item Ensure patches are approved by the CISO
		\end{enumerate}
	\end{frame}
	\begin{frame}{D3-Q265: Answer}
		\begin{block}{Correct Answer: A}
			A reliable system rollback plan allows administrators to roll back patches in case they affect the system negatively.
		\end{block}
	\end{frame}
	
	\begin{frame}{D3-Q266: Mitigating Control Risk}
		\begin{block}{Question}
			Which BEST mitigates control risk?
		\end{block}
		\begin{enumerate}[(A)]
			\item Continuous monitoring
			\item An effective security awareness program
			\item Effective change management procedures
			\item Senior management support for control implementation
		\end{enumerate}
	\end{frame}
	\begin{frame}{D3-Q266: Answer}
		\begin{block}{Correct Answer: A}
			Continuous monitoring tests controls that mitigate the risk of the control being less effective over time.
		\end{block}
	\end{frame}
	
	\begin{frame}{D3-Q267: Aligning Controls with Business Objectives}
		\begin{block}{Question}
			When aligning controls with business objectives, what is MOST important?
		\end{block}
		\begin{enumerate}[(A)]
			\item Monitoring control activities periodically
			\item Ensuring ownership of key control activities
			\item Reviewing the risk management strategy
			\item Prioritizing control activities based on residual risk
		\end{enumerate}
	\end{frame}
	\begin{frame}{D3-Q267: Answer}
		\begin{block}{Correct Answer: B}
			Ensuring ownership of key control activities is the most important factor in assigning control responsibility and control accountability.
		\end{block}
	\end{frame}
	
	\begin{frame}{D3-Q268: Purchasing Insurance}
		\begin{block}{Question}
			Purchasing insurance is a form of:
		\end{block}
		\begin{enumerate}[(A)]
			\item Risk avoidance
			\item Risk mitigation
			\item Risk acceptance
			\item Risk transfer
		\end{enumerate}
	\end{frame}
	\begin{frame}{D3-Q268: Answer}
		\begin{block}{Correct Answer: D}
			Transferring risk typically involves insurance policies to share the financial risk.
		\end{block}
	\end{frame}
	
	\begin{frame}{D3-Q269: Selecting and Prioritizing Risk Responses}
		\begin{block}{Question}
			Which would PRIMARILY help an enterprise select and prioritize risk responses?
		\end{block}
		\begin{enumerate}[(A)]
			\item A cost-benefit analysis of available risk mitigation options
			\item The level of acceptable risk per risk appetite
			\item The potential to transfer or eliminate the risk
			\item The number of controls necessary to reduce the risk
		\end{enumerate}
	\end{frame}
	\begin{frame}{D3-Q269: Answer}
		\begin{block}{Correct Answer: A}
			The selection and priorities of a risk response will consider the cost-benefit of the various risk mitigation options in order to get the highest return on investment.
		\end{block}
	\end{frame}
	
	\begin{frame}{D3-Q270: Validating IT Risk Response Effectiveness}
		\begin{block}{Question}
			Business stakeholders and decision makers reviewing the effectiveness of IT risk responses would PRIMARILY validate whether:
		\end{block}
		\begin{enumerate}[(A)]
			\item IT controls eliminate the risk in question
			\item IT controls are continuously monitored
			\item IT controls achieve the desired objectives
			\item IT risk indicators are formally documented
		\end{enumerate}
	\end{frame}
	\begin{frame}{D3-Q270: Answer}
		\begin{block}{Correct Answer: C}
			The stakeholders are most interested in whether the control meets the stated objectives.
		\end{block}
	\end{frame}
	
	\begin{frame}{D3-Q271: Addressing Risk Exceeding Risk Acceptance}
		\begin{block}{Question}
			A risk assessment indicates a risk to the enterprise that exceeds the risk acceptance level set by senior management. What is the BEST way to address this risk?
		\end{block}
		\begin{enumerate}[(A)]
			\item Ensure that the risk is quickly brought within acceptable limits, regardless of cost
			\item Recommend mitigating controls if the cost and benefit would justify the controls
			\item Recommend that senior management revise the risk acceptance level
			\item Ensure that risk calculations are performed to revalidate the controls
		\end{enumerate}
	\end{frame}
	\begin{frame}{D3-Q271: Answer}
		\begin{block}{Correct Answer: B}
			Risk mitigating controls should be implemented based on cost and benefit. Controls are not justified if the cost of the control exceeds the benefit obtained.
		\end{block}
	\end{frame}
	
	\begin{frame}{D3-Q272: Critical Risk Without Resources}
		\begin{block}{Question}
			Which action is BEST when a critical risk has been identified and the resources to mitigate are not immediately available?
		\end{block}
		\begin{enumerate}[(A)]
			\item Log the risk in the risk register and review it with senior management regularly
			\item Capture the risk in the risk register once resources are available
			\item Escalate the risk report to senior management to obtain the resources to mitigate the risk
			\item Review the risk level with senior management and determine whether risk calculations are correct
		\end{enumerate}
	\end{frame}
	\begin{frame}{D3-Q272: Answer}
		\begin{block}{Correct Answer: C}
			If resources are not available or priorities need to be adjusted, it is important to engage senior management to assist in escalating the remediation.
		\end{block}
	\end{frame}
	
	\begin{frame}{D3-Q273: Cost-Benefit Analysis Stage in Risk Management}
		\begin{block}{Question}
			During which stage of the overall risk management process is the cost-benefit analysis PRIMARILY performed? During the:
		\end{block}
		\begin{enumerate}[(A)]
			\item Initial risk assessment
			\item Information asset classification
			\item Definition of the risk profile
			\item Risk response selection
		\end{enumerate}
	\end{frame}
	\begin{frame}{D3-Q273: Answer}
		\begin{block}{Correct Answer: D}
			During risk response selection, a range of controls that can mitigate the risk will be identified; the cost-benefit analysis in this process will help identify the right controls to address the risk at acceptable levels within the budget.
		\end{block}
	\end{frame}
	
	\begin{frame}{D3-Q274: Addressing Significant System Vulnerabilities}
		\begin{block}{Question}
			Which approach BEST helps address significant system vulnerabilities discovered during a network scan?
		\end{block}
		\begin{enumerate}[(A)]
			\item All significant vulnerabilities must be mitigated in a timely fashion
			\item Treatment should be based on threat, impact, and cost considerations
			\item Compensating controls must be implemented for major vulnerabilities
			\item Mitigation options should be proposed for management approval
		\end{enumerate}
	\end{frame}
	\begin{frame}{D3-Q274: Answer}
		\begin{block}{Correct Answer: B}
			The treatment should consider the degree of exposure and potential impact and the costs of various treatment options.
		\end{block}
	\end{frame}
	
	\begin{frame}{D3-Q275: Risk Response for Low Likelihood, High Impact}
		\begin{block}{Question}
			What is the BEST risk response for risk scenarios where the likelihood is low and financial impact is high?
		\end{block}
		\begin{enumerate}[(A)]
			\item Transfer the risk to a third party
			\item Accept the high cost of protection
			\item Implement detective controls
			\item Implement compensating controls
		\end{enumerate}
	\end{frame}
	\begin{frame}{D3-Q275: Answer}
		\begin{block}{Correct Answer: A}
			High-impact, low-likelihood situations are typically most cost-effectively covered by transferring the risk to a third party (e.g., insurance).
		\end{block}
	\end{frame}
	
	\begin{frame}{D3-Q276: Prioritization of Treatment Options}
		\begin{block}{Question}
			Faced with numerous risk scenarios, the prioritization of treatment options will be MOST effective when based on:
		\end{block}
		\begin{enumerate}[(A)]
			\item The existence of identified threats and vulnerabilities
			\item The likelihood of compromise and subsequent impact
			\item The results of vulnerability scans and exposure
			\item The exposure of corporate assets and operational risk
		\end{enumerate}
	\end{frame}
	\begin{frame}{D3-Q276: Answer}
		\begin{block}{Correct Answer: B}
			The probability of compromise and the likely impact will be the most important considerations for selecting treatment options.
		\end{block}
	\end{frame}
	
	\begin{frame}{D3-Q277: Example of Risk Sharing}
		\begin{block}{Question}
			Which activity is an example of risk sharing?
		\end{block}
		\begin{enumerate}[(A)]
			\item Moving a function to another department
			\item Selling a product or service to another company
			\item Deploying redundant firewalls
			\item Contracting with a third party
		\end{enumerate}
	\end{frame}
	\begin{frame}{D3-Q277: Answer}
		\begin{block}{Correct Answer: D}
			Contracting with a third party to share the responsibility for supporting activities can provide a form of risk transference as long as it is documented in the outsourcing contract.
		\end{block}
	\end{frame}
	
	\begin{frame}{D3-Q278: Risk Response Parameter Decreasing Impact}
		\begin{block}{Question}
			Which risk response selection parameter results in a decrease in magnitude of an event?
		\end{block}
		\begin{enumerate}[(A)]
			\item Efficiency of response
			\item Cost of response
			\item Effectiveness of response
			\item Capability to implement response
		\end{enumerate}
	\end{frame}
	\begin{frame}{D3-Q278: Answer}
		\begin{block}{Correct Answer: C}
			Effectiveness of response is the extent to which the response reduces the likelihood and impact.
		\end{block}
	\end{frame}
	
	\begin{frame}{D3-Q279: Risk Treatment Plan Primary Purpose}
		\begin{block}{Question}
			The risk treatment plan PRIMARILY provides treatment for:
		\end{block}
		\begin{enumerate}[(A)]
			\item High risk areas reported to senior management
			\item Identified risk that exceeds risk tolerance
			\item Every risk identified on the risk register
			\item Risk that has already materialized
		\end{enumerate}
	\end{frame}
	\begin{frame}{D3-Q279: Answer}
		\begin{block}{Correct Answer: B}
			Risk treatment plans would cover all the risk identified in the risk register that exceeds the enterprise risk tolerance level and, therefore, needs to be further treated.
		\end{block}
	\end{frame}
	
	\begin{frame}{D3-Q280: Risk Treatment Plan Primary Specification}
		\begin{block}{Question}
			A risk treatment plan should PRIMARILY specify the:
		\end{block}
		\begin{enumerate}[(A)]
			\item Responsibility for implementing the chosen risk treatment
			\item Approach to integrate risk treatment into day-to-day operations
			\item Risk acceptance decisions made by the risk owners
			\item Best approach to implement all the identified risk treatment options
		\end{enumerate}
	\end{frame}
	\begin{frame}{D3-Q280: Answer}
		\begin{block}{Correct Answer: A}
			A risk treatment plan will provide the details needed to implement the risk treatment options and the selected controls, including the responsibility for implementing the chosen risk treatment.
		\end{block}
	\end{frame}
	
	% ============================================================
	% DOMAIN 4 --- INFORMATION TECHNOLOGY AND SECURITY (22%)
	% ============================================================
	\section{Domain 4: Information Technology and Security}
	
	\begin{frame}{D4-Q281: Information No Longer Required}
		\begin{block}{Question}
			Information no longer required to support the main purpose of the enterprise should be managed:
		\end{block}
		\begin{enumerate}[(A)]
			\item Under the retention policy
			\item Under the information classification policy
			\item Under the backup policy
			\item Under the business impact analysis
		\end{enumerate}
	\end{frame}
	\begin{frame}{D4-Q281: Answer}
		\begin{block}{Correct Answer: A --- Retention policy}
			Information no longer required should be analyzed under the retention policy to determine whether the enterprise is required to maintain the data for business, legal, or regulatory reasons.
		\end{block}
	\end{frame}
	
	\begin{frame}{D4-Q282: Protecting Breakthrough Technology}
		\begin{block}{Question}
			An enterprise developed a breakthrough technology providing a competitive edge. Which FIRST governs how this information is protected from within the enterprise?
		\end{block}
		\begin{enumerate}[(A)]
			\item The data classification policy
			\item The acceptable use policy
			\item Encryption standards
			\item The access control policy
		\end{enumerate}
	\end{frame}
	\begin{frame}{D4-Q282: Answer}
		\begin{block}{Correct Answer: A}
			A data classification policy describes the data classification categories, the level of protection for each category, and roles and responsibilities.
		\end{block}
	\end{frame}
	
	\begin{frame}{D4-Q283: Most Important IS Infrastructure Requirement}
		\begin{block}{Question}
			Which is the MOST important requirement for setting up an information security infrastructure for a new system?
		\end{block}
		\begin{enumerate}[(A)]
			\item Performing a business impact analysis
			\item Considering personal devices as part of the security policy
			\item Basing the information security infrastructure on a risk assessment
			\item Initiating an IT security awareness campaign
		\end{enumerate}
	\end{frame}
	\begin{frame}{D4-Q283: Answer}
		\begin{block}{Correct Answer: C}
			The information security infrastructure should be based on a risk assessment.
		\end{block}
	\end{frame}
	
	\begin{frame}{D4-Q284: Data Retention Policy Primary Concern}
		\begin{block}{Question}
			The PRIMARY concern of a risk practitioner reviewing a formal data retention policy is:
		\end{block}
		\begin{enumerate}[(A)]
			\item Storage availability
			\item Applicable organizational standards
			\item Generally accepted industry good practices
			\item Regulatory and business requirements
		\end{enumerate}
	\end{frame}
	\begin{frame}{D4-Q284: Answer}
		\begin{block}{Correct Answer: D}
			In determining the retention policy, regulatory requirements are of primary importance along with business requirements.
		\end{block}
	\end{frame}
	
	\begin{frame}{D4-Q285: Security Awareness Programs Reason}
		\begin{block}{Question}
			Which is the MOST important reason for conducting security awareness programs throughout an enterprise?
		\end{block}
		\begin{enumerate}[(A)]
			\item Reducing the risk of social engineering attacks
			\item Training personnel to respond to security incidents
			\item Informing business units about the security strategy
			\item Maintaining evidence of training records
		\end{enumerate}
	\end{frame}
	\begin{frame}{D4-Q285: Answer}
		\begin{block}{Correct Answer: A}
			Security awareness reduces the risk of successful social engineering attacks by sensitizing employees to security policies and risk.
		\end{block}
	\end{frame}
	
	\begin{frame}{D4-Q286: Incident Response Training Effective}
		\begin{block}{Question}
			Which is the BEST indicator that incident response training is effective?
		\end{block}
		\begin{enumerate}[(A)]
			\item Decreased reporting of security incidents
			\item Increased reporting of security incidents
			\item Decreased number of password resets
			\item Increased number of identified system vulnerabilities
		\end{enumerate}
	\end{frame}
	\begin{frame}{D4-Q286: Answer}
		\begin{block}{Correct Answer: B}
			Increased reporting of valid incidents is the best indicator because it shows users are aware of security rules and know how to report incidents.
		\end{block}
	\end{frame}
	
	\begin{frame}{D4-Q287: Greatest Information Security Awareness}
		\begin{block}{Question}
			Which provides the GREATEST level of information security awareness?
		\end{block}
		\begin{enumerate}[(A)]
			\item Job descriptions
			\item A security manual
			\item Security training
			\item An organizational diagram
		\end{enumerate}
	\end{frame}
	\begin{frame}{D4-Q287: Answer}
		\begin{block}{Correct Answer: C}
			Security training is the best way to inform all employees about changes to the risk landscape and enhance awareness of enterprise risk management strategy.
		\end{block}
	\end{frame}
	
	\begin{frame}{D4-Q288: Shared Generic Account}
		\begin{block}{Question}
			Which is the MAIN concern when two or more staff members are allowed to use the same generic account?
		\end{block}
		\begin{enumerate}[(A)]
			\item Segregation of duties
			\item Inability to change the password
			\item Repudiation
			\item Inability to trace account activities
		\end{enumerate}
	\end{frame}
	\begin{frame}{D4-Q288: Answer}
		\begin{block}{Correct Answer: C --- Repudiation}
			Because username and password are the same for generic accounts, repudiation becomes an issue. None of the users can be held accountable because each can deny accountability for transactions.
		\end{block}
	\end{frame}
	
	\begin{frame}{D4-Q289: IT Risk Awareness Campaign Approach}
		\begin{block}{Question}
			Which is the BEST approach when conducting an IT risk awareness campaign?
		\end{block}
		\begin{enumerate}[(A)]
			\item Provide technical detail on exploits
			\item Provide common messages tailored for different groups
			\item Target system administrators and help desk staff
			\item Target senior managers and business process owners
		\end{enumerate}
	\end{frame}
	\begin{frame}{D4-Q289: Answer}
		\begin{block}{Correct Answer: B}
			Groups differ in level of responsibility and expertise; tailor common messages to each group's role and understanding.
		\end{block}
	\end{frame}
	
	\begin{frame}{D4-Q290: Social Media Outlet Greatest Concern}
		\begin{block}{Question}
			Which principle of information security is of GREATEST concern to a social media outlet for political or news-related content?
		\end{block}
		\begin{enumerate}[(A)]
			\item Integrity
			\item Confidentiality
			\item Availability
			\item Nonrepudiation
		\end{enumerate}
	\end{frame}
	\begin{frame}{D4-Q290: Answer}
		\begin{block}{Correct Answer: C --- Availability}
			For a social media outlet, availability is of the greatest concern.
		\end{block}
	\end{frame}
	
	\begin{frame}{D4-Q291: Most Significant Threat to Project}
		\begin{block}{Question}
			Which poses the MOST significant threat to a project?
		\end{block}
		\begin{enumerate}[(A)]
			\item A lack of feedback upon project completion
			\item A lack of unit testing
			\item Missed opportunities from lessons learned
			\item Misunderstanding the requirements
		\end{enumerate}
	\end{frame}
	\begin{frame}{D4-Q291: Answer}
		\begin{block}{Correct Answer: D}
			Projects exist to deliver specific outcomes. If requirements are misunderstood, a project can succeed internally (schedule, budget) yet fail the business because it did not deliver business value.
		\end{block}
	\end{frame}
	
	\begin{frame}{D4-Q292: SDLC Phase for Risk Practitioner}
		\begin{block}{Question}
			In which phase of the SDLC should a risk practitioner FIRST become involved?
		\end{block}
		\begin{enumerate}[(A)]
			\item Analysis
			\item Design
			\item Planning
			\item Implementation
		\end{enumerate}
	\end{frame}
	\begin{frame}{D4-Q292: Answer}
		\begin{block}{Correct Answer: C --- Planning}
			The risk practitioner should become involved as early as possible. When participating in planning, they can influence resource requests to meet risk objectives most efficiently and effectively.
		\end{block}
	\end{frame}
	
	\begin{frame}{D4-Q293: Project Going Over Budget}
		\begin{block}{Question}
			Which is MOST likely to result in a project going over budget?
		\end{block}
		\begin{enumerate}[(A)]
			\item Testing
			\item Scope creep
			\item Loss of sponsorship
			\item Quality control
		\end{enumerate}
	\end{frame}
	\begin{frame}{D4-Q293: Answer}
		\begin{block}{Correct Answer: B --- Scope creep}
			Scope creep refers to uncontrolled changes in a project's scope. Unless scope is controlled, duration and budget cannot be held to account.
		\end{block}
	\end{frame}
	
	\begin{frame}{D4-Q294: Project Risk Identification}
		\begin{block}{Question}
			Project-related risk that can threaten the successful completion of a software development project is BEST identified by the:
		\end{block}
		\begin{enumerate}[(A)]
			\item Project manager
			\item Project sponsor
			\item Project steering committee
			\item Quality assurance team
		\end{enumerate}
	\end{frame}
	\begin{frame}{D4-Q294: Answer}
		\begin{block}{Correct Answer: A --- Project manager}
			The project manager is responsible for identification, analysis, and monitoring of project-related risk.
		\end{block}
	\end{frame}
	
	\begin{frame}{D4-Q295: VPN Implementation Concern}
		\begin{block}{Question}
			Which is of MOST concern in a review of a VPN implementation?
		\end{block}
		\begin{enumerate}[(A)]
			\item Computers on the network are located at the enterprise's remote offices
			\item Computers are on the enterprise's internal network
			\item Computers are at the backup site
			\item Computers are in employees' homes
		\end{enumerate}
	\end{frame}
	\begin{frame}{D4-Q295: Answer}
		\begin{block}{Correct Answer: D}
			Home computers are least often subject to the enterprise security policy and therefore are high-risk. Once a computer is compromised, any network that trusts it is at risk.
		\end{block}
	\end{frame}
	
	\begin{frame}{D4-Q296: BCM Meeting Stakeholder Expectations}
		\begin{block}{Question}
			Which BEST supports business continuity management in meeting external stakeholder expectations?
		\end{block}
		\begin{enumerate}[(A)]
			\item Prioritizing applications based on business criticality
			\item Ensuring backup data are available for restoration
			\item Disclosing the crisis management strategy statement
			\item Obtaining risk assessment by an independent party
		\end{enumerate}
	\end{frame}
	\begin{frame}{D4-Q296: Answer}
		\begin{block}{Correct Answer: A}
			External parties expect their information assets are secured. Prioritizing applications based on business criticality allows external expectations to be met optimally with limited resources.
		\end{block}
	\end{frame}
	
	\begin{frame}{D4-Q297: Cloud Computing Risk Controls}
		\begin{block}{Question}
			Which BEST identifies controls addressing risk related to cloud computing?
		\end{block}
		\begin{enumerate}[(A)]
			\item Data encryption, tenant isolation, controlled change management
			\item Data encryption, customizing the application template
			\item Use of open standards technology, data encryption, tenant isolation
			\item Tenant isolation, controlled change management, custom widgets
		\end{enumerate}
	\end{frame}
	\begin{frame}{D4-Q297: Answer}
		\begin{block}{Correct Answer: A}
			Encryption facilitates separation of data among tenants. Tenant isolation ensures that one tenant's data are sequestered from other tenants. Controlled change management ensures all changes are planned and tenant dependencies are mapped.
		\end{block}
	\end{frame}
	
	\begin{frame}{D4-Q298: Self-Signed Digital Certificates}
		\begin{block}{Question}
			A financial institution conducts high-value electronic transactions using customers' self-signed digital certificates. Which is of MOST concern?
		\end{block}
		\begin{enumerate}[(A)]
			\item Self-signed certificates may not be accepted by internal browsers
			\item Nonrepudiation requirements may not be enforced
			\item The certificate's private key may be susceptible to attack
			\item Self-signed certificates may lead to session hijacking
		\end{enumerate}
	\end{frame}
	\begin{frame}{D4-Q298: Answer}
		\begin{block}{Correct Answer: B}
			Nonrepudiation of transactions is critical in financial transactions. This can only be satisfied by using third party-trusted certificates signed by an accredited authority.
		\end{block}
	\end{frame}
	
	\begin{frame}{D4-Q299: Hardware vs Software Risk}
		\begin{block}{Question}
			Which risk consideration is of MOST concern to the risk practitioner regarding hardware as opposed to software?
		\end{block}
		\begin{enumerate}[(A)]
			\item Misconfiguration
			\item Access control
			\item Skilled resources
			\item Version control
		\end{enumerate}
	\end{frame}
	\begin{frame}{D4-Q299: Answer}
		\begin{block}{Correct Answer: A --- Misconfiguration}
			Misconfiguration can lead to potential vulnerabilities, leaving the system open to external parties, and is a larger concern for hardware than software.
		\end{block}
	\end{frame}
	
	\begin{frame}{D4-Q300: System Requirements Risk Phase}
		\begin{block}{Question}
			At which phase of the SDLC should risk related to system requirements be determined?
		\end{block}
		\begin{enumerate}[(A)]
			\item Development
			\item Implementation
			\item Initiation
			\item Maintenance
		\end{enumerate}
	\end{frame}
	\begin{frame}{D4-Q300: Answer}
		\begin{block}{Correct Answer: C --- Initiation}
			Upon initiation of the SDLC, identified risk should inform system requirements, including security requirements and a security concept of operations.
		\end{block}
	\end{frame}
	
	\begin{frame}{D4-Q301: Wrong Patient File}
		\begin{block}{Question}
			A healthcare practitioner providing care is given the file of a patient scheduled for an appointment later that day. What is the PRIMARY type of risk faced?
		\end{block}
		\begin{enumerate}[(A)]
			\item Relevance
			\item Integrity
			\item Security
			\item Availability
		\end{enumerate}
	\end{frame}
	\begin{frame}{D4-Q301: Answer}
		\begin{block}{Correct Answer: A --- Relevance risk}
			Relevance risk is the risk of not getting the right information to the right person at the right time for the correct action. The file provided is not relevant to that patient.
		\end{block}
	\end{frame}
	
	\begin{frame}{D4-Q302: Best Cloud Model for Confidential Info}
		\begin{block}{Question}
			An enterprise accepts the risk of confidential information in the cloud. Which cloud deployment model offers the most safeguards?
		\end{block}
		\begin{enumerate}[(A)]
			\item Public cloud
			\item Private cloud
			\item Community cloud
			\item Hybrid cloud
		\end{enumerate}
	\end{frame}
	\begin{frame}{D4-Q302: Answer}
		\begin{block}{Correct Answer: B --- Private cloud}
			The private cloud model operates solely for the enterprise and will have controls in place needed to keep enterprise information confidential.
		\end{block}
	\end{frame}
	
	\begin{frame}{D4-Q303: BCP Authorship}
		\begin{block}{Question}
			Business continuity plans should be written and maintained by:
		\end{block}
		\begin{enumerate}[(A)]
			\item The information security and IT functions
			\item Representatives from all functional units
			\item The risk management function
			\item Executive management
		\end{enumerate}
	\end{frame}
	\begin{frame}{D4-Q303: Answer}
		\begin{block}{Correct Answer: B}
			Business continuity planning is an enterprise-wide activity; it is successful only if all business owners collaborate in developing, testing, and maintaining the plan.
		\end{block}
	\end{frame}
	
	\begin{frame}{D4-Q304: Verifying Up-to-Date Antivirus}
		\begin{block}{Question}
			What is the BEST way to verify that critical production servers are using up-to-date antivirus signature files?
		\end{block}
		\begin{enumerate}[(A)]
			\item Check a sample of servers
			\item Verify the date signature files were last pushed out
			\item Use a benign virus to test quarantine
			\item Research the most recent signature file and compare to the console
		\end{enumerate}
	\end{frame}
	\begin{frame}{D4-Q304: Answer}
		\begin{block}{Correct Answer: A}
			The only effective way to verify currency of signature files is to look at a sample of servers.
		\end{block}
	\end{frame}
	
	\begin{frame}{D4-Q305: Security Awareness Program Effectiveness}
		\begin{block}{Question}
			Which is MOST important for measuring the effectiveness of a security awareness program?
		\end{block}
		\begin{enumerate}[(A)]
			\item Increased interest in focus groups on security
			\item A reduced number of security violation reports
			\item A quantitative evaluation to ensure user comprehension
			\item An increased number of security violation reports
		\end{enumerate}
	\end{frame}
	\begin{frame}{D4-Q305: Answer}
		\begin{block}{Correct Answer: D}
			An increase in the number of violation reports is the best indicator of a high level of security awareness.
		\end{block}
	\end{frame}
	
	\begin{frame}{D4-Q306: Major Risk of Single Sign-On}
		\begin{block}{Question}
			A MAJOR risk of single sign-on is that it:
		\end{block}
		\begin{enumerate}[(A)]
			\item Uses complex technologies for password management
			\item May potentially bypass the enterprise firewall
			\item Is prone to distributed denial-of-service attacks
			\item May be a potential single point of compromise
		\end{enumerate}
	\end{frame}
	\begin{frame}{D4-Q306: Answer}
		\begin{block}{Correct Answer: D}
			Because SSO uses a single username and password for access to multiple systems, loss or disclosure of the credentials may compromise a broad range of systems.
		\end{block}
	\end{frame}
	
	\begin{frame}{D4-Q307: Changeover Consistency of Data}
		\begin{block}{Question}
			Which changeover technique would have the GREATEST risk regarding consistency of data when moving users from one system to another?
		\end{block}
		\begin{enumerate}[(A)]
			\item Parallel
			\item Outsourcing
			\item Abrupt
			\item Phased
		\end{enumerate}
	\end{frame}
	\begin{frame}{D4-Q307: Answer}
		\begin{block}{Correct Answer: C --- Abrupt}
			Abrupt changeover occurs when users are converted immediately upon availability. It is usually least expensive but involves high risk of data loss and system failure.
		\end{block}
	\end{frame}
	
	\begin{frame}{D4-Q308: Insider Threat to Confidential Info}
		\begin{block}{Question}
			Which is the MOST effective measure against insider threats to confidential information?
		\end{block}
		\begin{enumerate}[(A)]
			\item Audit trail monitoring
			\item A privacy policy
			\item Role-based access control
			\item Defense in depth
		\end{enumerate}
	\end{frame}
	\begin{frame}{D4-Q308: Answer}
		\begin{block}{Correct Answer: C}
			Role-based access control provides access according to business need, reducing unnecessary access rights and enforcing accountability.
		\end{block}
	\end{frame}
	
	\begin{frame}{D4-Q309: Improving Awareness After Spear Phishing}
		\begin{block}{Question}
			Despite a comprehensive awareness program, a breach occurred via spear phishing. What is the MOST effective way to improve security awareness?
		\end{block}
		\begin{enumerate}[(A)]
			\item Review and improve coverage of social engineering threats
			\item Launch a disciplinary process against those who leaked information
			\item Perform periodic social engineering tests and communicate summary results
			\item Implement a DLP system that points users to corporate policies
		\end{enumerate}
	\end{frame}
	\begin{frame}{D4-Q309: Answer}
		\begin{block}{Correct Answer: C}
			Users who are aware of security threats may need a reminder that these threats are real. Periodic social engineering tests help maintain a level of alertness.
		\end{block}
	\end{frame}
	
	\begin{frame}{D4-Q310: Increased Malware Attacks Recommendation}
		\begin{block}{Question}
			IT operations indicates a heavy increase in malware attacks. Which is the MOST appropriate recommendation?
		\end{block}
		\begin{enumerate}[(A)]
			\item Contract with a new anti-malware vendor
			\item Close down the Internet connection to prevent employees from visiting infected sites
			\item Make the number of malware attacks part of each employee's performance metrics
			\item Increase employee awareness training, including roles and responsibilities
		\end{enumerate}
	\end{frame}
	\begin{frame}{D4-Q310: Answer}
		\begin{block}{Correct Answer: D}
			Employee awareness training helps the enterprise avoid and more quickly detect malware attacks.
		\end{block}
	\end{frame}
	
	\begin{frame}{D4-Q311: When Data Warehouse Should Remain In-House}
		\begin{block}{Question}
			Which MOST likely indicates that a customer data warehouse should remain in-house rather than outsourced offshore?
		\end{block}
		\begin{enumerate}[(A)]
			\item Telecommunications costs may be higher in the first year
			\item Privacy laws may prevent a cross-border flow of information
			\item Time zone differences may impede communications
			\item Software development may require more detailed specifications
		\end{enumerate}
	\end{frame}
	\begin{frame}{D4-Q311: Answer}
		\begin{block}{Correct Answer: B}
			Privacy laws prohibiting the cross-border flow of personally identifiable information may make it impossible to locate a data warehouse containing customer information in another country.
		\end{block}
	\end{frame}
	
	\begin{frame}{D4-Q312: Responsible for Information Classification}
		\begin{block}{Question}
			The person responsible for ensuring that information is classified is the:
		\end{block}
		\begin{enumerate}[(A)]
			\item Security manager
			\item Technology group
			\item Data owner
			\item Senior management
		\end{enumerate}
	\end{frame}
	\begin{frame}{D4-Q312: Answer}
		\begin{block}{Correct Answer: C --- Data owner}
			The data owner is responsible for applying the proper classification to the data.
		\end{block}
	\end{frame}
	
	\begin{frame}{D4-Q313: Transmitting Personal Information}
		\begin{block}{Question}
			Which is the MOST important when transmitting personal information across networks?
		\end{block}
		\begin{enumerate}[(A)]
			\item Encrypting the personal information
			\item Obtaining consent to transfer personal information
			\item Ensuring the privacy of the personal information
			\item Ensuring effective change management
		\end{enumerate}
	\end{frame}
	\begin{frame}{D4-Q313: Answer}
		\begin{block}{Correct Answer: C}
			Privacy protection is necessary to ensure that the receiving party has the appropriate level of protection of personal data.
		\end{block}
	\end{frame}
	
	\begin{frame}{D4-Q314: Excessive Workstation Images KRI}
		\begin{block}{Question}
			An excessive number of standard workstation images can be categorized as a key risk indicator for:
		\end{block}
		\begin{enumerate}[(A)]
			\item Change management
			\item Configuration management
			\item IT operations management
			\item Data management
		\end{enumerate}
	\end{frame}
	\begin{frame}{D4-Q314: Answer}
		\begin{block}{Correct Answer: B --- Configuration management}
			An excessive number of unique workstation images indicates poor configuration management processes and insufficient attention to actual business requirements.
		\end{block}
	\end{frame}
	
	\begin{frame}{D4-Q315: Addressing Data Leakage}
		\begin{block}{Question}
			Which BEST addresses the risk of data leakage?
		\end{block}
		\begin{enumerate}[(A)]
			\item Capacity planning procedures
			\item File backup procedures
			\item Acceptable use policies
			\item Database integrity checks
		\end{enumerate}
	\end{frame}
	\begin{frame}{D4-Q315: Answer}
		\begin{block}{Correct Answer: C --- Acceptable use policies}
			Acceptable use policies are an important control for preventing unauthorized disclosure of confidential information.
		\end{block}
	\end{frame}
	
	% ============================================================
	% FINAL SLIDES
	% ============================================================
	\section{Final Review}
	
	\begin{frame}{Exam Summary: Expanded Question Bank}
		\begin{block}{Domain Coverage}
			\begin{itemize}
				\item \textbf{Domain 1: Governance} --- 100 questions
				\item \textbf{Domain 2: IT Risk Assessment} --- 100 questions
				\item \textbf{Domain 3: Risk Response and Reporting} --- 80 questions
				\item \textbf{Domain 4: IT and Security} --- 35 questions
			\end{itemize}
		\end{block}
		\begin{exampleblock}{How to Use This Deck}
			Read each question, commit to an answer, then advance to the next slide to verify against the ISACA justification.
		\end{exampleblock}
	\end{frame}
	
	\begin{frame}{Key Takeaways Across Domains}
		\begin{block}{Golden Rules from the QAE Manual}
			\begin{itemize}
				\item Business management owns risk; senior management accepts it; the board is ultimately accountable
				\item Risk identification precedes assessment; you cannot protect what you do not know exists
				\item Residual risk = remaining risk after controls; acceptable risk = residual within tolerance
				\item BIA establishes RTO, RPO, MTD; risk assessment evaluates likelihood and impact
				\item Control types: preventive (stop), detective (find), corrective (fix), deterrent (discourage), compensating (alternate)
				\item KRIs predict risk; KPIs measure performance; KCIs measure control effectiveness
				\item Cost-benefit analysis justifies controls; when cost exceeds benefit, accept the risk
				\item Encryption = confidentiality; hashing = integrity; digital signature = integrity + nonrepudiation
			\end{itemize}
		\end{block}
	\end{frame}
	
	\begin{frame}{Final Advice}
		\begin{block}{Exam Strategy}
			\centering
			\vspace{0.3cm}
			\begin{itemize}
				\item Read the entire stem; look for \textbf{BEST}, \textbf{MOST}, \textbf{FIRST}
				\item Eliminate wrong choices, then choose the \textbf{business-focused} answer
				\item Understand the \textit{why} behind each justification
				\item Practice scenario-based reasoning --- application beats memorization
			\end{itemize}
			\par\vspace{0.5cm}
			\Large{\textbf{Good Luck on Your CRISC Exam!}}
		\end{block}
	\end{frame}
	
\end{document}